% Leçon 32 — Expression écrite : affiche publicitaire sur un produit (marketing) — Chinois Tle A — Cours signé OnBuch+
% Programme officiel MINESEC. Compiler avec : tectonic ZH32-ecrit-affiche-publicitaire-sur-un-produit-marketin.tex
\newcommand{\DOCMATIERE}{Chinois}
\newcommand{\DOCNIVEAU}{Tle A}
\newcommand{\DOCMODULE}{Module 4 : Activités économiques et monde du travail}
\newcommand{\DOCLECON}{ZH32}
\newcommand{\DOCTITRE}{Leçon 32 — Expression écrite : affiche publicitaire sur un produit (marketing)}
% OnBuch+ — préambule « cours signés » ENRICHI (compilé avec tectonic / XeLaTeX)
% Système visuel « Pop » d'OnBuch+ : fond crème (jamais blanc), bordures encre
% épaisses, ombres « dures » décalées, titres Archivo Black en capitales.
% Version MATHÉMATIQUES : graphes (pgfplots/tikz), boîtes pédagogiques
% (définition, propriété, méthode, attention, le savais-tu, exercice résolu,
% exercice, TP, bilan), page de garde et signature OnBuch+ ancrée en bas.
%
% Macros attendues dans main.tex AVANT \input{preamble} :
%   \DOCMATIERE  \DOCNIVEAU  \DOCMODULE  \DOCLECON (numéro)  \DOCTITRE
\documentclass[11pt]{article}

\usepackage{fontspec}
\usepackage[french]{babel} % français = langue principale ; le chinois via \zh{..} / textebox (xeCJK)
\usepackage{xeCJK}
\usepackage{pifont} % \ding{51} \ding{55} utilisés par les modèles
\usepackage[a4paper,margin=2cm,top=3.4cm,bottom=2.8cm,headheight=1.4cm,footskip=1.3cm]{geometry}
\usepackage{amsmath,amssymb,amsfonts}
\usepackage{mathtools} % \xmapsto, \coloneqq... souvent utilisés par les modèles
\usepackage{cancel} % simplifications barrées (bilans de forces)
% \abs{z} (module, valeur absolue) et \norm{u} (norme), utilisés par les modèles
\providecommand{\abs}{}\let\abs\relax\DeclarePairedDelimiter{\abs}{\lvert}{\rvert}
\providecommand{\norm}{}\let\norm\relax\DeclarePairedDelimiter{\norm}{\lVert}{\rVert}
\usepackage[table]{xcolor} % [table] : fournit aussi \rowcolors (alternance de lignes)
\usepackage{pagecolor}
\usepackage{tikz}
\usetikzlibrary{arrows.meta,positioning,calc,shapes.geometric,shapes.misc,shapes.arrows,decorations.pathmorphing,decorations.pathreplacing,decorations.markings,patterns,shadows,mindmap,trees,fit,backgrounds,matrix,intersections,angles,quotes}
\usepackage{pgfplots}
\usepackage[version=4]{mhchem} % équations nucléaires \ce{^{235}_{92}U}
\usepackage[european,siunitx]{circuitikz} % schémas électriques
\pgfplotsset{compat=1.18}
\usepgfplotslibrary{fillbetween,groupplots} % aires entre courbes (calcul intégral)
\usepackage{enumitem}
\usepackage{fancyhdr}
% [most] charge toutes les bibliothèques tcolorbox (listings, minted, raster,
% documentation...) et gonfle inutilement la mémoire TeX sur un document long
% (~300 boîtes) : on ne charge que ce qui est réellement utilisé.
\usepackage[skins,breakable]{tcolorbox}
\usepackage{graphicx}
\usepackage{colortbl}
\usepackage{booktabs}
\usepackage{tabularx}
\usepackage{diagbox} % case d'en-tête diagonale des tableaux à double entrée
% Colonnes extensibles centrée / gauche / droite (C, L, R), souvent utilisées par les modèles
\newcolumntype{C}{>{\centering\arraybackslash}X}
\newcolumntype{L}{>{\raggedright\arraybackslash}X}
\newcolumntype{R}{>{\raggedleft\arraybackslash}X}
\usepackage{array}
\usepackage{multicol}
\usepackage{multirow}
\usepackage{siunitx}
% parse-numbers=false : accepte aussi \SI{2,5 \times 10^{-3}}{...}
\sisetup{locale=FR,output-decimal-marker={,},group-minimum-digits=4,parse-numbers=false}
\DeclareSIUnit\ohm{\ensuremath{\symup{\Omega}}} % U+2126 absent de la police
\DeclareSIUnit\division{div} % réglages d'oscilloscope (ms/div, V/div)
\DeclareSIUnit\tr{tr} % tours (vitesse de rotation, ex. \SI{1500}{\tr\per\minute})
\DeclareSIUnit\molar{M} % concentration molaire (ex. \SI{3}{\molar}, \SI{1}{\micro\molar})
\DeclareSIUnit\an{an} % année (le modèle confond parfois avec \year, un compteur TeX) — ex. \SI{5}{\kilo\meter\per\an}
\DeclareSIUnit\FCFA{FCFA} % franc CFA (ex. \SI{500}{\FCFA\per\kilo\gram})
\DeclareSIUnit\degreeBrix{\textdegree Brix} % teneur en sucre d'un sirop/jus (réfractométrie)
\DeclareSIUnit\month{mois} % le modèle utilise parfois \month, un compteur TeX, comme unité de durée
\usepackage{qrcode}
% \degree : commande fréquemment utilisée par les modèles pour un angle ou une
% température en dehors de \SI{}{} (non fournie par les paquets chargés).
\providecommand{\degree}{^{\circ}}
% \qed (fin de démonstration), utilisé par les modèles sans amsthm
\providecommand{\qed}{\hfill$\square$}
% \barre : double barre des valeurs interdites dans un tableau de variations (tkz-tab)
\providecommand{\barre}{\ensuremath{\|}}
% environnement proof (amsthm non chargé) utilisé par les modèles
\ifdefined\proof\else\newenvironment{proof}[1][Démonstration]{\par\noindent\textit{#1.}\ }{\qed\par}\fi
\usepackage{lastpage}
\usepackage{hyperref}
\hypersetup{hidelinks}

% ============================================================
% Palette « Pop » OnBuch+ — mêmes hex que la classe P (plus_ui.dart)
% ============================================================
\definecolor{poporange}{HTML}{F59321}
\definecolor{popdark}{HTML}{E07A0C}
\definecolor{popcream}{HTML}{FFF6E8}
\definecolor{popink}{HTML}{1C1714}
\definecolor{poppurple}{HTML}{7A5AE0}
\definecolor{popgreen}{HTML}{2E9E6A}
\definecolor{poppink}{HTML}{E0457B}
\definecolor{popblue}{HTML}{2F7DE1}
\definecolor{popgold}{HTML}{C98A12}
\definecolor{popmuted}{HTML}{8A7F76}
\definecolor{popline}{HTML}{EADFD2}
\definecolor{popgray}{HTML}{8A7F76}
\colorlet{popgrey}{popgray}
% Alias tolérants (le modèle nomme parfois les couleurs différemment)
\colorlet{popwhite}{white}
\colorlet{popblack}{popink}
\colorlet{popred}{poppink}
\colorlet{popyellow}{popgold}
% Teintes claires (fonds de tableaux/encadrés) — le modèle nomme parfois
% les couleurs avec un suffixe L (ex. popblueL) sans les avoir définies.
\colorlet{poporangeL}{poporange!20}
\colorlet{popdarkL}{popdark!20}
\colorlet{poppurpleL}{poppurple!20}
\colorlet{popgreenL}{popgreen!20}
\colorlet{poppinkL}{poppink!20}
\colorlet{popblueL}{popblue!20}
\colorlet{popgoldL}{popgold!20}
\colorlet{popredL}{popred!20}
\colorlet{popgrayL}{popgray!20}
% Teintes claires pour fonds de tableaux / zones
\colorlet{popblueL}{popblue!10!white}
\colorlet{poporangeL}{poporange!14!white}
\colorlet{popgreenL}{popgreen!12!white}
\colorlet{poppurpleL}{poppurple!12!white}
\colorlet{poppinkL}{poppink!12!white}
\colorlet{popgoldL}{popgold!14!white}

\pagecolor{popcream}

% ============================================================
% Typographie
% ============================================================
\defaultfontfeatures{Ligatures=TeX}
\setmainfont{PlusJakartaSans-Medium.ttf}[Path=../../../fonts/,BoldFont=PlusJakartaSans-Bold.ttf]
% Symboles, API et emojis absents de la police principale : repli sur DejaVu Sans (caractères actifs)
\newfontface\symfont{DejaVuSans.ttf}[Path=../../../fonts/]
\makeatletter
\newcommand\symactive[1]{\catcode#1=13 \begingroup\lccode`\~=#1 \lowercase{\endgroup\def~}{{\symfont\char#1}}}
\newcommand\symblank[1]{\catcode#1=13 \begingroup\lccode`\~=#1 \lowercase{\endgroup\def~}{}}
\newcommand\symhyph[1]{\catcode#1=13 \begingroup\lccode`\~=#1 \lowercase{\endgroup\def~}{\mbox{-}}}
\newcount\symc
\newcommand\symrange[2]{\symc=#1 \@whilenum\symc<#2 \do{\expandafter\symactive\expandafter{\the\symc}\advance\symc by 1 }}
\newcommand\symblankrange[2]{\symc=#1 \@whilenum\symc<#2 \do{\expandafter\symblank\expandafter{\the\symc}\advance\symc by 1 }}
\makeatother
\symrange{"0250}{"02FF}\symactive{"2713}\symactive{"2714}\symactive{"2717}\symactive{"2718}\symactive{"2610}\symactive{"2611}\symactive{"2612}
\symrange{"2600}{"27BF}
\catcode"2705=13 \begingroup\lccode`\~="2705 \lowercase{\endgroup\def~}{{\symfont\char"2713}}\symblank{"26BD}\symblank{"26A1}
\symrange{"2460}{"257F}\symactive{"223C}\symactive{"2248}\symactive{"2260}
\symactive{"01F8}\symactive{"01F9}\symactive{"1E3E}\symactive{"1E3F}
\symhyph{"2011}\symhyph{"2E1A}\symblankrange{"1F300}{"1FAFF}\symblank{"FE0F}
% Caractères chinois : WenQuanYi Zen Hei (sous-ensemble GB2312 + pinyin), gras/italique simulés
\setCJKmainfont{WenQuanYiZenHei-subset.ttf}[Path=../../../fonts/,AutoFakeBold=3,AutoFakeSlant=0.2]
\xeCJKsetup{CJKspace=false}
% Pinyin : voyelles à caron (ǎ ǐ ǒ ǔ ǖ ǘ ǚ ǜ) absentes de la police principale -> caractères actifs
\newfontface\pyfont{WenQuanYiZenHei-subset.ttf}[Path=../../../fonts/]
\newcommand\pyactive[1]{\catcode#1=13 \begingroup\lccode`\~=#1 \lowercase{\endgroup\def~}{{\pyfont\char#1}}}
\pyactive{"01CD}\pyactive{"01CE}\pyactive{"01CF}\pyactive{"01D0}\pyactive{"01D1}\pyactive{"01D2}\pyactive{"01D3}\pyactive{"01D4}\pyactive{"01D5}\pyactive{"01D6}\pyactive{"01D7}\pyactive{"01D8}\pyactive{"01D9}\pyactive{"01DA}\pyactive{"01DB}\pyactive{"01DC}
% Maths en sans-serif (Fira Math) avec chiffres et lettres droites de Plus
% Jakarta, pour que les formules se fondent dans le texte.
\usepackage[math-style=ISO,bold-style=ISO]{unicode-math}
\setmathfont{FiraMath-Regular.otf}[Path=../../../fonts/]
\setmathfont{PlusJakartaSans-Medium.ttf}[Path=../../../fonts/,range={up/num,up/{Latin,latin},\mathup}]
\usepackage{icomma} % virgule décimale sans espace en mode math
% Caractères mathématiques Unicode absents des polices de texte (Plus Jakarta,
% Archivo Black) : rendus par leur équivalent mathématique, en texte comme en
% formule — sinon ils s'affichent en carré vide (ex. « Arithmétique dans ℤ »).
\usepackage{newunicodechar}
\newunicodechar{ℝ}{\ensuremath{\mathbb{R}}}
\newunicodechar{ℤ}{\ensuremath{\mathbb{Z}}}
\newunicodechar{ℂ}{\ensuremath{\mathbb{C}}}
\newunicodechar{ℕ}{\ensuremath{\mathbb{N}}}
\newunicodechar{ℚ}{\ensuremath{\mathbb{Q}}}
\newunicodechar{𝔻}{\ensuremath{\mathbb{D}}}
\newunicodechar{α}{\ensuremath{\alpha}}
\newunicodechar{β}{\ensuremath{\beta}}
\newunicodechar{γ}{\ensuremath{\gamma}}
\newunicodechar{δ}{\ensuremath{\delta}}
\newunicodechar{ε}{\ensuremath{\varepsilon}}
\newunicodechar{θ}{\ensuremath{\theta}}
\newunicodechar{λ}{\ensuremath{\lambda}}
\newunicodechar{μ}{\ensuremath{\mu}}
\newunicodechar{σ}{\ensuremath{\sigma}}
\newunicodechar{τ}{\ensuremath{\tau}}
\newunicodechar{φ}{\ensuremath{\varphi}}
\newunicodechar{ψ}{\ensuremath{\psi}}
\newunicodechar{ω}{\ensuremath{\omega}}
\newunicodechar{Σ}{\ensuremath{\Sigma}}
% majuscules produites par \MakeUppercase dans les titres (α-aminés -> Α)
\newunicodechar{Α}{\ensuremath{\alpha}}
\newunicodechar{Β}{\ensuremath{\beta}}
\newunicodechar{Γ}{\ensuremath{\Gamma}}
\newunicodechar{Δ}{\ensuremath{\Delta}}
\newunicodechar{Ε}{\ensuremath{\varepsilon}}
\newunicodechar{Θ}{\ensuremath{\Theta}}
\newunicodechar{Λ}{\ensuremath{\Lambda}}
\newunicodechar{Μ}{\ensuremath{\mu}}
\newunicodechar{Τ}{\ensuremath{\tau}}
\newunicodechar{Φ}{\ensuremath{\Phi}}
\newunicodechar{Ψ}{\ensuremath{\Psi}}
\newunicodechar{Ω}{\ensuremath{\Omega}}
\newunicodechar{↦}{\ensuremath{\mapsto}}
\newunicodechar{∈}{\ensuremath{\in}}
\newunicodechar{∉}{\ensuremath{\notin}}
\newunicodechar{∧}{\ensuremath{\wedge}}
\newunicodechar{⇔}{\ensuremath{\Leftrightarrow}}
\newunicodechar{⟺}{\ensuremath{\Longleftrightarrow}}
\newunicodechar{⇒}{\ensuremath{\Rightarrow}}
\newunicodechar{⟹}{\ensuremath{\Longrightarrow}}
\newunicodechar{∘}{\ensuremath{\circ}}
\newunicodechar{∪}{\ensuremath{\cup}}
\newunicodechar{∩}{\ensuremath{\cap}}
\newunicodechar{≡}{\ensuremath{\equiv}}
\newunicodechar{⊂}{\ensuremath{\subset}}
\newunicodechar{⊥}{\ensuremath{\perp}}
\newunicodechar{∃}{\ensuremath{\exists}}
\newunicodechar{∀}{\ensuremath{\forall}}
\newunicodechar{∛}{\ensuremath{\sqrt[3]{\,}}}
\newunicodechar{∜}{\ensuremath{\sqrt[4]{\,}}}
\newunicodechar{⃗}{\ensuremath{{}^{\to}}}
\newunicodechar{ⁿ}{\ifmmode^{n}\else\textsuperscript{\ensuremath{n}}\fi}
\newunicodechar{ˣ}{\ifmmode^{x}\else\textsuperscript{\ensuremath{x}}\fi}
\newunicodechar{ᵇ}{\ifmmode^{b}\else\textsuperscript{\ensuremath{b}}\fi}
\newunicodechar{ʳ}{\ifmmode^{r}\else\textsuperscript{\ensuremath{r}}\fi}
\newunicodechar{ᵘ}{\ifmmode^{u}\else\textsuperscript{\ensuremath{u}}\fi}
\newunicodechar{ʸ}{\ifmmode^{y}\else\textsuperscript{\ensuremath{y}}\fi}
\newunicodechar{ᵃ}{\ifmmode^{a}\else\textsuperscript{\ensuremath{a}}\fi}
\newunicodechar{ᵉ}{\ifmmode^{e}\else\textsuperscript{\ensuremath{e}}\fi}
\newunicodechar{ᵏ}{\ifmmode^{k}\else\textsuperscript{\ensuremath{k}}\fi}
\newunicodechar{ᵖ}{\ifmmode^{p}\else\textsuperscript{\ensuremath{p}}\fi}
\newunicodechar{ᴺ}{\ifmmode^{N}\else\textsuperscript{\ensuremath{N}}\fi}
\newunicodechar{ᶜ}{\ifmmode^{c}\else\textsuperscript{\ensuremath{c}}\fi}
\newunicodechar{ʰ}{\ifmmode^{h}\else\textsuperscript{\ensuremath{h}}\fi}
\newunicodechar{ᵠ}{\ifmmode^{q}\else\textsuperscript{\ensuremath{q}}\fi}
\newunicodechar{ₙ}{\ifmmode_{n}\else\textsubscript{\ensuremath{n}}\fi}
\newunicodechar{ₐ}{\ifmmode_{a}\else\textsubscript{\ensuremath{a}}\fi}
\newunicodechar{ᵢ}{\ifmmode_{i}\else\textsubscript{\ensuremath{i}}\fi}
\newunicodechar{ᵣ}{\ifmmode_{r}\else\textsubscript{\ensuremath{r}}\fi}
\newunicodechar{ₖ}{\ifmmode_{k}\else\textsubscript{\ensuremath{k}}\fi}
\newunicodechar{ᵦ}{\ifmmode_{\beta}\else\textsubscript{\ensuremath{\beta}}\fi}
\newunicodechar{ₓ}{\ifmmode_{x}\else\textsubscript{\ensuremath{x}}\fi}
\newfontfamily\popdisplay{ArchivoBlack-Regular.ttf}[Path=../../../fonts/]
\newfontfamily\poplogo{SpaceGrotesk-Bold.ttf}[Path=../../../fonts/]
\newfontfamily\popsemi{PlusJakartaSans-SemiBold.ttf}[Path=../../../fonts/]
\newfontfamily\popxbold{PlusJakartaSans-ExtraBold.ttf}[Path=../../../fonts/]
\newfontfamily\popxboldls{PlusJakartaSans-ExtraBold.ttf}[Path=../../../fonts/,LetterSpace=3.0]

\setlist[enumerate]{leftmargin=1.7em,itemsep=5pt,topsep=4pt}
\setlist[itemize]{leftmargin=1.7em,itemsep=4pt,topsep=4pt,label={\color{poporange}\textbullet}}
\setlength{\parindent}{0pt}
\setlength{\parskip}{5pt}
\renewcommand{\arraystretch}{1.25}

% pgfplots : style Pop par défaut
\pgfplotsset{
  every axis/.append style={axis line style={popink,line width=1pt},tick style={popink},
    grid style={popline,line width=0.5pt},label style={font=\small},tick label style={font=\footnotesize}},
  popcurve/.style={poporange,line width=1.8pt,no markers,smooth},
}
% Même style disponible dans un \draw TikZ simple (hors pgfplots)
\tikzset{popcurve/.style={poporange,line width=1.8pt}}
\tikzset{popgrid/.style={popline,very thin}}
\colorlet{popcurve}{poporange} % le nom du style est parfois utilisé comme couleur

% ============================================================
% Pill / squiggle
% ============================================================
\newcommand{\poppill}[3]{%
  \tikz[baseline=-0.55ex]\node[draw=popink,line width=1.2pt,rounded corners=2.6pt,%
    fill=#2,inner xsep=6.5pt,inner ysep=3pt]{%
    \popxboldls\fontsize{7}{7}\selectfont\color{#3}\MakeUppercase{#1}};%
}
\newcommand{\popsquiggle}[1]{%
  \tikz[baseline=-0.4ex]{
    \draw[#1,line width=2.6pt,line cap=round]
      plot [smooth] coordinates {(0,0.09) (0.34,0.01) (0.68,0.21) (1.02,0.01) (1.36,0.10)};
  }%
}
% Mot-clé surligné (marqueur orange sous le texte)
\newcommand{\cle}[1]{{\popxbold\color{popdark}#1}}

% ============================================================
% En-tête / pied
% ============================================================
\pagestyle{fancy}
\fancyhf{}
\renewcommand{\headrulewidth}{0pt}
\renewcommand{\footrulewidth}{0pt}
\lhead{\raisebox{-0.15ex}{\poplogo\fontsize{13}{13}\selectfont\color{popink}OnBuch\color{poporange}+}}
\rhead{\poppill{\DOCMATIERE\ \textbullet\ \DOCNIVEAU\ \textbullet\ Leçon \DOCLECON}{white}{popink}}
\lfoot{\popxbold\fontsize{7.6}{7.6}\selectfont\color{popmuted}\MakeUppercase{Cours signé OnBuch+}\ \textbullet\ {\popsemi\DOCTITRE}}
\rfoot{\tikz[baseline=-0.6ex]\node[rounded corners=8.5pt,fill=popink,minimum height=17pt,minimum width=17pt,inner xsep=5pt,inner sep=0]{\popxbold\fontsize{8}{8}\selectfont\color{popcream}\thepage\,/\,\pageref*{LastPage}};}

% ============================================================
% Titres
% ============================================================
\newcommand{\chaptitre}[2]{%
  \begin{tcolorbox}[enhanced,colback=poporange,colframe=popink,boxrule=2.6pt,arc=11pt,
      drop shadow=popink,left=15pt,right=15pt,top=12pt,bottom=11pt,before skip=3pt,after skip=18pt]
    \poppill{#2}{popink}{popcream}\par\vspace{9pt}
    {\raggedright\hyphenpenalty=10000\exhyphenpenalty=10000\popdisplay\fontsize{22}{24}\selectfont\color{popcream}\MakeUppercase{#1}\par}
  \end{tcolorbox}
}
% Section numérotée : pastille orange avec le numéro + titre Archivo
\newcounter{popsec}
\newcommand{\coursec}[1]{%
  \stepcounter{popsec}\setcounter{popsub}{0}%
  \par\vspace{16pt}\noindent
  \begin{minipage}{\linewidth}
  \tikz[baseline=-0.7ex]\node[draw=popink,line width=1.6pt,fill=poporange,rounded corners=4pt,minimum size=22pt,inner sep=0]{\popdisplay\fontsize{12}{12}\selectfont\color{popcream}\thepopsec};%
  \hspace{8pt}{\popdisplay\fontsize{15}{17}\selectfont\color{popink}\MakeUppercase{#1}}\par
  \vspace{4pt}\noindent{\color{poporange}\rule{36pt}{3.6pt}}
  \end{minipage}\par\nopagebreak\vspace{8pt}
}
\newcounter{popsub}
\newcommand{\courssub}[1]{%
  \stepcounter{popsub}%
  \par\vspace{9pt}\noindent{\popxbold\fontsize{11.5}{13}\selectfont\color{popink}{\color{poporange}\thepopsec.\thepopsub}\hspace{6pt}#1}\par\nopagebreak\vspace{4pt}
}

% ============================================================
% Boîtes pédagogiques — même recette : fond blanc, bordure épaisse colorée,
% ombre dure encre, badge plein. Titre optionnel [#1].
% ============================================================
\tcbset{popbase/.style={breakable,enhanced,colback=white,boxrule=2.3pt,arc=9pt,
  shadow={3pt}{-3pt}{0pt}{popink},left=14pt,right=14pt,top=11pt,bottom=11pt,before skip=11pt,after skip=15pt,
  fontupper=\popsemi\fontsize{10.6}{14.5}\selectfont\color{popink},
  fontlower=\popsemi\fontsize{10.6}{14.5}\selectfont\color{popink}}}
% Coche dessinée (le glyphe ✓ n'existe pas dans les polices du document)
\newcommand{\popcheck}[1]{\tikz[baseline=-0.5ex]\draw[#1,line width=1.6pt,line cap=round,line join=round](-0.13,0.0)--(-0.04,-0.09)--(0.13,0.1);}
\providecommand{\coche}{\popcheck{popgreen}}
\providecommand{\resultat}[1]{\par\smallskip\noindent\fcolorbox{popgreen}{popgreenL}{\parbox{\dimexpr\linewidth-2\fboxsep-2\fboxrule\relax}{\textbf{Résultat.} #1}}\par\smallskip}
\newcommand{\popboxhead}[3]{\poppill{#1}{#2}{white}\ifx\relax#3\relax\else\hspace{6pt}{\popxbold\fontsize{10.6}{12}\selectfont\color{popink}#3}\fi\par\vspace{7pt}}

\newtcolorbox{aretenir}[1][]{popbase,colframe=popgreen,before upper={\popboxhead{À retenir}{popgreen}{#1}}}
% Texte en chinois (document de lecture, dialogue, extrait) : cadre
% d'étude. Un titre optionnel (source, auteur) s'affiche dans l'en-tête.
\newtcolorbox{textebox}[1][]{popbase,colframe=popblue,before upper={\popboxhead{文本 · Texte}{popblue}{#1}},after upper={}}
% Marques ✓ / ✗ (forme correcte / fautive) souvent utilisées par les modèles
\providecommand{\cross}{\textcolor{poppink}{\ensuremath{\times}}}
\providecommand{\xmark}{\cross}\providecommand{\crossmark}{\cross}\providecommand{\cmark}{\ensuremath{\checkmark}}
% Ligne de réponse / justification à compléter (inventée par les modèles)
\providecommand{\justification}[1]{\par\noindent\textit{Justification~:} \dotfill\par}
% \textipa (package tipa non chargé) : la police ne couvre pas l'API -> simple passage
\providecommand{\textipa}[1]{#1}
% Soulignement (ulem non chargé) : \uline, \ul
\providecommand{\uline}[1]{\underline{#1}}\providecommand{\ul}[1]{\underline{#1}}
% \glossaire{\item ...} inventée par les modèles : liste de glossaire
\providecommand{\glossaire}[1]{\begin{itemize}#1\end{itemize}}
% \alert (beamer) inventée par les modèles : texte mis en évidence
\providecommand{\alert}[1]{\textbf{\textcolor{poppink}{#1}}}
% variantes ASCII d'ordinaux écrites en macro
\providecommand{\iere}{\textsuperscript{re}}\providecommand{\ieme}{\textsuperscript{e}}\providecommand{\ere}{\textsuperscript{re}}\providecommand{\eme}{\textsuperscript{e}}\providecommand{\er}{\textsuperscript{er}}\providecommand{\ie}{\textsuperscript{e}}
% Ordinaux français écrits en macro par les modèles : 1\ière, 3\ième
\newcommand{\ière}{\textsuperscript{re}}\newcommand{\ième}{\textsuperscript{e}}
% \exemple (inventée par les modèles) : étiquette « Exemple : »
\providecommand{\exemple}{\textbf{Exemple~:}\ }
% \square utilisable aussi hors mode math (cases à cocher)
\AtBeginDocument{\let\squareMath\square\renewcommand{\square}{\ensuremath{\squareMath}}}
% Mot ou phrase en chinois à l'intérieur d'un paragraphe français
\newcommand{\zh}[1]{\ifmmode\text{#1}\else{#1}\fi}
\let\eng\zh \let\esp\zh \let\all\zh % alias tolérés
% flèches d'intonation inventées par les modèles
\providecommand{\textsearrow}{}\renewcommand{\textsearrow}{\ensuremath{\searrow}}\providecommand{\textnearrow}{}\renewcommand{\textnearrow}{\ensuremath{\nearrow}}
% \fr{..} : mot français (inventé par les modèles)
\providecommand{\fr}[1]{\textit{#1}}
% zhbox : encadré de texte chinois inventé par les modèles
\newenvironment{zhbox}{\begin{quote}}{\end{quote}}
% \courssubsub : titre de niveau 3 inventé par les modèles
\providecommand{\courssubsub}[1]{\par\medskip\noindent\textbf{#1}\par\smallskip}
% \vs : « versus » inventé par les modèles
\providecommand{\vs}{\textit{vs}\ }
% \blank : case à compléter inventée par les modèles
\providecommand{\blank}[1][]{\underline{\hspace{1.6em}}}
\newtcolorbox{exemplebox}[1][]{popbase,colframe=poporange,before upper={\popboxhead{Exemple}{poporange}{#1}}}
\newtcolorbox{definition}[1][]{popbase,colframe=popblue,before upper={\popboxhead{Définition}{popblue}{#1}}}
\newtcolorbox{propriete}[1][]{popbase,colframe=poppurple,before upper={\popboxhead{Propriété}{poppurple}{#1}}}
\newtcolorbox{methode}[1][]{popbase,colframe=popgold,before upper={\popboxhead{Méthode}{popgold}{#1}}}
\newtcolorbox{attention}[1][]{popbase,colframe=poppink,before upper={\popboxhead{Attention}{poppink}{#1}}}
\newtcolorbox{experience}[1][]{popbase,colframe=popblue,colback=popblueL,before upper={\popboxhead{Expérience}{popblue}{#1}}}
% « Le savais-tu ? » : carte violette pleine (contexte, culture, Cameroun)
\newtcolorbox{savaistu}[1][]{popbase,colback=poppurple,colframe=popink,
  fontupper=\popsemi\fontsize{10.4}{14.2}\selectfont\color{white},
  before upper={\poppill{Le savais-tu\,?}{white}{poppurple}\ifx\relax#1\relax\else\hspace{6pt}{\popxbold\color{white}#1}\fi\par\vspace{7pt}}}
% Exercice résolu : énoncé en haut, \tcblower puis la solution sur fond crème
\newcounter{popexo}\newcounter{popexr}
\newtcolorbox{exoresolu}[1][]{popbase,colframe=poporange,colbacklower=popcream,
  segmentation style={popink,line width=1.2pt,dashed},
  before upper={\stepcounter{popexr}\popboxhead{Exercice résolu \thepopexr}{poporange}{#1}},
  before lower={\poppill{Solution}{popink}{popcream}\par\vspace{6pt}}}
% Exercice (non résolu en place) — la correction est dans la partie Corrigés
\newtcolorbox{exercice}[1][]{popbase,colframe=popink,
  before upper={\stepcounter{popexo}\popboxhead{Exercice \thepopexo}{popink}{#1}}}
% Corrigé
\newtcolorbox{corrige}[1][]{popbase,colframe=popgreen,colback=popgreenL,
  before upper={\popboxhead{Corrigé}{popgreen}{#1}}}
% Objectifs / prérequis / bilan
\newtcolorbox{objectifs}{popbase,colframe=popink,colback=poporangeL,
  before upper={\popboxhead{Objectifs de la leçon}{popink}{}}}
\newtcolorbox{prerequis}{popbase,colframe=popink,colback=popblueL,
  before upper={\popboxhead{Prérequis}{popblue}{}}}
\newtcolorbox{bilan}{popbase,colframe=popink,colback=white,boxrule=2.8pt,
  before upper={\popboxhead{Fiche bilan}{poporange}{L'essentiel à connaître pour l'examen}}}
% Figure encadrée avec légende
\newtcolorbox{popfigure}[1][]{enhanced,breakable=false,colback=white,colframe=popline,boxrule=1.4pt,arc=8pt,
  left=10pt,right=10pt,top=10pt,bottom=8pt,before skip=10pt,after skip=12pt,halign=center,
  lower separated=false,halign lower=center,
  fontlower=\popsemi\fontsize{9}{11.5}\selectfont\color{popmuted},
  #1}
\newcommand{\legende}[1]{\par\vspace{6pt}{\popsemi\fontsize{9}{11.5}\selectfont\color{popmuted}#1}}

% ============================================================
% Signature OnBuch+
% ============================================================
\newcommand{\poplogoinline}[1]{{\poplogo\fontsize{#1}{#1}\selectfont\color{popink}OnBuch\color{poporange}+}}
\newcommand{\signatureonbuch}{%
  \begin{tcolorbox}[enhanced,colback=popink,colframe=popink,boxrule=2.2pt,arc=10pt,
      drop shadow=poporange,left=14pt,right=14pt,top=10pt,bottom=10pt,before skip=0pt,after skip=0pt]
    \tikz[baseline=-0.7ex]\node[circle,fill=popgreen,draw=popcream,line width=1.3pt,minimum size=22pt,inner sep=0]{\popcheck{white}};%
    \hspace{9pt}\begin{minipage}[c]{0.62\linewidth}
      {\popxbold\fontsize{12}{13}\selectfont\color{popcream}Signé\ \textbullet\ {\poplogo OnBuch\color{poporange}+}}\par\vspace{2pt}
      {\raggedright\popsemi\fontsize{8}{10.5}\selectfont\color{popline}\MakeUppercase{Équipe pédagogique OnBuch+}\\\MakeUppercase{Conforme au programme officiel MINESEC}\par}
    \end{minipage}\hfill
    {\poplogo\fontsize{22}{22}\selectfont\color{popcream}OnBuch\color{poporange}+}
  \end{tcolorbox}%
}

% ============================================================
% Page de garde
% #1 titre · #2 matière · #3 niveau · #4 module · #5 n° leçon · #6 durée ·
% #7 description / objectifs (texte ou itemize)
% ============================================================
\newcommand{\pagegarde}[7]{%
  \thispagestyle{empty}
  \newgeometry{a4paper,margin=1.7cm,top=1.6cm,bottom=1.4cm}%
  \begin{tikzpicture}[remember picture,overlay]
    \fill[poporange!18] ([xshift=-3.2cm,yshift=-3.6cm]current page.north east) circle (4.6cm);
    \fill[poppurple!14] ([xshift=2.4cm,yshift=7.6cm]current page.south west) circle (3.8cm);
  \end{tikzpicture}%
  {\poplogo\fontsize{30}{30}\selectfont\color{popink}OnBuch\color{poporange}+}\hfill
  \raisebox{0.6ex}{\poppill{Cours signé \textbullet\ Premium}{popink}{popcream}}\par
  \vspace{1pt}\popsquiggle{poppurple}\par
  \vspace{8pt}
  \begin{tcolorbox}[enhanced,colback=poporange,colframe=popink,boxrule=2.8pt,arc=13pt,
      drop shadow=popink,left=18pt,right=18pt,top=12pt,bottom=13pt,after skip=10pt]
    \poppill{#2 \textbullet\ #3}{popink}{popcream}\hspace{5pt}\poppill{#4}{white}{popink}\par\vspace{8pt}
    {\popdisplay\fontsize{14}{14}\selectfont\color{popink}LEÇON #5}\par\vspace{4pt}
    {\raggedright\hyphenpenalty=10000\exhyphenpenalty=10000\popdisplay\fontsize{25}{27}\selectfont\color{popcream}\MakeUppercase{#1}\par}
    \vspace{8pt}
    {\popxbold\fontsize{9.5}{9.5}\selectfont\color{popink}\MakeUppercase{Durée conseillée : #6}}
  \end{tcolorbox}
  \begin{tcolorbox}[enhanced,colback=white,colframe=popink,boxrule=2.2pt,arc=10pt,
      drop shadow=popink,left=16pt,right=16pt,top=10pt,bottom=10pt,after skip=8pt]
    \poppill{Dans cette leçon}{poppurple}{white}\par\vspace{5pt}
    \popsemi\fontsize{9.3}{12.6}\selectfont\color{popink}#7
  \end{tcolorbox}
  \noindent\begin{minipage}[t]{0.487\linewidth}
    \begin{tcolorbox}[enhanced,colback=popgreen,colframe=popink,boxrule=2.2pt,arc=10pt,
        drop shadow=popink,left=11pt,right=11pt,top=10pt,bottom=10pt,height=3.1cm,valign=center]
      \tcbox[colback=white,colframe=popink,boxrule=1.3pt,arc=5pt,boxsep=0pt,nobeforeafter,
        left=3pt,right=3pt,top=3pt,bottom=3pt,valign=center]{\qrcode[height=1.9cm]{https://play.google.com/store/apps/details?id=com.luvvix.onbuch&hl=fr}}%
      \hspace{8pt}\begin{minipage}[c]{0.5\linewidth}
        {\popdisplay\fontsize{11}{12.5}\selectfont\color{white}\MakeUppercase{Télécharge OnBuch}}\par\vspace{4pt}
        {\raggedright\popsemi\fontsize{8.4}{11}\selectfont\color{white}Gratuit \textbullet\ Google Play\\Tuteur IA, annales, concours, résultats\par}
      \end{minipage}
    \end{tcolorbox}
  \end{minipage}\hfill
  \begin{minipage}[t]{0.487\linewidth}
    \begin{tcolorbox}[enhanced,colback=poppurple,colframe=popink,boxrule=2.2pt,arc=10pt,
        drop shadow=popink,left=11pt,right=11pt,top=10pt,bottom=10pt,height=3.1cm,valign=center]
      \tikz[baseline=-0.7ex]\node[circle,fill=white,draw=popink,line width=1.3pt,minimum size=1.6cm,inner sep=0]{\popdisplay\fontsize{24}{24}\selectfont\color{poppurple}+};%
      \hspace{8pt}\begin{minipage}[c]{0.6\linewidth}
        {\popdisplay\fontsize{11}{12.5}\selectfont\color{white}\MakeUppercase{Passe à OnBuch+}}\par\vspace{4pt}
        {\raggedright\popsemi\fontsize{8.4}{11}\selectfont\color{white}Tous les cours signés, Tuteur IA illimité, fiches PDF — onglet OnBuch+ de l'app\par}
      \end{minipage}
    \end{tcolorbox}
  \end{minipage}
  \vspace{8pt}
  \popcontenu
  \vfill
  \signatureonbuch
  \restoregeometry
  \newpage
}

% Rangée « Ce cours contient » (page de garde)
\newcommand{\poptile}[3]{% #1 couleur #2 gros texte #3 légende
  \begin{tcolorbox}[enhanced,colback=#1,colframe=popink,boxrule=2pt,arc=9pt,shadow={2.5pt}{-2.5pt}{0pt}{popink},
      left=6pt,right=6pt,top=9pt,bottom=9pt,halign=center,valign=center,height=1.75cm,width=\linewidth,nobeforeafter]
    {\popdisplay\fontsize{12.5}{13}\selectfont\color{white}\MakeUppercase{#2}}\par\vspace{3pt}
    {\centering\popsemi\fontsize{7.8}{9.5}\selectfont\color{white}#3\par}
  \end{tcolorbox}}
\newcommand{\popcontenu}{%
  \par\noindent{\popxboldls\fontsize{7.5}{7.5}\selectfont\color{popmuted}CE COURS SIGNÉ CONTIENT}\par\vspace{7pt}
  \noindent\begin{minipage}[t]{0.235\linewidth}\poptile{popblue}{Cours}{complet, illustré, pas à pas}\end{minipage}\hfill
  \begin{minipage}[t]{0.235\linewidth}\poptile{poporange}{Méthodes}{et exercices résolus}\end{minipage}\hfill
  \begin{minipage}[t]{0.235\linewidth}\poptile{poppink}{Exercices}{type examen + corrigés}\end{minipage}\hfill
  \begin{minipage}[t]{0.235\linewidth}\poptile{popgreen}{Fiche bilan}{l'essentiel à retenir}\end{minipage}\par}

% Sommaire
\newtcolorbox{sommaire}{enhanced,colback=white,colframe=popink,boxrule=2.2pt,arc=10pt,
  drop shadow=popink,left=18pt,right=18pt,top=13pt,bottom=13pt,before skip=6pt,after skip=20pt,
  before upper={\poppill{Sommaire}{poppurple}{white}\par\vspace{9pt}\popsemi\fontsize{10.6}{14.5}\selectfont\color{popink}}}

% Signature de fin de cours — poussée tout en bas de la dernière page
\newcommand{\findecours}{%
  \par\vfill\nopagebreak
  \begin{center}\popsquiggle{poporange}\end{center}\vspace{4pt}
  \signatureonbuch
}
\AtBeginDocument{\def\llbracket{\mathopen{[\mkern-3.5mu[}}\def\rrbracket{\mathclose{]\mkern-3.5mu]}}} % intervalles d'entiers (glyphe ⟦ absent de la police)
% Glyphes absents de Fira Math : redéfinis à partir de glyphes disponibles
\AtBeginDocument{%
  \def\setminus{\mathbin{\backslash}}\let\smallsetminus\setminus
  \def\vdots{\mathord{\vbox{\baselineskip3.2pt\lineskiplimit0pt\kern4pt\hbox{.}\hbox{.}\hbox{.}}}}%
  \def\ddots{\mathinner{\mkern1mu\raise6.5pt\hbox{.}\mkern2mu\raise3.6pt\hbox{.}\mkern2mu\raise0.7pt\hbox{.}\mkern1mu}}%
  \def\triangle{\mathord{\scalebox{0.9}{$\symup{\Delta}$}}}%
  \def\nabla{\mathord{\raisebox{\depth}{\scalebox{1}[-1]{$\symup{\Delta}$}}}}%
  \def\checkmark{\popcheck{popdark}}%
  \def\star{\mathbin{\ast}}\def\bigstar{\mathord{\ast}}%
  \def\diamond{\mathbin{\scalebox{0.6}{$\Diamond$}}}%
  \def\bigcup{\mathop{\mathchoice{\raisebox{-0.3em}{\scalebox{1.7}{$\cup$}}}{\scalebox{1.25}{$\cup$}}{\cup}{\cup}}\displaylimits}%
  \def\bigcap{\mathop{\mathchoice{\raisebox{-0.3em}{\scalebox{1.7}{$\cap$}}}{\scalebox{1.25}{$\cap$}}{\cap}{\cap}}\displaylimits}%
  \def\lesseqqgtr{\mathrel{\lessgtr}}\def\bigoplus{\mathop{\oplus}\displaylimits}%
}
\providecommand{\ssi}{si et seulement si} % abréviation fréquente
\AtBeginDocument{\providecommand{\wideparen}{\overparen}} % arc de cercle

%% FIGFIX-BEGIN v4 — correctifs globaux des figures (docs/onbuch-generation/tools/figfix)
\usepackage{adjustbox}\usepackage{etoolbox}
% toute figure TikZ/pgfplots plus large que la ligne est réduite à la largeur disponible
\BeforeBeginEnvironment{tikzpicture}{\begin{adjustbox}{max width=\linewidth}}
\AfterEndEnvironment{tikzpicture}{\end{adjustbox}}
% légendes pgfplots lisibles par-dessus les courbes
\pgfplotsset{every axis/.append style={legend style={fill=white,fill opacity=0.92,text opacity=1,draw=black!15,inner sep=3pt}}}
% étiquettes d'arêtes (syntaxe edge["..."]) avec fond blanc
\tikzset{every edge quotes/.append style={fill=white,inner sep=1.2pt}}
%% FIGFIX-END
\begin{document}

\pagegarde{Leçon 32 — Expression écrite : affiche publicitaire sur un produit (marketing)}{Chinois}{Tle A}{Module 4 : Activités économiques et monde du travail}{ZH32}{2 h}{Cette leçon apprend à rédiger une affiche publicitaire efficace en chinois : analyse de modèle, structures persuasives, écriture des caractères clés, exercices de thème/version et grille d'évaluation type Bac.
\vspace{4pt}
\begin{multicols}{2}\small
\begin{itemize}
\item Identifier la structure type d'une affiche publicitaire chinoise (slogan, arguments, appel à l'action).
\item Utiliser les connecteurs logiques et rhétoriques adaptés au marketing (不但…而且…, 只要…就…, 首先、此外、最重要的是).
\item Écrire correctement les 12 caractères clés du thème (traits, radicaux, confusions fréquentes).
\item Traduire du français vers le chinois des phrases publicitaires en respectant l'ordre des mots et le registre.
\item Traduire du chinois vers le français un texte publicitaire en rendant les nuances persuasives.
\item Appliquer une grille d'évaluation (vocabulaire, grammaire, caractères, mise en page, créativité) pour s'auto-corriger.
\end{itemize}\end{multicols}}

\begin{sommaire}
\begin{multicols}{2}
\begin{enumerate}
\item 1. Analyse d'un modèle d'affiche commenté
\item 2. Structures et connecteurs pour la persuasion
\item 3. Écriture correcte des caractères du thème
\item 4. Exercices de thème (français → chinois)
\item 5. Exercices de version (chinois → français)
\item 6. Grille d'évaluation et production finale
\item Activité d'intégration
\item Méthodes et exercices résolus
\item Exercices
\item Corrigés des exercices
\item Fiche bilan
\end{enumerate}
\end{multicols}
\end{sommaire}

\begin{objectifs}
\begin{itemize}
\item Identifier la structure type d'une affiche publicitaire chinoise (slogan, arguments, appel à l'action).
\item Utiliser les connecteurs logiques et rhétoriques adaptés au marketing (不但…而且…, 只要…就…, 首先、此外、最重要的是).
\item Écrire correctement les 12 caractères clés du thème (traits, radicaux, confusions fréquentes).
\item Traduire du français vers le chinois des phrases publicitaires en respectant l'ordre des mots et le registre.
\item Traduire du chinois vers le français un texte publicitaire en rendant les nuances persuasives.
\item Appliquer une grille d'évaluation (vocabulaire, grammaire, caractères, mise en page, créativité) pour s'auto-corriger.
\item Produire une affiche bilingue chinois-français pour un produit camerounais ciblant le marché chinois.
\end{itemize}
\end{objectifs}

\begin{prerequis}
\begin{itemize}
\item Vocabulaire de base sur les produits et le commerce (produit, prix, qualité, acheter, vendre) acquis en Première.
\item Maîtrise des structures comparatives (比, 没有, 更, 最) et causales (因为…所以…, 由于…因此…).
\item Connaissance des connecteurs d'énumération et de progression (首先、其次、最后、而且).
\item Capacité à écrire environ 300 caractères courants et à en déduire l'ordre des traits.
\item Notions de base sur la culture de consommation en Chine (fêtes promotionnelles, codes QR, WeChat).
\end{itemize}
\end{prerequis}

\newpage

\coursec{Analyse d'un modèle d'affiche commenté}

Dans les rues de Yaoundé ou de Douala, vous croisez chaque jour des panneaux vantant le \emph{Poivre de Penja}, l'\emph{Eau minérale du Mont Cameroun} ou le \emph{Café robusta de l'Ouest}. Ces affiches parlent aux passants camerounais. Imaginez maintenant la situation inverse : une coopérative camerounaise vous demande de concevoir une affiche \textbf{en chinois} pour attirer les touristes de Shanghai ou de Guangzhou vers ces produits. Comment accrocher le regard en quelques caractères ? Quels arguments mettre en avant ? Quelle formule pour pousser à l'achat ?

\textbf{Problématique :} quels sont les ingrédients linguistiques et culturels d'une affiche publicitaire chinoise efficace ? Pour y répondre, nous allons d'abord \emph{observer} un modèle authentique, puis \emph{analyser} ses quatre zones, son vocabulaire, ses connecteurs et son registre, avant de le comparer aux habitudes publicitaires camerounaises.

\courssub{Observation : une affiche originale}

Lisons d'abord l'affiche sans chercher à tout comprendre. Repérez simplement : les caractères en grand, les phrases courtes, la présence d'un appel final.

\begin{textebox}[Affiche : 云南绿茶 — Thé vert du Yunnan]
绿色有机，纯天然！\\
Lǜsè yǒujī, chún tiānrán!\\
\emph{Vert et biologique, 100 \% naturel !}\\[4pt]
不但口感鲜爽，而且营养丰富。\\
Bùdàn kǒugǎn xiānshuǎng, érqiě yíngyǎng fēngfù.\\
\emph{Non seulement son goût est frais et vif, mais il est aussi très nutritif.}\\[4pt]
只要扫码下单，即享限时优惠。\\
Zhǐyào sǎomǎ xiàdān, jí xiǎng xiànshí yōuhuì.\\
\emph{Il suffit de scanner le code pour commander et profiter immédiatement de l'offre limitée.}\\[4pt]
与其犹豫，不如立即尝试！\\
Yǔqí yóuyù, bùrú lìjí chángshì!\\
\emph{Plutôt que d'hésiter, essayez tout de suite !}\\[4pt]
—— 云南绿茶，最纯正的选择\\
—— Yúnnán lǜchá, zuì chúnzhèng de xuǎnzé\\
\emph{—— Le thé vert du Yunnan, le choix le plus authentique.}
\end{textebox}

\textbf{Questions de compréhension :}
\begin{enumerate}
\item Quel produit est présenté et d'où vient-il ?
\item Quels sont les deux avantages du produit mentionnés dans la deuxième phrase ?
\item Que doit faire le client pour bénéficier de la réduction ?
\item Quelle est la dernière phrase et quel rôle joue-t-elle ?
\end{enumerate}

\begin{corrige}[Questions de compréhension]
\begin{enumerate}
\item Il s'agit du thé vert du Yunnan (\zh{云南绿茶}), une province du sud-ouest de la Chine célèbre pour son thé.
\item Son goût frais et vif (\zh{口感鲜爽}) et sa richesse nutritionnelle (\zh{营养丰富}).
\item Il doit scanner le code QR et passer commande (\zh{扫码下单}).
\item C'est la signature-slogan finale : \zh{最纯正的选择}, « le choix le plus authentique », qui résume la promesse de la marque.
\end{enumerate}
\end{corrige}

\courssub{Les quatre zones d'une affiche chinoise}

Une affiche publicitaire chinoise, même très courte, est organisée en \textbf{quatre zones fonctionnelles}. Apprenez à les repérer : c'est le squelette de toute production écrite de ce type.

\begin{center}
\begin{tabularx}{\linewidth}{|X|X|X|X|}
\hline
\rowcolor{popblueL} Zone & Fonction & Exemple de l'affiche & Traduction \\
\hline
1. Accroche (slogan) & Attirer l'œil, mémoriser & \zh{绿色有机，纯天然！} & Vert et bio, 100 \% naturel ! \\
\hline
2. Description du produit & Dire ce que c'est, ses qualités & \zh{不但口感鲜爽，而且营养丰富。} & Goût frais et nutritif \\
\hline
3. Arguments de vente & Donner une raison d'acheter & \zh{只要扫码下单，即享限时优惠。} & Offre limitée par QR code \\
\hline
4. Appel à l'action & Pousser à agir maintenant & \zh{与其犹豫，不如立即尝试！} & Essayez tout de suite ! \\
\hline
\end{tabularx}
\end{center}

\begin{definition}[Le slogan : 标语]
Le \cle{slogan} (\zh{标语} \emph{biāoyǔ}) est une phrase courte, percutante, souvent rythmée ou rimée, placée en haut de l'affiche. En chinois, il exploite volontiers les \textbf{parallélismes de syllabes} : \zh{绿色有机} (4 syllabes) + \zh{纯天然} (3 syllabes) crée un rythme sec et mémorable. Le chinois, langue monosyllabique à tons, se prête naturellement à ces formules frappées.
\end{definition}

\begin{propriete}[Structure type d'une affiche publicitaire chinoise]
La structure canonique à reproduire dans vos productions est :\\
\textbf{[Accroche / slogan]} + \textbf{[Description avec 不但…而且…]} + \textbf{[Argument avec 只要…就…]} + \textbf{[Incitation avec 与其…不如…]} + \textbf{[Impératif doux + QR code]}\\
Chaque phrase est courte (moins de 12 caractères si possible), sans sujet explicite quand le produit est évident, et ponctuée de ！ pour l'énergie.
\end{propriete}

\begin{popfigure}
\begin{tikzpicture}[mindmap, grow cyclic, every node/.style={concept, align=center, font=\small},
root concept/.append style={concept color=popblue, font=\bfseries},
level 1/.append style={level distance=3.2cm, sibling angle=90},
level 2/.append style={level distance=2.2cm, sibling angle=50, font=\footnotesize}]
\node[root concept, align=center]{Affiche\\publicitaire\\海报}
child[concept color=poporange]{node[align=center]{Slogan\\标语\\绿色有机}
child{node{纯天然}}
child{node{最纯正}}}
child[concept color=popgreen]{node[align=center]{Description\\描述\\口感鲜爽}
child{node{营养丰富}}
child{node{香气浓郁}}}
child[concept color=poppurple]{node[align=center]{Arguments\\论据\\限时优惠}
child{node{买一送一}}
child{node{免费品尝}}}
child[concept color=poppink]{node[align=center]{Appel à\\l'action\\扫码下单}
child{node{请尝试}}
child{node{立即购买}}};
\end{tikzpicture}
\legende{Carte mentale : les quatre zones de l'affiche et leurs mots-clés}
\end{popfigure}

\courssub{Le vocabulaire spécifique du marketing}

Repérons les expressions-clés de l'affiche. Ce sont des \textbf{formules figées du marketing chinois} que vous devez savoir réemployer.

\begin{center}
\begin{tabularx}{\linewidth}{|X|X|X|X|}
\hline
\rowcolor{popblueL} Caractères & Pinyin & Nature & Traduction \\
\hline
绿色有机 & lǜsè yǒujī & loc. adj. & vert et biologique \\
\hline
纯天然 & chún tiānrán & loc. adj. & 100 \% naturel \\
\hline
口感鲜爽 & kǒugǎn xiānshuǎng & loc. adj. & goût frais et vif (en bouche) \\
\hline
营养丰富 & yíngyǎng fēngfù & loc. adj. & très nutritif, riche en nutriments \\
\hline
限时优惠 & xiànshí yōuhuì & nom & offre promotionnelle limitée dans le temps \\
\hline
扫码下单 & sǎomǎ xiàdān & verbe & scanner le code (QR) et commander \\
\hline
立即尝试 & lìjí chángshì & verbe & essayer immédiatement \\
\hline
最纯正的选择 & zuì chúnzhèng de xuǎnzé & nom & le choix le plus authentique \\
\hline
\end{tabularx}
\end{center}

\begin{aretenir}[Mots-clés à réemployer]
Pour vos propres affiches (café de l'Ouest, poivre de Penja, miel de l'Adamaoua), gardez ces briques : \zh{纯天然} (100 \% naturel), \zh{限时优惠} (offre limitée), \zh{扫码下单} (scannez et commandez), \zh{请品尝} (goûtez, s'il vous plaît). Ce sont des blocs prêts à l'emploi.
\end{aretenir}

\coursec{Structures et connecteurs pour la persuasion}

L'affiche utilise trois connecteurs corrélatifs majeurs du niveau HSK 3-4. Observons-les avant d'énoncer la règle.

\courssub{不但…而且… « non seulement… mais aussi… »}

\zh{不但口感鲜爽，而且营养丰富。} \emph{Bùdàn kǒugǎn xiānshuǎng, érqiě yíngyǎng fēngfù.} « Non seulement le goût est frais, mais en plus il est nutritif. »

\begin{propriete}[Structure 不但…而且…]
\textbf{Structure :} \textbf{不但 + Proposition 1，而且 + Proposition 2}.\\
\textbf{Emploi :} La deuxième proposition ajoute un argument plus fort, souvent décisif.\\
\textbf{Sujet :} Si le sujet est identique, il se place \textbf{avant} \zh{不但} (ex: \zh{这茶不但…而且…}) ; s'il est différent, chaque proposition a son propre sujet.\\
\textbf{Niveau :} HSK 3.
\end{propriete}

\begin{exemplebox}[不但…而且… en contexte]
\zh{这种咖啡不但便宜，而且好喝。} \emph{Zhè zhǒng kāfēi bùdàn piányi, érqiě hǎohē.} « Ce café est non seulement bon marché, mais aussi délicieux. »\\[3pt]
\zh{喀麦隆的芒果不但甜，而且个大。} \emph{Kāmàilóng de mángguǒ bùdàn tián, érqiě gè dà.} « Les mangues du Cameroun sont non seulement sucrées, mais aussi grosses. »
\end{exemplebox}

\courssub{只要…就… « il suffit que… pour que… »}

\zh{只要扫码下单，即享限时优惠。} \emph{Zhǐyào sǎomǎ xiàdān, jí xiǎng xiànshí yōuhuì.} « Il suffit de scanner et commander pour profiter immédiatement de l'offre. »

\begin{propriete}[Structure 只要…就…]
\textbf{Structure :} \textbf{只要 + Condition，就 / 即 + Résultat}.\\
\textbf{Emploi :} Présente l'action (achat) comme une condition minime pour un bénéfice immédiat. \zh{即} (jí) est plus formel / écrit que \zh{juste}, fréquent dans la publicité.\\
\textbf{Distinction :} Ne pas confondre avec \zh{只有…才…} (« seulement si… alors »), qui exprime une condition \textbf{nécessaire} et restrictive. \zh{只要} exprime une condition \textbf{suffisante} et facile.
\end{propriete}

\begin{exemplebox}[只要…就… en contexte]
\zh{只要买一送一，顾客就会来。} \emph{Zhǐyào mǎi yī sòng yī, gùkè jiù huì lái.} « Dès qu'on offre un article pour un acheté, les clients viennent. »\\[3pt]
\zh{只要关注公众号，就送小礼品。} \emph{Zhǐyào guānzhù gōngzhònghào, jiù sòng xiǎo lǐpǐn.} « Il suffit de s'abonner au compte officiel pour recevoir un petit cadeau. »
\end{exemplebox}

\begin{attention}[Piège : 只要 vs 只有]
\begin{itemize}
\item \zh{只要你去，我就去。} « \textbf{Si seulement} tu y vas, j'irai. » (Condition suffisante, facile).
\item \zh{只有你去，我才去。} « \textbf{Seulement si} tu y vas, j'irai (sinon pas). » (Condition nécessaire, restrictive).
\end{itemize}
En publicité, on utilise presque toujours \zh{只要} pour baisser la barrière psychologique à l'achat.
\end{attention}

\courssub{与其…不如… « plutôt que… mieux vaut… »}

\zh{与其犹豫，不如立即尝试！} \emph{Yǔqí yóuyù, bùrú lìjí chángshì!} « Plutôt que d'hésiter, mieux vaut essayer tout de suite ! »

\begin{propriete}[Structure 与其…不如…]
\textbf{Structure :} \textbf{与其 + Option rejetée (verbe/adj.), 不如 + Option recommandée (verbe/adj.)}.\\
\textbf{Emploi :} Structure rhétorique forte. Elle balaye l'objection (hésitation, attente, comparaison) et impose l'action souhaitée. Très utilisée en slogan final.
\end{propriete}

\begin{exemplebox}[与其…不如… en contexte]
\zh{与其等待，不如现在就下单！} \emph{Yǔqí děngdài, bùrú xiànzài jiù xiàdān!} « Plutôt que d'attendre, commandez maintenant ! »\\[3pt]
\zh{与其买别的，不如选这个。} \emph{Yǔqí mǎi biéde, bùrú xuǎn zhège.} « Plutôt que d'acheter un autre, mieux vaut choisir celui-ci. »
\end{exemplebox}

\courssub{Le registre publicitaire : concision, impératif doux, superlatif}

Trois traits stylistiques caractérisent l'affiche chinoise :

\begin{itemize}
\item \textbf{Style concis (言简意赅)} : phrases sans sujet, sans verbe copule, souvent réduites à des adjectifs juxtaposés (\zh{绿色有机，纯天然！}). En français, nous dirions « C'est un produit vert, bio et naturel » : le chinois supprime tout ce qui n'est pas essentiel.
\item \textbf{Impératif doux (委婉祈使)} : \zh{请尝试} (« veuillez essayer »), \zh{请品尝} (« veuillez goûter »), \zh{欢迎选购} (« bienvenue pour choisir/acheter »). La particule \zh{请} adoucit l'ordre : le client est invité, pas commandé. Souvent, \zh{请} est omis mais l'intonation reste invitante (\zh{立即尝试}).
\item \textbf{Superlatif (最高级)} : \zh{最纯正的选择} (« le choix le plus authentique »). Le superlatif chinois se forme simplement avec \zh{最} + adjectif (+ \zh{的} + nom). Pas de complément de comparaison (\emph{de tous}) nécessaire.
\end{itemize}

\begin{attention}[Piège n°1 : pas de 是 devant l'adjectif prédicatif]
En français on dit « le goût \emph{est} frais ». En chinois, l'adjectif fonctionne directement comme prédicat, \textbf{sans} \zh{是} :\\
\zh{口感鲜爽。} \emph{Kǒugǎn xiānshuǎng.} « Le goût est frais. » (correct)\\
\zh{口感是鲜爽。} (incorrect en phrase descriptive neutre)\\
De même : \zh{这个茶很好。} \emph{Zhège chá hěn hǎo.} « Ce thé est très bon. » — \zh{很} joue ici le rôle de lien (souvent affaibli), jamais \zh{这个茶是很好}. Réservez \zh{是} aux phrases d'identification avec un nom : \zh{这是绿茶。} « C'est du thé vert. »
\end{attention}

\begin{attention}[Piège n°2 : le superlatif sans complément de comparaison]
Le français dit « le choix le plus authentique \emph{de tous} ». Le chinois n'a pas besoin de ce complément : \zh{最纯正的选择} suffit. N'ajoutez pas de structure lourde calquée du français (ex: \zh{在所有选择中最纯正的选择} — trop long pour une affiche).
\end{attention}

\begin{exoresolu}[Transformer une phrase française en slogan chinois]
Énoncé : transformez en deux slogans chinois adaptés à une affiche : « Notre miel est naturel et délicieux » et « Achetez maintenant, l'offre est limitée ! »
\tcblower
Solution détaillée :\\
1. On supprime le verbe « est » (adjectif prédicatif sans 是) et on condense en parallélisme : \zh{天然蜂蜜，香甜可口！} \emph{Tiānrán fēngmì, xiāngtián kěkǒu!} « Miel naturel, doux et savoureux ! » — deux blocs de 4 syllabes, rythme publicitaire typique.\\
2. On utilise l'impératif doux + nom : \zh{立即购买，优惠有限！} \emph{Lìjí gòumǎi, yōuhuì yǒuxiàn!} « Achetez immédiatement, l'offre est limitée ! »\\
\emph{Erreurs à éviter :} ne pas écrire \zh{蜂蜜是香甜} (\zh{是} interdit devant l'adjectif) ; ne pas traduire mot à mot « l'offre est limitée » par une phrase longue avec \zh{是} : le style affiche exige la concision nominalisée.
\end{exoresolu}

\courssub{Comparaison culturelle : affiche chinoise et affiche camerounaise}

\begin{center}
\begin{tabularx}{\linewidth}{|X|X|X|}
\hline
\rowcolor{popblueL} Critère & Affiche chinoise & Affiche camerounaise typique \\
\hline
Part du texte & Importante : slogans écrits, jeux de mots sur les caractères & Réduite : le visuel domine \\
\hline
Image & Souvent sobre, calligraphie valorisée & Très imagée : photo du produit, couleurs vives \\
\hline
Appel à l'action & QR code (\zh{扫码}) omniprésent & Numéro de téléphone, adresse du marché \\
\hline
Rythme & Parallélismes de 4 caractères (\zh{绿色有机}) & Slogans oraux, souvent en français ou en pidgin \\
\hline
\end{tabularx}
\end{center}

En Chine, l'écrit elle-même est un argument de vente : un beau caractère calligraphié inspire confiance. Au Cameroun, l'affiche s'adresse d'abord à l'œil et à l'oreille du passant pressé au carrefour. Pour votre production finale, retenez l'exigence chinoise : \textbf{peu de mots, mais des mots justes, rythmés et correctement écrits}.

\begin{savaistu}[La loi encadre les superlatifs]
En Chine, la loi sur la publicité (广告法) interdit dans les annonces commerciales les termes absolus comme \zh{国家级} (« de niveau national »), \zh{最高级} (« le plus haut niveau ») et \zh{最佳} (« le meilleur »), sous peine d'amendes importantes. Les marques contournent l'interdiction par des formules subjectives comme \zh{最纯正的选择} (« le choix le plus authentique »), qui qualifient le choix du consommateur et non le produit lui-même. Voilà pourquoi notre affiche utilise précisément cette tournure !
\end{savaistu}

\begin{exercice}[Repérage dans l'affiche]
Dans l'affiche du thé vert du Yunnan : 1) soulignez les quatre zones et nommez-les ; 2) entourez les trois connecteurs corrélatifs ; 3) relevez deux adjectifs prédicatifs employés sans \zh{是} ; 4) identifiez l'impératif doux et le superlatif.
\end{exercice}

\begin{corrige}[Exercice de repérage]
1) Zone 1 (accroche) : \zh{绿色有机，纯天然！} ; zone 2 (description) : \zh{不但口感鲜爽，而且营养丰富。} ; zone 3 (argument) : \zh{只要扫码下单，即享限时优惠。} ; zone 4 (appel à l'action) : \zh{与其犹豫，不如立即尝试！} + signature \zh{最纯正的选择}.\\
2) \zh{不但…而且…}, \zh{只要…就…} (ici \zh{即}), \zh{与其…不如…}.\\
3) \zh{口感鲜爽} et \zh{营养丰富} (adjectifs prédicatifs, sans \zh{是}).\\
4) Impératif doux : \zh{尝试} (sous-entendu \zh{请}) ; superlatif : \zh{最纯正的}.
\end{corrige}

\coursec{Écriture correcte des caractères du thème}

Maîtriser l'écriture des caractères clés du marketing est indispensable pour produire une affiche lisible et professionnelle. Nous travaillons l'ordre des traits (haut $\rightarrow$ bas, gauche $\rightarrow$ droite, horizontal avant vertical, cadre avant contenu), les radicaux (clés) et les risques de confusion.

\begin{definition}[Caractères clés de la leçon — Analyse graphique]
\begin{center}
\begin{tabularx}{\linewidth}{|c|X|X|X|X|}
\hline
\rowcolor{popblueL} Car. & Pinyin & Radical (Clé) & Nb traits & Points de vigilance / Caractères proches \\
\hline
\zh{有} & yǒu & \zh{月} (viande/chair) & 6 & Ne pas confondre avec \zh{右} (yòu, droite) : \zh{有} a \zh{月} à gauche, \zh{右} a \zh{口}. \\
\hline
\zh{机} & jī & \zh{木} (bois) & 6 & Partie droite \zh{幂} (gōng + petit trait). Ne pas écrire \zh{几} (jǐ, tablette/quelques) seul. \\
\hline
\zh{天} & tiān & \zh{大} (grand) & 4 & Un trait horizontal au-dessus de \zh{大}. Différent de \zh{夫} (fū, mari) qui a un trait en plus en bas. \\
\hline
\zh{然} & rán & \zh{火} (feu, forme \zh{灬} en bas) & 12 & Haut : \zh{肰} (shàn). Ne pas oublier les 4 points du feu en bas. \\
\hline
\zh{鲜} & xiān & \zh{鱼} (poisson) & 14 & Structure haut-bas : \zh{鱼} + \zh{羊}. \zh{羊} a un point en haut, pas un trait horizontal long. \\
\hline
\zh{爽} & shuǎng & \zh{月} (viande) & 11 & Haut : \zh{ㄟ} + \zh{大}. Bas : \zh{月}. Ne pas confondre avec \zh{亮} (liàng, clair) qui a \zh{京} en haut. \\
\hline
\zh{营} & yíng & \zh{艹} (herbe) & 11 & Simplifié de \zh{營}. Haut : \zh{艹} + \zh{呈}. Ne pas écrire \zh{营} avec \zh{火} en bas (forme traditionnelle). \\
\hline
\zh{养} & yǎng & \zh{月} (viande) & 8 & Simplifié de \zh{養}. Bas : \zh{羊} sans le point du haut $\rightarrow$ \zh{(trait brisé)} + \zh{月}. \\
\hline
\zh{扫} & sǎo & \zh{扌} (main) & 6 & Phonétique \zh{帚} (zhǒu, balai) à droite. \zh{帚} = \zh{巾} + \zh{由}. \\
\hline
\zh{码} & mǎ & \zh{石} (pierre) & 10 & Phonétique \zh{马} (mǎ, cheval). Ne pas confondre \zh{码} (code/unité) et \zh{马} (cheval). \\
\hline
\zh{下} & xià & \zh{一} (un) & 3 & Symétrique de \zh{上} (shàng). Trait central vertical traversant. \\
\hline
\zh{单} & dān & \zh{十} (dix) & 9 & Simplifié de \zh{單}. Ne pas confondre avec \zh{禅} (chán, méditation) qui a \zh{示} à gauche. \\
\hline
\zh{尝} & cháng & \zh{小} (petit) & 8 & Simplifié de \zh{嘗}. Haut : \zh{尚} (shàng) sans le trait horizontal du bas ? Non, \zh{尝} = \zh{小} + \zh{尚} (simplifié). Attention : \zh{尚} a 8 traits. \\
\hline
\zh{试} & shì & \zh{讠} (parole) & 8 & Simplifié de \zh{試}. Droite : \zh{式} (shì). \zh{式} = \zh{工} + \zh{弋}. \\
\hline
\zh{纯} & chún & \zh{纟} (soie) & 7 & Phonétique \zh{屯} (tún). Ne pas confondre avec \zh{准} (zhǔn, exact) qui a \zh{冰} à gauche. \\
\hline
\zh{正} & zhèng & \zh{止} (arrêt) & 5 & 5 traits : \zh{一} \zh{丨} \zh{一} \zh{丨} \zh{一}. Compter les traits pour ne pas en mettre 6. \\
\hline
\zh{选} & xuǎn & \zh{辶} (marche) & 9 & Simplifié de \zh{選}. Droite : \zh{巽} (xùn) simplifié en \zh{⺈} + \zh{彐} + \zh{ㄏ} ? Non, forme simplifiée : \zh{辶} + \zh{⺈} + \zh{人} + \zh{一} (ressemble à \zh{旦} sans trait). \\
\hline
\zh{择} & zé & \zh{扌} (main) & 10 & Phonétique \zh{咋} (zǎ) en haut à droite. \zh{咋} = \zh{口} + \zh{乍}. \\
\hline
\end{tabularx}
\end{center}
\end{definition}

\begin{methode}[Écrire un caractère complexe : la méthode des 3 étapes]
\begin{enumerate}
\item \textbf{Identifier le radical (clé) :} Il donne souvent l'indice sémantique (ex: \zh{扌} main pour \zh{扫}, \zh{石} pierre pour \zh{码}, \zh{鱼} poisson pour \zh{鲜}). L'écrire en premier (souvent à gauche ou en haut).
\item \textbf{Construire le phonétique / le corps :} Écrire la partie droite ou basse selon l'ordre standard (haut $\rightarrow$ bas, gauche $\rightarrow$ droite). Pour les caractères à structure \textbf{enclosure} (ex: \zh{选} avec \zh{辶}), écrire l'extérieur \textbf{avant} l'intérieur, et fermer le cadre en dernier.
\item \textbf{Vérifier le nombre de traits et l'alignement :} Compter les traits (ex: \zh{正} = 5 traits). Vérifier que les horizontaux sont parallèles et que le trait central vertical est bien centré.
\end{enumerate}
\end{methode}

\begin{experience}[Atelier calligraphie : « Mon slogan en caractères »]
\begin{enumerate}
\item Choisissez un produit camerounais (Poivre de Penja, Café de l'Ouest, Miel d'Adamaoua, Banane plantain).
\item Rédigez un slogan de 4 caractères (ex: \zh{纯天然香} — pur, naturel, parfumé).
\item Au tableau, un élève écrit le slogan au feutre-pinceau (ou craie) en respectant l'ordre des traits. La classe valide ou corrige (traits manquants, ordre inversé, radical mal formé).
\item Photographiez les meilleures productions pour le mur de la classe / le portfolio numérique.
\end{enumerate}
\end{experience}

\coursec{Exercices de thème (français → chinois)}

Traduisez les phrases suivantes en chinois (caractères + pinyin). Utilisez les structures \zh{不但…而且…}, \zh{只要…就…}, \zh{与其…不如…} et le vocabulaire de la leçon.

\begin{exercice}[Thème : Phrases isolées]
\begin{enumerate}
\item Ce café est non seulement bio, mais aussi très parfumé.
\item Il suffit de scanner le code pour obtenir un échantillon gratuit.
\item Plutôt que d'hésiter, goûtez-le tout de suite !
\item Notre poivre est 100 \% naturel, le choix le plus authentique.
\item Achetez maintenant, l'offre se termine demain.
\end{enumerate}
\end{exercice}

\begin{corrige}[Exercice de thème]
\begin{enumerate}
\item \zh{这种咖啡不但有机，而且香气浓郁。} \emph{Zhè zhǒng kāfēi bùdàn yǒujī, érqiě xiāngqì nóngyù.}
\item \zh{只要扫码，就送免费样品。} \emph{Zhǐyào sǎomǎ, jiù sòng miǎnfèi yàngpǐn.}
\item \zh{与其犹豫，不如立即品尝！} \emph{Yǔqí yóuyù, bùrú lìjí pǐncháng!}
\item \zh{我们的胡椒纯天然，最纯正的选择。} \emph{Wǒmen de hújiāo chún tiānrán, zuì chúnzhèng de xuǎnzé.}
\item \zh{立即购买，优惠明日结束。} \emph{Lìjí gòumǎi, yōuhuì míngrì jiéshù.}
\end{enumerate}
\end{corrige}

\begin{exercice}[Thème : Mini-texte publicitaire (Poivre de Penja)]
Rédigez une affiche complète (4 zones) pour le \emph{Poivre de Penja} en vous inspirant du modèle du thé vert. Contraintes : utiliser \zh{不但…而且…}, \zh{只要…就…}, \zh{与其…不如…}, un superlatif avec \zh{最}, un impératif doux. 40-60 caractères chinois.
\end{exercice}

\begin{corrige}[Exercice de thème — Proposition de correction]
\zh{黑色黄金，纯天然！}\\
\emph{Hēisè huángjīn, chún tiānrán!}\\
\zh{不但麻香浓郁，而且提味增鲜。}\\
\emph{Bùdàn máxiāng nóngyù, érqiě tíwèi zēngxiān.}\\
\zh{只要扫码下单，即享包邮优惠。}\\
\emph{Zhǐyào sǎomǎ xiàdān, jí xiǎng bāoyóu yōuhuì.}\\
\zh{与其寻找，不如拥有！}\\
\emph{Yǔqí xúnzhǎo, bùrú yǒngyǒu!}\\
\zh{—— 笔佳胡椒，厨房最佳拍档}\\
\emph{—— Bǐjiā hújiāo, chúfáng zuìjiā pāidàng.}
\emph{Note : « Meilleur partenaire » (\zh{最佳拍档}) est toléré comme superlatif subjectif (voir \emph{Le saviez-vous ?}).}
\end{corrige}

\coursec{Exercices de version (chinois → français)}

Traduisez en français les textes publicitaires suivants. Respectez le ton marketing (accrocheur, persuasif).

\begin{exercice}[Version : Phrases]
\begin{enumerate}
\item \zh{新鲜直达，营养锁住！}
\item \zh{不但价格实惠，而且品质保证。}
\item \zh{只要关注我们，就送会员价。}
\item \zh{与其羡慕别人，不如自己拥有。}
\item \zh{云南普洱，岁月陈香，最值得收藏的茶。}
\end{enumerate}
\end{exercice}

\begin{corrige}[Exercice de version]
\begin{enumerate}
\item « Fraîcheur directe, nutriments préservés ! » (Structure nominale parallèle, style slogan).
\item « Non seulement le prix est abordable, mais la qualité est garantie. »
\item « Il suffit de nous suivre pour bénéficier du prix membre. »
\item « Plutôt que d'envier les autres, mieux vaut le posséder soi-même. »
\item « Thé Pu'er du Yunnan, parfum vieilli par les années, le thé le plus digne d'être collectionné. »
\end{enumerate}
\end{corrige}

\begin{exercice}[Version : Texte complet — Affiche « Miel de l'Adamaoua »]
Traduisez l'affiche suivante en français, en restituant les 4 zones.

\zh{天然蜂蜜，甘甜润喉！}\\
\zh{不但增强免疫，而且美容养颜。}\\
\zh{只要下单满百，就享九折优惠。}\\
\zh{与其观望，不如立即囤货！}\\
\zh{—— 阿达马瓦蜂蜜，大自然最甜的礼物}
\end{exercice}

\begin{corrige}[Exercice de version — Texte complet]
\textbf{Accroche :} Miel naturel, doux et apaisant pour la gorge !\\
\textbf{Description :} Non seulement il renforce l'immunité, mais en plus il embellit le teint et nourrit la peau.\\
\textbf{Argument :} Il suffit de commander pour 100 yuans (ou « à partir de 100 ») pour profiter de 10 \% de réduction.\\
\textbf{Appel à l'action :} Plutôt que d'attendre, faites des réserves tout de suite !\\
\textbf{Signature :} — Miel de l'Adamaoua, le cadeau le plus doux de la nature.
\end{corrige}

\coursec{Grille d'évaluation et production finale}

\begin{aretenir}[Grille d'évaluation de l'affiche publicitaire (20 points)]
\begin{center}
\begin{tabularx}{\linewidth}{|X|c|X|}
\hline
\rowcolor{popblueL} Critères & Points & Indicateurs de réussite \\
\hline
\textbf{Structure (4 zones)} & 4 & Les 4 zones présentes, dans l'ordre, délimitées par \\\\ ou sauts de ligne. \\
\hline
\textbf{Grammaire (Connecteurs)} & 5 & Utilisation correcte d'au moins 2 structures parmi : \zh{不但…而且…}, \zh{只要…就…}, \zh{与其…不如…}. Pas de \zh{是} + adj. \\
\hline
\textbf{Vocabulaire marketing} & 4 & Emploi pertinent de 4+ termes du tableau (ex: \zh{纯天然}, \zh{限时优惠}, \zh{扫码下单}, \zh{最...选择}). \\
\hline
\textbf{Registre \& Style} & 3 & Phrases concises, rythme (parallélisme 4+4 ou 3+3), impératif doux (\zh{请} ou implicite), superlatif \zh{最}. \\
\hline
\textbf{Écriture des caractères} & 2 & Caractères lisibles, traits dans l'ordre, radicaux corrects, pas de confusion visuelle (ex: \zh{有}/\zh{右}, \zh{码}/\zh{马}). \\
\hline
\textbf{Adaptation culturelle} & 2 & Produit camerounais mis en valeur, argumentaire adapté au public chinois (QR code, naturel, santé). \\
\hline
\textbf{Total} & \textbf{20} & \\
\hline
\end{tabularx}
\end{center}
\end{aretenir}

\begin{experience}[Production finale : « Mon affiche pour le Cameroun »]
\textbf{Consigne :} Vous travaillez pour l'agence \emph{« Cameroun Export »}. Concevez une affiche publicitaire \textbf{en chinois} (caractères + pinyin) pour l'un de ces produits : \emph{Café Arabica de l'Ouest}, \emph{Poivre blanc de Penja}, \emph{Miel blanc de l'Adamaoua}, \emph{Banane plantain chips}.

\textbf{Étapes :}
\begin{enumerate}
\item \textbf{Brouillon (10 min) :} Remplissez le tableau des 4 zones en français, puis transposez en chinois en utilisant les structures imposées.
\item \textbf{Rédaction soignée (15 min) :} Écrivez la version finale sur une feuille A5 (format affiche). Soignez l'écriture (taille, alignement). Ajoutez le pinyin en petit sous chaque ligne.
\item \textbf{Présentation orale (5 min / groupe) :} Présentez votre affiche à la classe en chinois : lisez le texte, expliquez le choix du slogan et de l'argumentaire. La classe note sur la grille d'évaluation (pair-évaluation).
\end{enumerate}

\textbf{Modèle de fiche de brouillon :}
\begin{center}
\begin{tabularx}{\linewidth}{|X|X|}
\hline
\rowcolor{popblueL} Zone & Contenu chinois (Caractères + Pinyin) \\
\hline
1. Slogan (4-8 car.) & \\
\hline
2. Description (\zh{不但…而且…}) & \\
\hline
3. Argument (\zh{只要…就…}) & \\
\hline
4. Appel à l'action (\zh{与其…不如…} + Impératif) & \\
\hline
Signature (\zh{最...选择}) & \\
\hline
\end{tabularx}
\end{center}
\end{experience}

\begin{methode}[Auto-correction avant remise : Check-list 5 points]
\begin{enumerate}
\item \textbf{Connecteurs :} Ai-je utilisé \zh{不但…而且…} / \zh{只要…就…} / \zh{与其…不如…} correctement (pas de \zh{是} après, ordre des propositions) ?
\item \textbf{Adjectifs prédicatifs :} Mes phrases du type « Produit + Adj » n'ont-elles pas de \zh{是} ? (Ex: \zh{茶香} et non \zh{茶是香}).
\item \textbf{Superlatif :} Mon \zh{最...} est-il suivi de \zh{的} + Nom ? N'ai-je pas ajouté « de tous » ?
\item \textbf{Caractères pièges :} Ai-je vérifié \zh{有/右}, \zh{码/马}, \zh{鲜/鱼+羊}, \zh{选/辶+...}, \zh{正/5 traits} ?
\item \textbf{Ponctuation :} Ai-je utilisé les signes chinois pleins \zh{，} \zh{。} \zh{！} \zh{：} ? Pas d'espaces entre caractères.
\end{enumerate}
\end{methode}

\coursec{Structures et connecteurs pour la persuasion}

Dans la section précédente, nous avons analysé un modèle d'affiche publicitaire : nous avons repéré son titre accrocheur, son slogan, sa description de produit et son appel à l'action. Nous allons maintenant entrer dans l'atelier du rédacteur publicitaire : quelles \cle{structures grammaticales} et quels \cle{connecteurs} utiliser pour convaincre un client ? C'est toute la grammaire de la persuasion que nous allons construire ici, pas à pas.

\courssub{Révision des patterns de persuasion}

\textbf{Le pattern progressif : 不但 A 而且 B}\par

Observons d'abord deux phrases tirées d'affiches :

\zh{这款手机不但便宜，而且质量好。} \emph{Zhè kuǎn shǒujī bùdàn piányi, érqiě zhìliàng hǎo.} « Ce téléphone est non seulement bon marché, mais aussi de bonne qualité. »

\zh{本店不但卖咖啡，而且卖喀麦隆的可可。} \emph{Běn diàn bùdàn mài kāfēi, érqiě mài Kāmàilóng de kěkě.} « Notre magasin vend non seulement du café, mais aussi du cacao du Cameroun. »

\begin{propriete}[Structure : 不但 A 而且 B]
\cle{Forme} : Sujet + 不但 + A，而且 + B。\\
\cle{Emploi} : exprime une progression (« non seulement... mais aussi... »). L'argument B vient renforcer l'argument A. En publicité, on place souvent l'argument le plus fort en deuxième position.\\
\cle{Exception} : si les deux parties ont des sujets différents, 不但 se place avant le premier sujet : \zh{不但价格低，而且服务也好。} \emph{Bùdàn jiàgé dī, érqiě fúwù yě hǎo.} « Non seulement les prix sont bas, mais le service est bon aussi. »
\end{propriete}

\begin{attention}[Erreur fréquente des francophones]
Le français dit « non seulement... mais aussi... » avec deux mots séparés ; en chinois, 不但 et 而且 forment un couple inséparable : on ne peut pas utiliser 不但 seul sans 而且 (ou 也/还) dans la seconde partie. Phrase fautive : *\zh{这个商品不但很好。} — il manque la suite !
\end{attention}

\textbf{Le pattern conditionnel attractif : 只要 A 就 B}\par

Observons :

\zh{只要关注公众号，就送小礼品。} \emph{Zhǐyào guānzhù gōngzhònghào, jiù sòng xiǎo lǐpǐn.} « Il suffit de suivre notre compte officiel pour recevoir un petit cadeau. »

\zh{只要今天来本店，就能享受八折优惠。} \emph{Zhǐyào jīntiān lái běn diàn, jiù néng xiǎngshòu bā zhé yōuhuì.} « Venez aujourd'hui dans notre magasin et profitez de 20 \% de réduction. »

\begin{propriete}[只要 + condition + 就 + résultat positif]
\cle{Forme} : 只要 + condition，就 + résultat。\\
\cle{Emploi} : « du moment que..., il suffit que... pour que... ». La condition est présentée comme facile à remplir, et le résultat comme automatique. C'est LA structure préférée du marketing chinois : elle minimise l'effort du client et maximise le bénéfice.\\
\cle{Point de comparaison} : le français utilise « il suffit de... pour... » ; le chinois garde l'ordre condition d'abord, résultat ensuite — jamais l'inverse.\\
\cle{Culture du chiffre} : \zh{八折} \emph{bā zhé} signifie littéralement « huit dixièmes » : le client paie 80 \% du prix, soit 20 \% de réduction. \zh{半价} \emph{bànjià} = moitié prix (–50 \%).
\end{propriete}

\textbf{Le pattern de choix préférentiel : 与其 A 不如 B}\par

Observons :

\zh{与其等待，不如行动。} \emph{Yǔqí děngdài, bùrú xíngdòng.} « Plutôt qu'attendre, agis. »

\zh{与其买贵的，不如买对的。} \emph{Yǔqí mǎi guì de, bùrú mǎi duì de.} « Plutôt qu'acheter cher, achète juste. »

\begin{propriete}[Structure : 与其 A 不如 B]
\cle{Forme} : 与其 + option rejetée，不如 + option préférée。\\
\cle{Emploi} : « plutôt que A, mieux vaut B ». Le rédacteur écarte une mauvaise habitude (attendre, payer cher, hésiter) pour orienter le client vers l'action d'achat. Structure très efficace pour les slogans courts.\\
\cle{Attention à l'ordre} : en français « mieux vaut B que A » commence par la bonne option ; en chinois, l'option rejetée (与其) vient TOUJOURS en premier.
\end{propriete}

\begin{center}
\begin{tabularx}{\linewidth}{|X|X|X|}
\hline
\rowcolor{popblueL} Structure & Exemple (caractères + pinyin) & Traduction \\
\hline
不但 A 而且 B & \zh{不但漂亮，而且实用。} \emph{Bùdàn piàoliang, érqiě shíyòng.} & Non seulement beau, mais aussi pratique. \\
\hline
只要 A 就 B & \zh{只要扫码，就有惊喜。} \emph{Zhǐyào sǎomǎ, jiù yǒu jīngxǐ.} & Scannez le code et une surprise vous attend. \\
\hline
与其 A 不如 B & \zh{与其犹豫，不如试试。} \emph{Yǔqí yóuyù, bùrú shìshi.} & Plutôt qu'hésiter, essayez. \\
\hline
\end{tabularx}
\end{center}

\courssub{Les connecteurs : organiser le message publicitaire}

Une affiche ou un texte promotionnel bien construit guide le lecteur. Trois familles de connecteurs sont indispensables.

\begin{definition}[Connecteurs d'énumération]
\zh{首先} \emph{shǒuxiān} « d'abord, en premier » ; \zh{其次} \emph{qícì} « ensuite, en second lieu » ; \zh{最后} \emph{zuìhòu} « enfin ». Ils ordonnent les arguments du produit, du plus important au dernier.
\end{definition}

\begin{definition}[Connecteurs d'addition et de conclusion]
Addition : \zh{此外} \emph{cǐwài} « de plus » ; \zh{另外} \emph{lìngwài} « en outre, par ailleurs » ; \zh{而且} \emph{érqiě} « et aussi ».\\
Conclusion : \zh{总之} \emph{zǒngzhī} « en résumé, bref » ; \zh{一句话} \emph{yí jù huà} « en un mot » (très utilisé pour clore un argumentaire : \zh{一句话，买它！} \emph{Yí jù huà, mǎi tā!} « En un mot : achetez-le ! »).
\end{definition}

\begin{exemplebox}[Un paragraphe argumentatif complet]
\zh{首先品质优，其次价格低，最后服务好。此外，本周还有八折优惠。总之，现在买最划算！}\\
\emph{Shǒuxiān pǐnzhì yōu, qícì jiàgé dī, zuìhòu fúwù hǎo. Cǐwài, běn zhōu hái yǒu bā zhé yōuhuì. Zǒngzhī, xiànzài mǎi zuì huásuàn!}\\
« D'abord, la qualité est excellente ; ensuite, les prix sont bas ; enfin, le service est bon. De plus, cette semaine, il y a encore 20 \% de réduction. Bref, c'est maintenant que c'est le plus avantageux ! »
\end{exemplebox}

\courssub{Les figures de style publicitaires}

\begin{definition}[Trois figures de style du marketing]
\cle{排比} \emph{páibǐ} : le parallélisme — répéter la même structure dans trois phrases ou plus pour créer un rythme.\\
\cle{对比} \emph{duìbǐ} : le contraste — opposer avant/après, cher/bon marché, attendre/agir.\\
\cle{夸张} \emph{kuāzhāng} : l'hyperbole — exagérer avec modération pour frapper l'esprit.
\end{definition}

\begin{exemplebox}[Figures de style en action]
Parallélisme (排比) : \zh{买得放心，用得舒心，吃得开心。} \emph{Mǎi de fàngxīn, yòng de shūxīn, chī de kāixīn.} « Achetez en confiance, utilisez en confort, mangez en joie. »\\
Contraste (对比) : \zh{昨天很贵，今天半价。} \emph{Zuótiān hěn guì, jīntiān bànjià.} « Hier c'était cher, aujourd'hui c'est à moitié prix. »\\
Hyperbole modérée (夸张) : \zh{全雅温得最好喝的咖啡！} \emph{Quán Yǎwēndé zuì hǎohē de kāfēi!} « Le meilleur café de tout Yaoundé ! »
\end{exemplebox}

\begin{attention}[Hyperbole : avec modération]
En chinois comme en français, l'hyperbole publicitaire doit rester crédible. On préfère \zh{很好} « très bon » ou \zh{最好的之一} « l'un des meilleurs » à des promesses impossibles. Une hyperbole trop forte rend l'affiche ridicule.
\end{attention}

\courssub{Les particules modales de l'incitation}

\begin{propriete}[吧 et 呢 dans le discours commercial]
\cle{吧} \emph{ba} : adoucit l'impératif et le transforme en suggestion. \zh{快来吧！} \emph{Kuài lái ba!} « Venez vite, allez ! »\\
\cle{请 + verbe + 吧} : politesse commerciale. \zh{请进店看看吧！} \emph{Qǐng jìn diàn kànkan ba!} « Entrez jeter un œil, je vous en prie ! »\\
\cle{呢} \emph{ne} : crée une attente, une incitation interrogative. \zh{这么好的价格，你还在等什么呢？} \emph{Zhème hǎo de jiàgé, nǐ hái zài děng shénme ne?} « Un si bon prix, qu'attendez-vous encore ? »
\end{propriete}

\begin{attention}[吧 n'est pas un point d'exclamation]
Le francophone traduit souvent « ! » par un ton brusque. En chinois, l'impératif nu (\zh{买！} « Achète ! ») sonne agressif ; la particule 吧 le rend amical. Sur une affiche, on écrit presque toujours \zh{快来看看吧！} et jamais \zh{看！} seul.
\end{attention}

\courssub{La place des compléments de temps et de lieu}

\begin{propriete}[Temps et lieu AVANT le verbe]
En chinois, les compléments de temps et de lieu se placent \textbf{avant le verbe}, alors qu'en français ils se placent souvent après.\\
\cle{Schéma} : Sujet + temps + lieu + verbe + complément.\\
\zh{我们今天在本店举行促销活动。} \emph{Wǒmen jīntiān zài běn diàn jǔxíng cùxiāo huódòng.} « Nous organisons une promotion dans notre magasin aujourd'hui. »\\
Ordre français : verbe → lieu → temps. Ordre chinois : temps → lieu → verbe. C'est l'inverse !
\end{propriete}

\begin{exemplebox}[Compare les deux ordres]
\zh{明天在杜阿拉市场见！} \emph{Míngtiān zài Dù'ālā shìchǎng jiàn!} « Rendez-vous demain au marché de Douala ! » (chinois : demain + au marché + voir ; français : voir + demain + au marché).
\end{exemplebox}

\courssub{Piège lexical : 优惠 ou 优质 ?}

\begin{attention}[优惠, 优质, 优秀 : trois mots à ne pas confondre]
\zh{优惠} \emph{yōuhuì} « avantageux, réduction » (prix) : \zh{优惠价} « prix promotionnel ».\\
\zh{优质} \emph{yōuzhì} « de qualité supérieure » (produit) : \zh{优质产品} « produit de qualité supérieure ».\\
\zh{优秀} \emph{yōuxiù} « excellent » (personne, élève) : \zh{优秀学生} « excellent élève ».\\
优惠 et 优秀 partagent le même premier caractère 优 (radical 人, « personne ») mais leur sens est radicalement différent : on ne dit jamais *\zh{优秀价} pour « prix réduit » !
\end{attention}

\begin{center}
\begin{tabularx}{\linewidth}{|X|X|X|X|}
\hline
\rowcolor{popblueL} Caractères & Pinyin & Nature & Traduction \\
\hline
\zh{优惠} & yōuhuì & adj./n. & avantageux ; réduction \\
\hline
\zh{优质} & yōuzhì & adj. & de qualité supérieure \\
\hline
\zh{优秀} & yōuxiù & adj. & excellent (personne) \\
\hline
\zh{打折} & dǎzhé & v. & faire une réduction \\
\hline
\zh{划算} & huásuàn & adj. & avantageux, rentable \\
\hline
\zh{促销} & cùxiāo & v./n. & promouvoir ; promotion \\
\hline
\zh{扫码} & sǎomǎ & v. & scanner un code QR \\
\hline
\zh{礼品} & lǐpǐn & n. & cadeau \\
\hline
\end{tabularx}
\end{center}

\courssub{Le flux logique de la persuasion}

Une phrase publicitaire chinoise suit souvent un enchaînement en trois temps : on pose une condition facile, on promet un avantage, on lance un appel à l'action.

\begin{popfigure}
\begin{tikzpicture}[node distance=1.6cm, >=Stealth]
\node[draw=popblue, fill=popblueL, rounded corners, align=center, minimum width=3.4cm, minimum height=1.1cm] (cond) {\zh{只要}\,+\,condition\\ \emph{zhǐyào} : « il suffit de »};
\node[draw=popgreen, fill=popgreenL, rounded corners, align=center, minimum width=3.4cm, minimum height=1.1cm, right=of cond] (av) {\zh{就}\,+\,avantage\\ \emph{jiù} : « alors, aussitôt »};
\node[draw=poporange, fill=poporangeL, rounded corners, align=center, minimum width=3.4cm, minimum height=1.1cm, right=of av] (appel) {\zh{请}\,+\,appel à l'action\\ \emph{qǐng} : « veuillez »};
\draw[->, very thick, popblue] (cond) -- node[above, align=center]{condition facile} (av);
\draw[->, very thick, popgreen] (av) -- node[above, align=center]{bénéfice promis} (appel);
\end{tikzpicture}
\legende{Flux logique : Condition (\zh{只要}) → Avantage (\zh{就}) → Appel (\zh{请扫码}). Exemple : \zh{只要关注公众号，就送小礼品，请扫码吧！}}
\end{popfigure}

\courssub{Texte d'application : une annonce promotionnelle}

\begin{textebox}[Annonce du magasin « 好朋友 » à Yaoundé]
\zh{好消息！好朋友商店本周举行大促销。我们的商品不但质量好，而且价格低。首先，喀麦隆咖啡八折优惠；其次，优质茶叶半价；最后，买两件衣服送一件。此外，只要关注我们的公众号，就送小礼品。与其等待，不如行动！这么好的价格，你还在等什么呢？一句话：今天在本店，买得放心，用得舒心。请快来看看吧！}\\
\emph{Hǎo xiāoxi! Hǎo Péngyou shāngdiàn běn zhōu jǔxíng dà cùxiāo. Wǒmen de shāngpǐn bùdàn zhìliàng hǎo, érqiě jiàgé dī. Shǒuxiān, Kāmàilóng kāfēi bā zhé yōuhuì; qícì, yōuzhì cháyè bànjià; zuìhòu, mǎi liǎng jiàn yīfu sòng yí jiàn. Cǐwài, zhǐyào guānzhù wǒmen de gōngzhònghào, jiù sòng xiǎo lǐpǐn. Yǔqí děngdài, bùrú xíngdòng! Zhème hǎo de jiàgé, nǐ hái zài děng shénme ne? Yí jù huà: jīntiān zài běn diàn, mǎi de fàngxīn, yòng de shūxīn. Qǐng kuài kuài lái kànkan ba!}\\
« Bonne nouvelle ! Le magasin Bon Ami organise cette semaine une grande promotion. Nos produits sont non seulement de bonne qualité, mais aussi à bas prix. D'abord, le café du Cameroun à moins 20 \% ; ensuite, le thé de qualité supérieure à moitié prix ; enfin, deux vêtements achetés, un offert. De plus, il suffit de suivre notre compte officiel pour recevoir un petit cadeau. Plutôt qu'attendre, agissez ! Un si bon prix, qu'attendez-vous ? En un mot : aujourd'hui dans notre magasin, achetez en confiance, utilisez en confort. Venez vite voir ! »
\end{textebox}

\begin{methode}[Écrire un slogan publicitaire en 4 étapes]
\begin{enumerate}
\item \cle{Choisir} 2 ou 3 mots-clés du produit (ex. : \zh{咖啡} café, \zh{香} parfumé, \zh{便宜} bon marché).
\item \cle{Former} une phrase courte de 4 à 7 caractères (ex. : \zh{好咖啡，好价格}).
\item \cle{Vérifier} la rime ou l'allitération : les sons proches rendent le slogan mémorable (\zh{放心 / 舒心} riment en -xīn).
\item \cle{Tester} à l'oral : lire le slogan à voix haute ; s'il se retient facilement, il est bon.
\end{enumerate}
\end{methode}

\begin{exoresolu}[Transformer une phrase neutre en phrase publicitaire]
Énoncé : transforme la phrase neutre « Notre thé est bon et pas cher » en une phrase publicitaire utilisant 不但...而且... et une particule modale.\\
\tcblower
Solution détaillée :\\
Étape 1 : phrase neutre — \zh{我们的茶很好，也不贵。} \emph{Wǒmen de chá hěn hǎo, yě bù guì.}\\
Étape 2 : on insère le couple 不但...而且... pour la progression : \zh{我们的茶不但好喝，而且不贵。}\\
Étape 3 : on ajoute l'appel avec 吧 : \zh{我们的茶不但好喝，而且不贵，快来买吧！} \emph{Wǒmen de chá bùdàn hǎohē, érqiě bù guì, kuài lái mǎi ba!} « Notre thé est non seulement délicieux, mais aussi pas cher : venez vite l'acheter ! »\\
La phrase est passée d'un constat à une incitation : c'est l'effet du pattern progressif + particule 吧.
\end{exoresolu}

\begin{exercice}[À toi de jouer]
\begin{enumerate}
\item Complète avec 不但 / 而且 : \zh{这个商店\_\_\_\_卖咖啡，\_\_\_\_卖茶叶。}
\item Traduis en chinois avec 只要...就... : « Il suffit de venir aujourd'hui pour avoir un cadeau. »
\item Traduis en français : \zh{与其买贵的，不如买好的。}
\item Remets dans l'ordre chinois correct : 我们 / 明天 / 在市场 / 举行 / 促销活动。
\item Choisis le bon mot (优惠 / 优质 / 优秀) : \zh{\_\_\_\_产品} (produit de qualité supérieure) ; \zh{\_\_\_\_价格} (prix avantageux).
\end{enumerate}
\end{exercice}

\begin{corrige}[Exercice « À toi de jouer »]
\begin{enumerate}
\item \zh{这个商店不但卖咖啡，而且卖茶叶。} « Ce magasin vend non seulement du café, mais aussi du thé. »
\item \zh{只要今天来，就有礼品。} \emph{Zhǐyào jīntiān lái, jiù yǒu lǐpǐn.} (condition d'abord, résultat ensuite).
\item « Plutôt qu'acheter cher, mieux vaut acheter bon. »
\item \zh{我们明天在市场举行促销活动。} (sujet + temps + lieu + verbe).
\item \zh{优质产品} ; \zh{优惠价格}。
\end{enumerate}
\end{corrige}

\begin{savaistu}[Les couplets publicitaires (对联)]
En Chine, les grands magasins affichent souvent des \cle{couplets} (\zh{对联} \emph{duìlián}) en période de fête : deux phrases parallèles de même longueur, placées de chaque côté de la porte. Un classique du commerce : \zh{买得放心，用得舒心} \emph{Mǎi de fàngxīn, yòng de shūxīn} « Achetez en confiance, utilisez en confort ». Remarque le parallélisme parfait : même structure \zh{verbe + 得 + adjectif} dans les deux moitiés, et la rime en -xīn. Au Cameroun, les commerçants utilisent plutôt des slogans peints sur les murs des boutiques — deux cultures, un même art de convaincre !
\end{savaistu}

\begin{aretenir}[L'essentiel de la section]
\begin{itemize}
\item Trois patterns de persuasion : \zh{不但...而且...} (progression), \zh{只要...就...} (condition facile → avantage), \zh{与其...不如...} (choix préférentiel, option rejetée d'abord).
\item Connecteurs : \zh{首先、其次、最后} (énumération), \zh{此外、另外} (addition), \zh{总之、一句话} (conclusion).
\item Figures de style : \zh{排比} (parallélisme), \zh{对比} (contraste), \zh{夸张} (hyperbole modérée).
\item Particules : \zh{吧} adoucit et suggère, \zh{呢} incite, \zh{请...吧} est la politesse commerciale.
\item Temps et lieu se placent AVANT le verbe en chinois.
\item Ne pas confondre \zh{优惠} (prix), \zh{优质} (produit), \zh{优秀} (personne).
\end{itemize}
\end{aretenir}

\coursec{Écriture correcte des caractères du thème}

Dans la section précédente, nous avons appris à construire des phrases de persuasion avec des connecteurs comme \zh{因为} (yīnwèi, « parce que ») et \zh{所以} (suǒyǐ, « donc »). Mais une affiche publicitaire, aussi bien argumentée soit-elle, perd toute crédibilité si les caractères sont mal tracés : un trait manquant peut transformer \zh{产品} (chǎnpǐn, « produit ») en un caractère inexistant, et un client chinois le remarquera immédiatement. Cette section est donc consacrée à l'\cle{écriture correcte} des douze caractères clés de notre thème du marketing.

\courssub{Les douze caractères cibles : observation générale}

Avant d'écrire, observons. Voici les douze caractères que vous devez savoir écrire de mémoire à la fin de cette leçon, avec leur radical (la \cle{clé} qui sert à classer le caractère dans un dictionnaire) et leur nombre de traits.

\begin{center}
\begin{tabularx}{\linewidth}{|X|X|X|X|X|}\hline
\rowcolor{popblueL} Caractère & Pinyin & Radical & Traits & Traduction \\ \hline
广 & guǎng & 广 (toit) & 3 & large, vaste \\ \hline
告 & gào & 口 (bouche) & 7 & annoncer, informer \\ \hline
产 & chǎn & 亠 (tête) & 6 & produire \\ \hline
品 & pǐn & 口 (bouche) & 9 & article, produit, qualité \\ \hline
质 & zhì & 贝 (coquillage) & 8 & qualité, nature \\ \hline
优 & yōu & 亻 (homme) & 6 & excellent, supérieur \\ \hline
惠 & huì & 心 (cœur) & 12 & faveur, avantage \\ \hline
买 & mǎi & 乙 / 乛 & 6 & acheter \\ \hline
卖 & mài & 十 (dix) & 8 & vendre \\ \hline
赠 & zèng & 贝 (coquillage) & 16 & offrir, donner en cadeau \\ \hline
限 & xiàn & 阝 (colline) & 8 & limiter \\ \hline
量 & liàng & 里 (village, dedans) & 12 & quantité, mesure \\ \hline
\end{tabularx}
\end{center}

Observez attentivement : trois de ces caractères (\zh{质}, \zh{赠}, et aussi \zh{货} huò, « marchandise ») contiennent le radical \zh{贝} (bèi, « coquillage »). Ce n'est pas un hasard : dans la Chine ancienne, les coquillages servaient de monnaie. Tout caractère contenant \zh{贝} a donc souvent un rapport avec l'\cle{argent}, le commerce ou la valeur. C'est une aide précieuse pour la mémorisation : quand vous voyez \zh{贝} dans un caractère, pensez « argent ».

\begin{definition}[Le radical 广 (yǎn) : « toit, abri »]
Le radical \zh{广}, appelé \emph{yǎn} en chinois traditionnel de classement, représente un \cle{toit} ou un \cle{abri} appuyé contre une paroi. Il indique souvent un \textbf{lieu} ou une \textbf{action qui se déroule sous un toit} : \zh{店} (diàn, « magasin »), \zh{库} (kù, « entrepôt »), \zh{座} (zuò, « siège, place »). Dans \zh{广告} (guǎnggào, « publicité »), on peut mémoriser : « une annonce faite \emph{largement} (广), sous tous les toits ».
\end{definition}

\courssub{L'ordre des traits : la règle d'or de l'écriture chinoise}

En français, nous écrivons les lettres de gauche à droite, et l'ordre des traits d'un « a » importe peu. En chinois, l'\cle{ordre des traits} est rigoureusement codifié : il garantit des caractères équilibrés, rapides à tracer et lisibles. Les règles générales sont :

\begin{enumerate}
\item De \textbf{haut en bas} : on écrit d'abord la partie supérieure.
\item De \textbf{gauche à droite} : on écrit d'abord la partie gauche.
\item Le trait \textbf{horizontal avant le vertical} quand ils se croisent (ex. \zh{十}).
\item L'\textbf{extérieur avant l'intérieur}, puis on « ferme la porte » (ex. \zh{国}).
\item Le \textbf{point final} (丶) se trace souvent en dernier, comme la cerise sur le gâteau.
\end{enumerate}

\begin{methode}[Mémoriser l'ordre des traits : deux formules magiques]
Pour ne jamais hésiter, récitez à voix haute en écrivant :
\begin{enumerate}
\item Pour \zh{广} : « \textbf{Point, horizontal, coude} » (丶 → 一 → 丿). Le point se pose d'abord, comme une goutte sur le toit ; puis le toit horizontal ; puis la paroi qui descend en courbe vers la gauche.
\item Pour \zh{产} : « \textbf{Horizontal, vertical, horizontal, vertical, horizontal, point} » : on alterne, comme les marches d'un escalier, et le long trait horizontal du bas vient soutenir l'ensemble avant le point final.
\end{enumerate}
Entraînez-vous en « écrivant dans l'air » avec le doigt, en récitant la formule : la mémoire musculaire fait le reste.
\end{methode}

\begin{popfigure}
\begin{center}
\begin{tikzpicture}[scale=1.1]
% --- 广 : 3 traits ---
\draw[popline, thin] (0,0) grid (2,2);
\draw[popline, dashed, thin] (1,0) -- (1,2);
\draw[popline, dashed, thin] (0,1) -- (2,1);
\draw[popink, very thick] (1.05,1.75) -- (1.2,1.6);
\draw[popink, very thick] (0.35,1.35) -- (1.75,1.35);
\draw[popink, very thick] (0.35,1.35) .. controls (0.3,0.9) and (0.35,0.5) .. (0.15,0.15);
\node[circle, fill=poporange, text=white, inner sep=1.2pt, font=\scriptsize] at (0.75,1.85) {1};
\node[circle, fill=poporange, text=white, inner sep=1.2pt, font=\scriptsize] at (1.05,1.15) {2};
\node[circle, fill=poporange, text=white, inner sep=1.2pt, font=\scriptsize] at (0.55,0.6) {3};
\node[below, font=\small] at (1,-0.1) {\zh{广} guǎng (3 traits)};
% --- 产 : 6 traits ---
\begin{scope}[xshift=3.2cm]
\draw[popline, thin] (0,0) grid (2,2);
\draw[popline, dashed, thin] (1,0) -- (1,2);
\draw[popline, dashed, thin] (0,1) -- (2,1);
\draw[popink, very thick] (0.55,1.4) -- (1.5,1.4);
\draw[popink, very thick] (1.0,1.4) -- (1.0,1.85);
\draw[popink, very thick] (0.75,1.15) -- (0.62,1.0);
\draw[popink, very thick] (1.25,1.15) -- (1.38,1.0);
\draw[popink, very thick] (0.3,0.75) -- (1.7,0.75);
\draw[popink, very thick] (1.0,0.75) .. controls (0.85,0.5) and (0.6,0.3) .. (0.35,0.15);
\node[circle, fill=poporange, text=white, inner sep=1.2pt, font=\scriptsize] at (0.35,1.55) {1};
\node[circle, fill=poporange, text=white, inner sep=1.2pt, font=\scriptsize] at (1.2,1.7) {2};
\node[circle, fill=poporange, text=white, inner sep=1.2pt, font=\scriptsize] at (0.42,1.05) {3};
\node[circle, fill=poporange, text=white, inner sep=1.2pt, font=\scriptsize] at (1.58,1.05) {4};
\node[circle, fill=poporange, text=white, inner sep=1.2pt, font=\scriptsize] at (1.0,0.55) {5};
\node[circle, fill=poporange, text=white, inner sep=1.2pt, font=\scriptsize] at (0.7,0.28) {6};
\node[below, font=\small] at (1,-0.1) {\zh{产} chǎn (6 traits)};
\end{scope}
% --- 赠 : schéma en deux blocs ---
\begin{scope}[xshift=6.4cm]
\draw[popline, thin] (0,0) grid (2,2);
\draw[popline, dashed, thin] (1,0) -- (1,2);
\draw[popline, dashed, thin] (0,1) -- (2,1);
\draw[popblue, thick, rounded corners=2pt] (0.12,0.15) rectangle (0.85,1.85);
\node[popblue, font=\small, align=center] at (0.48,1.0) {\zh{贝}\\ traits 1--4};
\draw[popgreen, thick, rounded corners=2pt] (0.95,0.15) rectangle (1.88,1.85);
\node[popgreen, font=\small, align=center] at (1.42,1.0) {\zh{曾}\\ traits 5--13};
\node[below, font=\small] at (1,-0.1) {\zh{赠} zèng (13 traits)};
\end{scope}
\end{tikzpicture}
\legende{Ordre des traits de \zh{广}, \zh{产} et \zh{赠} : les numéros rouges indiquent la succession des traits ; pour \zh{赠}, on écrit d'abord la partie gauche \zh{贝} (4 traits), puis la partie droite \zh{曾} (9 traits).}
\end{center}
\end{popfigure}

\courssub{Caractères proches : ne pas confondre}

Les francophones confondent souvent des caractères qui se ressemblent visuellement, comme nous confondons parfois « b » et « d ». Voici les trois paires les plus dangereuses de notre thème.

\textbf{1. 广 (guǎng, large) contre 庄 (zhuāng, village).} La différence tient à un seul trait : \zh{庄} contient \zh{土} (tǔ, « terre ») sous le toit. Un village, c'est de la \emph{terre} sous un toit. Si vous oubliez le \zh{土}, votre « village » devient « large ».

\textbf{2. 产 (chǎn, produire) contre 昌 (chāng, prospérer).} \zh{昌} est composé de deux \zh{日} (rì, « soleil ») superposés : deux soleils, c'est la prospérité. \zh{产} n'a aucun \zh{日} : il finit par un long trait horizontal et un trait qui descend vers la gauche. Attention aussi aux pinyin proches : \emph{chǎn} et \emph{chāng} se distinguent par la finale et le ton.

\textbf{3. 赠 (zèng, offrir) contre 赃 (zāng, butin).} Les deux contiennent \zh{贝} (l'argent), mais la partie droite diffère : \zh{曾} (céng) pour le cadeau, \zh{庄} pour le butin. Nuance socioculturelle amusante : le cadeau et le vol partagent le radical de l'argent — tout dépend de la manière dont on l'obtient !

\begin{attention}[Le piège du couple 买 / 卖]
\zh{卖} (mài, « vendre ») s'écrit avec \zh{买} (mǎi, « acheter ») plus un \zh{十} au-dessus. Deux pièges guettent les francophones :
\begin{enumerate}
\item Le trait horizontal du \zh{十} \textbf{ne dépasse pas à gauche} : il reste bien centré au-dessus du caractère. S'il déborde, le caractère devient illisible.
\item N'inversez pas les tons : \emph{mǎi} (3\up{e} ton) = acheter, \emph{mài} (4\up{e} ton) = vendre. Une erreur de ton dans une affiche peut transformer « nous vendons » en « nous achetons » !
\end{enumerate}
Astuce de mémorisation : celui qui \emph{vend} a quelque chose \emph{en plus} (le \zh{十}) à proposer ; celui qui \emph{achète} n'a rien au-dessus.
\end{attention}

\courssub{Les règles de proportion : l'esthétique du caractère}

Un caractère chinois s'inscrit dans un carré imaginaire. Deux règles de proportion s'appliquent à notre thème :

\begin{propriete}[Proportions : toit compact, base stable]
\begin{enumerate}
\item Pour les caractères à radical \zh{广} (comme \zh{广} lui-même, \zh{店}, \zh{库}) : la \textbf{partie supérieure est compacte} (le point et le toit restent serrés) et la \textbf{base est large}, bien ouverte, comme les fondations d'une maison. Un toit trop large écrase le caractère ; une base trop étroite le fait « tomber ».
\item Pour \zh{产} : le \textbf{dernier trait horizontal est long}, nettement plus long que les traits du haut. C'est lui qui « porte » tout le caractère, comme une étagère. S'il est trop court, le caractère paraît pencher.
\item Pour les caractères gauche-droite comme \zh{赠} et \zh{优} : la partie gauche (\zh{贝}, \zh{亻}) occupe environ \textbf{un tiers} de la largeur, la partie droite les deux tiers.
\end{enumerate}
\end{propriete}

\courssub{Du caractère au mot : former des mots composés}

En chinois moderne, la plupart des mots sont composés de deux caractères. Chacun de nos douze caractères est une brique qui s'assemble. Voici cinq mots composés essentiels pour votre affiche publicitaire :

\begin{exemplebox}[Cinq mots composés du marketing]
\begin{itemize}
\item \zh{广告} guǎnggào : « publicité » (large + annoncer = annoncer largement).
\item \zh{产品} chǎnpǐn : « produit » (produire + article).
\item \zh{优质} yōuzhì : « de qualité supérieure » (excellent + qualité).
\item \zh{优惠} yōuhuì : « promotion, prix préférentiel » (excellent + avantage).
\item \zh{赠品} zèngpǐn : « cadeau, article offert » (offrir + article).
\item \zh{限量} xiànliàng : « en édition limitée » (limiter + quantité).
\end{itemize}
Exemple de phrase d'affiche : \zh{优质产品，优惠价格，限量赠品！} \emph{Yōuzhì chǎnpǐn, yōuhuì jiàgé, xiànliàng zèngpǐn !} « Produits de qualité supérieure, prix promotionnels, cadeaux en édition limitée ! »
\end{exemplebox}

\begin{savaistu}[D'où vient le caractère 量 ?]
Le caractère \zh{量} (liàng, « quantité, mesure ») combine \zh{日} (rì, « soleil ») en haut et \zh{里} (lǐ, « dedans, village ») en bas. Une étymologie populaire l'interprète comme « mesurer le soleil depuis l'intérieur du village » : estimer, évaluer. On retrouve \zh{量} dans \zh{数量} (shùliàng, « quantité »), \zh{质量} (zhìliàng, « qualité » — littéralement « mesure de la nature ») et \zh{限量} (xiànliàng, « édition limitée »). Comprendre la logique interne des caractères est le meilleur moyen de ne plus jamais les oublier.
\end{savaistu}

\courssub{Copie guidée : de la grille à l'écriture libre}

\begin{methode}[S'entraîner efficacement en trois étapes]
\begin{enumerate}
\item \textbf{Grille de 9 cases} (九宫格 jiǔgōnggé) : tracez chaque caractère dans un grand carré divisé en neuf cases. Les lignes de la grille vous aident à placer les traits : le point de \zh{广} tombe dans la case centrale haute, la base occupe les deux cases du bas.
\item \textbf{Copie en cases simples} (田字格 tiánzìgé) : recopiez chaque caractère cinq fois en récitant l'ordre des traits à voix haute.
\item \textbf{Écriture libre avec auto-vérification} : écrivez le mot composé de mémoire (\zh{广告}, \zh{产品}…), puis vérifiez trois points : le radical est-il correct ? Le nombre de traits est-il bon ? Les proportions sont-elles respectées (base large, tiers/deux tiers) ?
\end{enumerate}
\end{methode}

\begin{exoresolu}[Identifier et corriger les erreurs d'écriture]
Énoncé : un élève de Terminale A a écrit une affiche contenant les mots suivants. Trois d'entre eux contiennent une erreur d'écriture ou de choix de caractère. Trouvez-les et corrigez-les : (1) \zh{广告} (2) \zh{庄品} (3) \zh{优质} (4) \zh{赃品} pour dire « cadeau » (5) \zh{限量}.
\tcblower
Solution détaillée :
\begin{enumerate}
\item \zh{广告} : correct. Radical \zh{广} en 3 traits, \zh{告} en 7 traits.
\item \zh{庄品} : \textbf{erreur}. L'élève a écrit \zh{庄} (zhuāng, « village ») au lieu de \zh{产} (chǎn, « produire »). Correction : \zh{产品} (chǎnpǐn, « produit »). Rappel : \zh{庄} contient \zh{土} (terre), \zh{产} finit par un long horizontal.
\item \zh{优质} : correct (excellent + qualité).
\item \zh{赃品} : \textbf{erreur de sens}. \zh{赃} (zāng) signifie « butin, biens volés » ! Pour « cadeau », il faut \zh{赠品} (zèngpǐn), avec \zh{曾} à droite du radical \zh{贝}.
\item \zh{限量} : correct (limiter + quantité = édition limitée).
\end{enumerate}
\end{exoresolu}

\begin{exercice}[À vous d'écrire]
\begin{enumerate}
\item Recopiez cinq fois chacun des caractères \zh{广}, \zh{产}, \zh{买}, \zh{卖}, \zh{赠} dans une grille de 9 cases, en récitant l'ordre des traits.
\item Complétez les mots composés avec le bon caractère : \zh{广\ \ } (publicité), \zh{\ \ 品} (produit), \zh{优\ \ } (promotion), \zh{\ \ 品} (cadeau), \zh{限\ \ } (édition limitée).
\item Écrivez de mémoire la phrase d'affiche : « Produits de qualité supérieure, cadeaux en édition limitée ! »
\end{enumerate}
\end{exercice}

\begin{corrige}[Exercice : À vous d'écrire]
\begin{enumerate}
\item Auto-vérification : \zh{广} = 3 traits (point, horizontal, coude) ; \zh{产} = 6 traits ; \zh{买} = 6 traits ; \zh{卖} = 8 traits (\zh{买} + \zh{十} centré, sans dépasser à gauche) ; \zh{赠} = 16 traits (\zh{贝} à gauche, un tiers de la largeur).
\item \zh{广告} (guǎnggào) ; \zh{产品} (chǎnpǐn) ; \zh{优惠} (yōuhuì) ; \zh{赠品} (zèngpǐn) ; \zh{限量} (xiànliàng).
\item \zh{优质产品，限量赠品！} \emph{Yōuzhì chǎnpǐn, xiànliàng zèngpǐn !} Vérifiez : \zh{优} a le radical \zh{亻} à gauche ; \zh{量} a \zh{日} en haut et \zh{里} en bas.
\end{enumerate}
\end{corrige}

\begin{aretenir}[L'essentiel de la section]
\begin{itemize}
\item Douze caractères cibles : \zh{广、告、产、品、质、优、惠、买、卖、赠、限、量}.
\item Le radical \zh{贝} (coquillage) signale l'argent et le commerce ; le radical \zh{广} signale le toit et le lieu.
\item Ordre des traits : haut avant bas, gauche avant droite, horizontal avant vertical, point souvent en dernier.
\item Paires à ne pas confondre : \zh{广}/\zh{庄}, \zh{产}/\zh{昌}, \zh{赠}/\zh{赃}, et le couple \zh{买}/\zh{卖} (tons 3 et 4).
\item Proportions : toit compact et base large pour \zh{广} ; long trait horizontal final pour \zh{产} ; gauche = un tiers pour \zh{赠} et \zh{优}.
\end{itemize}
\end{aretenir}

Maintenant que vos caractères sont nets, équilibrés et corrects, vous êtes prêts à les mettre en phrases : la section suivante vous entraînera à la traduction du français vers le chinois (thème), première étape de la rédaction de votre propre affiche.

\coursec{Exercices de thème (français → chinois)}

Dans la section précédente, nous avons appris à tracer correctement les caractères clés du marketing (\zh{优}, \zh{惠}, \zh{赠}, \zh{扫}, \zh{促}) en respectant les radicaux et l'ordre des traits. Savoir \emph{écrire} les caractères ne suffit pas : il faut maintenant savoir les \emph{employer} dans des phrases complètes. C'est exactement l'objet du \cle{thème} : traduire du français vers le chinois. C'est l'exercice le plus exigeant du Bac, car il oblige à \emph{penser en chinois} et non à calquer le français. Nous allons procéder méthodiquement, avec dix phrases progressives, du simple constat (« Ce produit est bio. ») jusqu'à la phrase complexe à double articulation (« Non seulement il est bio, mais en plus il est en promo. »).

\courssub{La méthode du thème en cinq étapes}

Avant de traduire la moindre phrase, observons deux exemples réussis.

\begin{exemplebox}[Deux traductions modèles]
\zh{这茶不但有机，而且好喝。}\\
\emph{Zhè chá búdàn yǒujī, érqiě hǎohē.}\\
« Ce thé est non seulement bio, mais aussi délicieux. »\\[0.4em]
\zh{买二赠一！}\\
\emph{Mǎi èr zèng yī!}\\
« Achetez deux, recevez un cadeau ! » (litt. « achetez deux, offrez un »)
\end{exemplebox}

Que remarque-t-on ? Dans la première phrase, le français « non seulement\ldots mais aussi » devient le couple de connecteurs \zh{不但\ldots 而且\ldots}. Dans la seconde, le français utilise deux verbes à l'infinitif ; le chinois condense tout en quatre caractères, sans sujet, avec deux verbes d'action. Le thème n'est donc jamais une traduction mot à mot : c'est une \emph{reconstruction} de la phrase selon la logique chinoise.

\begin{methode}[Traduire une phrase en thème : les 5 étapes]
\begin{enumerate}
\item \textbf{Identifier le noyau de la phrase} : qui fait quoi ? (sujet + verbe + complément essentiel). Dans « Ce produit camerounais est bio et pas cher », le noyau est « produit + est + bio ».
\item \textbf{Choisir la structure chinoise adaptée} : phrase adjectivale (\zh{这个产品很有机}), structure en \zh{被} pour le passif, topicalisation pour mettre l'objet en tête, couple de connecteurs pour les articulations logiques.
\item \textbf{Insérer le vocabulaire précis} : vérifier chaque mot dans le lexique du thème (\zh{有机} « bio », \zh{优惠} « promotion », \zh{优质} « de qualité »).
\item \textbf{Vérifier l'ordre des compléments} : en chinois, les compléments de temps et de lieu se placent \emph{avant} le verbe ; l'adverbe précède l'adjectif ou le verbe.
\item \textbf{Relire à voix haute} : si la phrase sonne bizarre à l'oreille, elle est probablement fautive. Vérifier les tons et les caractères tracés.
\end{enumerate}
\end{methode}

\begin{propriete}[Verbes d'action pour l'appel à l'action]
Dans une affiche publicitaire, le thème doit privilégier les \cle{verbes d'action} (\zh{买} « acheter », \zh{尝} « goûter », \zh{扫} « scanner », \zh{来} « venir », \zh{抢} « se précipiter pour acheter ») plutôt que les verbes d'état (\zh{是}, \zh{有}). Le français dit « Venez profiter de nos offres » ; le chinois dira \zh{快来买吧！} (\emph{Kuài lái mǎi ba!} « Venez vite acheter ! ») : plus direct, plus percutant, plus court. Sur une affiche, chaque caractère compte.
\end{propriete}

\courssub{Transposer les structures françaises : passif, relatives, topicalisation}

Le français et le chinois ne construisent pas leurs phrases de la même façon. Trois cas posent systématiquement problème aux élèves francophones.

\textbf{1. Le passif français.} « Ce miel est produit à Buea. » Le français utilise « être + participe passé ». Le chinois utilise la marque \zh{被} : \zh{这种蜂蜜是在布埃亚生产的。} (\emph{Zhè zhǒng fēngmì shì zài Bù'āiyà shēngchǎn de.}). Attention : en chinois, on n'utilise \zh{被} que lorsque la subordination passive est vraiment nécessaire ; souvent, une simple topicalisation suffit.

\textbf{2. La proposition relative.} « Le thé \emph{que nous vendons} est bio. » En chinois, le complément déterminatif se place \emph{avant} le nom, suivi de \zh{的} : \zh{我们卖的茶是有机的。} (\emph{Wǒmen mài de chá shì yǒujī de.} « Le thé que nous vendons est bio. »). L'ordre est inversé par rapport au français : \emph{que nous vendons} → \zh{我们卖的} + nom.

\textbf{3. La topicalisation.} « Ce produit, tout le monde l'aime. » Le chinois adore mettre l'objet en tête de phrase : \zh{这个产品，大家都喜欢。} (\emph{Zhège chǎnpǐn, dàjiā dōu xǐhuan.}). C'est la structure la plus naturelle pour une affiche, car elle met le produit en avant dès le premier mot.

\begin{center}
\begin{tabularx}{\linewidth}{|X|X|X|}
\hline
\rowcolor{popblueL} Structure française & Structure chinoise & Exemple chinois + traduction \\
\hline
Passif : « est produit à\ldots » & \zh{是 + 在\ldots + 生产的} ou \zh{被} & \zh{这咖啡是在西部省生产的。} « Ce café est produit dans l'Ouest. » \\
\hline
Relative : « le thé que\ldots » & complément + \zh{的} + nom & \zh{我们卖的茶很好。} « Le thé que nous vendons est bon. » \\
\hline
Mise en relief : « ce produit, \ldots » & objet en tête (topicalisation) & \zh{这个手机，又便宜又好用。} « Ce téléphone, il est à la fois bon marché et pratique. » \\
\hline
« non seulement\ldots mais aussi » & \zh{不但\ldots 而且\ldots} & \zh{不但便宜，而且优质。} « Non seulement bon marché, mais aussi de qualité. » \\
\hline
« de plus en plus + adj. » & \zh{越来越 + adj.} & \zh{越来越受欢迎。} « De plus en plus apprécié. » \\
\hline
\end{tabularx}
\end{center}

\begin{attention}[« Il y a une promotion » : le piège du mot à mot]
Ne traduisez jamais « Il y a une promotion » par \zh{有促销} : c'est trop vague et peu naturel. Préférez \zh{正在搞促销} (\emph{zhèngzài gǎo cùxiāo}, « une promotion est en cours ») ou \zh{有优惠活动} (\emph{yǒu yōuhuì huódòng}, « il y a une offre promotionnelle »). De même, ne confondez pas \zh{优惠} (\emph{yōuhuì}, « avantageux, promotion ») et \zh{优质} (\emph{yōuzhì}, « de haute qualité ») : les deux commencent par \zh{优}, mais \zh{优惠} parle du \emph{prix}, \zh{优质} parle de la \emph{qualité}. Erreur classique au Bac : « produit de qualité » se dit \zh{优质产品}, jamais \zh{优惠产品} (qui voudrait dire « produit en promo »).
\end{attention}

\courssub{Tableau de transposition : français, chinois attendu, pièges}

\begin{popfigure}
\begin{tikzpicture}[
  cell/.style={draw=popline, fill=popcream, rounded corners=2pt, inner sep=4pt, align=left, font=\small},
  head/.style={draw=popblue, fill=popblueL, rounded corners=2pt, inner sep=4pt, align=center, font=\small\bfseries}
]
\node[head, minimum width=4.2cm] (h1) at (0,0) {Français};
\node[head, minimum width=5.2cm, right=0.15cm of h1] (h2) {Chinois attendu};
\node[head, minimum width=4.6cm, right=0.15cm of h2] (h3) {Piège à éviter};
\node[cell, minimum width=4.2cm, below=0.12cm of h1] (a1) {« bio »};
\node[cell, minimum width=5.2cm, below=0.12cm of h2] (a2) {有机 \emph{yǒujī}};
\node[cell, minimum width=4.6cm, below=0.12cm of h3] (a3) {pas 生物 \emph{shēngwù} (= biologie)};
\node[cell, minimum width=4.2cm, below=0.12cm of a1] (b1) {« en promotion »};
\node[cell, minimum width=5.2cm, below=0.12cm of a2] (b2) {有优惠 / 打折 \emph{dǎzhé}};
\node[cell, minimum width=4.6cm, below=0.12cm of a3] (b3) {pas 优质 (= de qualité)};
\node[cell, minimum width=4.2cm, below=0.12cm of b1] (c1) {« acheter 2, 1 offert »};
\node[cell, minimum width=5.2cm, below=0.12cm of b2] (c2) {买二赠一 \emph{mǎi èr zèng yī}};
\node[cell, minimum width=4.6cm, below=0.12cm of b3] (c3) {pas de phrase longue mot à mot};
\node[cell, minimum width=4.2cm, below=0.12cm of c1] (d1) {« qualité garantie »};
\node[cell, minimum width=5.2cm, below=0.12cm of c2] (d2) {质量保证 \emph{zhìliàng bǎozhèng}};
\node[cell, minimum width=4.6cm, below=0.12cm of c3] (d3) {ordre : nom + nom, pas de 的};
\node[cell, minimum width=4.2cm, below=0.12cm of d1] (e1) {« scannez le code »};
\node[cell, minimum width=5.2cm, below=0.12cm of d2] (e2) {扫码 \emph{sǎo mǎ}};
\node[cell, minimum width=4.6cm, below=0.12cm of d3] (e3) {扫 = verbe d'action, pas « regarder »};
\end{tikzpicture}
\legende{Tableau de transposition français--chinois : la bonne réponse et le piège francophone correspondant.}
\end{popfigure}

\courssub{Série de dix phrases progressives (correction collective)}

Voici la série de thème de la leçon. Les phrases vont du plus simple au plus complexe. Traduis-les d'abord seul, puis compare avec la correction.

\begin{exoresolu}[Phrases 1 à 5 : niveau simple à intermédiaire]
Énoncé. Traduire en chinois :\\
1. Ce produit est bio.\\
2. Notre thé est délicieux et pas cher.\\
3. Ce miel est produit à Buea.\\
4. Il y a une promotion cette semaine.\\
5. Achetez deux, recevez un cadeau.\\
\tcblower
Solutions détaillées :\\
1. \zh{这个产品是有机的。} \emph{Zhège chǎnpǐn shì yǒujī de.} — Structure \zh{是\ldots 的} pour insister sur la qualité. Piège : ne pas écrire \zh{生物}.\\
2. \zh{我们的茶又好喝又便宜。} \emph{Wǒmen de chá yòu hǎohē yòu piányi.} — Couple \zh{又\ldots 又\ldots} pour deux adjectifs coordonnés.\\
3. \zh{这种蜂蜜是在布埃亚生产的。} \emph{Zhè zhǒng fēngmì shì zài Bù'āiyà shēngchǎn de.} — Passif transposé par \zh{是在\ldots 生产的} ; le lieu précède le verbe ; classificateur \zh{种} pour « ce type de ».\\
4. \zh{本周有优惠活动。} \emph{Běn zhōu yǒu yōuhuì huódòng.} — Le complément de temps \zh{本周} en tête ; jamais \zh{有促销} seul.\\
5. \zh{买二赠一！} \emph{Mǎi èr zèng yī!} — Formule publicitaire figée en quatre caractères ; inutile de traduire « recevez » mot à mot.
\end{exoresolu}

\begin{exoresolu}[Phrases 6 à 10 : niveau complexe]
Énoncé. Traduire en chinois :\\
6. Le café que nous vendons vient de l'Ouest du Cameroun.\\
7. Ce produit est de plus en plus apprécié des jeunes.\\
8. Non seulement il est bio, mais en plus il est en promo.\\
9. Si vous scannez ce code, vous recevrez un cadeau.\\
10. Venez vite goûter notre nouveau jus naturel !\\
\tcblower
Solutions détaillées :\\
6. \zh{我们卖的咖啡来自喀麦隆西部。} \emph{Wǒmen mài de kāfēi láizì Kāmàilóng xībù.} — Relative transposée : \zh{我们卖的} + \zh{咖啡} ; \zh{来自} « venir de ».\\
7. \zh{这个产品越来越受年轻人的欢迎。} \emph{Zhège chǎnpǐn yuèláiyuè shòu niánqīngrén de huānyíng.} — \zh{越来越} + structure \zh{受\ldots 的欢迎} « être apprécié de ».\\
8. \zh{这个产品不但有机，而且有优惠。} \emph{Zhège chǎnpǐn búdàn yǒujī, érqiě yǒu yōuhuì.} — Couple \zh{不但\ldots 而且\ldots} ; « en promo » = \zh{有优惠}, pas \zh{优质}.\\
9. \zh{如果您扫这个码，就可以得到礼物。} \emph{Rúguǒ nín sǎo zhège mǎ, jiù kěyǐ dédào lǐwù.} — Couple \zh{如果\ldots 就\ldots} ; verbe d'action \zh{扫}.\\
10. \zh{快来尝尝我们的新天然果汁吧！} \emph{Kuài lái chángchang wǒmen de xīn tiānrán guǒzhī ba!} — Impératif publicitaire : \zh{快来} + redoublement du verbe \zh{尝尝} (goûter un peu) + particule finale \zh{吧}.
\end{exoresolu}

\begin{attention}[Erreurs types relevées à la correction collective]
\begin{itemize}
\item \textbf{Ordre des mots} : « cette semaine » doit précéder le verbe : \zh{本周有优惠}, jamais \zh{有优惠本周}.
\item \textbf{Confusion lexicale} : \zh{优惠} (prix) contre \zh{优质} (qualité) ; \zh{有机} (bio) contre \zh{生物} (biologie).
\item \textbf{Calque du passif} : multiplier les \zh{被} alourdit la phrase ; préférer \zh{是\ldots 的} ou la topicalisation.
\item \textbf{Relative mal placée} : le déterminant précède toujours le nom avec \zh{的} : \zh{我们卖的茶}, jamais \zh{茶我们卖的}.
\item \textbf{Caractères fautifs} : vérifier le radical de \zh{赠} (radical 贝, l'argent, à gauche) et celui de \zh{惠}.
\end{itemize}
\end{attention}

\courssub{Auto-évaluation et extension « slogan »}

\begin{exercice}[Grille d'auto-évaluation simplifiée]
Après avoir traduit les dix phrases, évalue-toi sur 10 points :\\
\textbf{Vocabulaire (4 pts)} : \zh{有机}, \zh{优惠}, \zh{赠}, \zh{扫}, \zh{优质} correctement employés (0,8 pt par mot).\\
\textbf{Grammaire (4 pts)} : au moins 3 connecteurs corrects parmi \zh{不但\ldots 而且\ldots}, \zh{又\ldots 又\ldots}, \zh{越来越}, \zh{如果\ldots 就\ldots} ; ordre des compléments respecté.\\
\textbf{Caractères (2 pts)} : les 5 caractères cibles tracés sans erreur de radical.\\
Mon score : \ldots /10. Mon point à retravailler : \ldots
\end{exercice}

\begin{exercice}[Extension : transformer une phrase en slogan de 4 caractères]
Une affiche efficace se termine souvent par un slogan très court. Réécris la phrase 8 (« Non seulement il est bio, mais en plus il est en promo ») en un slogan de quatre caractères.\\
Pistes : garde les deux idées essentielles (bio + prix avantageux) et condense. Exemples possibles : \zh{有机又实惠} (\emph{yǒujī yòu shíhuì}, « bio et avantageux »), \zh{好货不贵} (\emph{hǎo huò bù guì}, « bon produit, pas cher »). À toi d'en créer un troisième !
\end{exercice}

\begin{corrige}[Extension slogan]
Il n'y a pas une seule bonne réponse : tout slogan de quatre caractères, grammaticalement correct, reprenant l'idée « bio + bon prix », est accepté. Critères : 4 caractères exactement, rythme équilibré (2+2), vocabulaire du thème. Exemple de correction : \zh{天然好价} (\emph{tiānrán hǎo jià}, « naturel, bon prix »). Le slogan doit pouvoir se lire d'une traite, à voix haute, en deux secondes.
\end{corrige}

\begin{savaistu}[Le thème au Bac]
Au baccalauréat, l'exercice de thème vaut en général 6 points, répartis ainsi : 2 points pour le vocabulaire, 2 points pour la grammaire, 1 point pour l'exactitude des caractères, 1 point pour la présentation générale. Autrement dit : une phrase parfaitement traduite mais écrite avec des caractères fautifs perd un point entier. D'où l'importance de la section 3 (écriture) avant de se lancer dans le thème !
\end{savaistu}

\begin{aretenir}[L'essentiel de la section]
\begin{itemize}
\item Le thème se fait en 5 étapes : noyau, structure, vocabulaire, ordre des compléments, relecture à voix haute.
\item Le passif français se transpose par \zh{被}, \zh{是\ldots 的} ou la topicalisation ; la relative se place avant le nom avec \zh{的}.
\item Les connecteurs de persuasion à maîtriser : \zh{不但\ldots 而且\ldots}, \zh{又\ldots 又\ldots}, \zh{越来越}, \zh{如果\ldots 就\ldots}.
\item Pièges lexicaux : \zh{优惠} (prix) contre \zh{优质} (qualité) ; \zh{有机} (bio) contre \zh{生物} (biologie).
\item Sur une affiche : verbes d'action, phrases courtes, slogan final de quatre caractères.
\end{itemize}
\end{aretenir}

\coursec{Exercices de version (chinois → français)}

Après avoir appris, dans la section précédente, à \emph{produire} du chinois publicitaire à partir du français (thème), nous inversons maintenant le sens du travail : la \cle{version} consiste à partir d'un texte chinois authentique — ici, deux extraits d'affiches publicitaires — pour le restituer en un français \textbf{fidèle au sens} et \textbf{naturel dans la langue}. C'est une compétence clé de l'épreuve du Baccalauréat : le correcteur juge à la fois la compréhension du chinois et la qualité du français rendu.

\courssub{Observer avant de traduire : deux affiches chinoises}

Lisons d'abord les deux textes supports, sans dictionnaire, pour en dégager le sens global.

\begin{textebox}[Affiche 1 — Un produit high-tech : des écouteurs]
新款无线耳机，音质超好，佩戴舒适。\\
\emph{Xīn kuǎn wúxiàn ěrjī, yīnzhì chāo hǎo, pèidài shūshì.}\\
限时优惠，扫码下单，立减五十元！\\
\emph{Xiànshí yōuhuì, sǎo mǎ xiàdān, lì jiǎn wǔshí yuán!}\\[4pt]
« Nouveaux écouteurs sans fil, qualité sonore excellente, port confortable. Offre à durée limitée : scannez le code pour commander, 50 yuans de réduction immédiate ! »
\end{textebox}

\begin{textebox}[Affiche 2 — Un produit local : les mangues du Cameroun]
喀麦隆芒果，又大又甜，天然种植。\\
\emph{Kāmàilóng mángguǒ, yòu dà yòu tián, tiānrán zhòngzhí.}\\
不但新鲜，而且便宜，快来选购吧！\\
\emph{Búdàn xīnxiān, érqiě piányi, kuài lái xuǎngòu ba!}\\[4pt]
« Mangues du Cameroun, à la fois grosses et sucrées, cultivées naturellement. Non seulement fraîches, mais encore bon marché : venez vite les choisir ! »
\end{textebox}

\begin{center}
\begin{tabularx}{\linewidth}{|X|X|X|X|}
\hline
\rowcolor{popblueL} Caractères & Pinyin & Nature & Traduction \\
\hline
耳机 & ěrjī & nom & écouteurs \\
\hline
无线 & wúxiàn & adj. & sans fil \\
\hline
音质 & yīnzhì & nom & qualité sonore \\
\hline
佩戴 & pèidài & verbe & porter (sur soi) \\
\hline
优惠 & yōuhuì & nom / adj. & promotion, avantage, prix réduit \\
\hline
扫码 & sǎo mǎ & verbe + compl. & scanner le code (QR) \\
\hline
下单 & xiàdān & verbe & passer commande \\
\hline
立减 & lì jiǎn & verbe & réduire immédiatement \\
\hline
打折 & dǎ zhé & verbe & faire une réduction, solder \\
\hline
芒果 & mángguǒ & nom & mangue \\
\hline
又大又甜 & yòu dà yòu tián & structure & à la fois gros et sucré \\
\hline
天然 & tiānrán & adj. & naturel \\
\hline
选购 & xuǎngòu & verbe & choisir et acheter \\
\hline
\end{tabularx}
\end{center}

\courssub{Repérer les marqueurs de structure avant de traduire}

Une affiche publicitaire chinoise, même courte, obéit à une architecture stable. Avant de traduire le moindre mot, on \textbf{segmente} le texte en zones fonctionnelles :

\begin{enumerate}
\item le \cle{slogan} (口号 \emph{kǒuhào}) : phrase choc, souvent nominale ou rythmée ;
\item les \cle{arguments} (卖点 \emph{màidiǎn}) : qualités du produit, souvent en adjectifs juxtaposés ;
\item la \cle{condition} ou l'offre (优惠 \emph{yōuhuì}) : prix, réduction, durée limitée ;
\item l'\cle{appel à l'action} (呼吁 \emph{hūyù}) : impératif invitant à acheter, souvent terminé par 吧 \emph{ba} ou un point d'exclamation.
\end{enumerate}

\begin{popfigure}
\begin{tikzpicture}[
  zone/.style={draw, rounded corners=3pt, align=left, text width=6.2cm, font=\small, inner sep=5pt},
  tag/.style={font=\small\bfseries, text=white, rounded corners=2pt, inner sep=3pt}
]
\node[zone, fill=poppurpleL, draw=poppurple] (slogan) at (0,0)
  {\textbf{新款无线耳机，音质超好，佩戴舒适。}\\ \emph{Xīn kuǎn wúxiàn ěrjī, yīnzhì chāo hǎo, pèidài shūshì.}};
\node[tag, fill=poppurple, anchor=west] at ($(slogan.east)+(0.3,0)$) {SLOGAN + ARGUMENTS};

\node[zone, fill=popgoldL, draw=popgold] (offre) at (0,-2.1)
  {\textbf{限时优惠，立减五十元！}\\ \emph{Xiànshí yōuhuì, lì jiǎn wǔshí yuán!}};
\node[tag, fill=popgold, anchor=west] at ($(offre.east)+(0.3,0)$) {CONDITION / OFFRE};

\node[zone, fill=popgreenL, draw=popgreen] (appel) at (0,-4.2)
  {\textbf{扫码下单！}\\ \emph{Sǎo mǎ xiàdān!}};
\node[tag, fill=popgreen, anchor=west] at ($(appel.east)+(0.3,0)$) {APPEL À L'ACTION};

\draw[-{Stealth}, thick, popmuted] (slogan.south) -- (offre.north);
\draw[-{Stealth}, thick, popmuted] (offre.south) -- (appel.north);
\end{tikzpicture}
\legende{Segmentation visuelle de l'affiche 1 : slogan et arguments (violet), offre (or), appel à l'action (vert).}
\end{popfigure}

Ce repérage préalable évite l'erreur la plus grave en version : traduire mot à mot sans comprendre la \textbf{fonction} de chaque segment. Un slogan se traduit comme un slogan, un appel comme un appel.

\begin{methode}[La version publicitaire en 5 étapes]
\begin{enumerate}
\item \textbf{Lire globalement} : lire tout le texte deux fois, identifier le produit, la cible et le ton (dynamique, chaleureux, urgent).
\item \textbf{Segmenter en unités de sens} : découper en zones (slogan / arguments / condition / appel) et repérer les structures connues (不但…而且…, 又…又…, 快来…吧).
\item \textbf{Traduire mot à mot} : produire une première version littérale, au brouillon, sans se soucier du style.
\item \textbf{Reformuler en français publicitaire} : remplacer les tournures lourdes par le registre de la pub française (impératifs, phrases nominales, superlatifs).
\item \textbf{Vérifier le ton} : relire à voix haute. Est-ce que cela « sonne » comme une vraie affiche française ? Est-ce que cela donne envie d'acheter ?
\end{enumerate}
\end{methode}

\courssub{Gérer la polysémie : un mot chinois, plusieurs traductions françaises}

La version exige de choisir le \textbf{bon équivalent français selon le contexte}. Deux mots de nos affiches illustrent ce problème.

\begin{definition}[Polysémie de 优惠 \emph{yōuhuì}]
Selon le contexte, \zh{优惠} se traduit par : « promotion » (限时优惠 → « offre à durée limitée »), « prix réduit » (优惠价 → « prix promotionnel »), ou « avantage » (学生优惠 → « tarif étudiant »). Le mot à mot « préférentiel » est rarement naturel en français publicitaire.
\end{definition}

\begin{definition}[Polysémie de 扫码 \emph{sǎo mǎ}]
Littéralement « balayer le code ». En français courant : « scanner le code », « flasher le QR code ». Dans \zh{微信扫码} (\emph{Wēixìn sǎo mǎ}), on traduira « scanner le code WeChat » ou, si le public français ne connaît pas WeChat, « scanner le code avec l'application ». On ne traduit jamais WeChat par « message minuscule » : les noms propres de marques restent tels quels.
\end{definition}

\begin{center}
\begin{tabularx}{\linewidth}{|X|X|X|}
\hline
\rowcolor{popblueL} Expression chinoise & Contexte & Traduction française adaptée \\
\hline
限时优惠 & affiche high-tech & « Offre à durée limitée » \\
\hline
优惠价 & étiquette de prix & « Prix promo » \\
\hline
扫码下单 & boutique en ligne & « Scannez le code pour commander » \\
\hline
微信扫码 & réseau social chinois & « Scannez avec WeChat » \\
\hline
立减五十元 & promotion immédiate & « 50 yuans de réduction immédiate » \\
\hline
快来选购吧 & appel à l'action & « Venez vite choisir ! » / « Foncez ! » \\
\hline
\end{tabularx}
\end{center}

\begin{attention}[La corrélation 不但…而且…]
\zh{不但新鲜，而且便宜} (\emph{Búdàn xīnxiān, érqiě piányi}) se traduit littéralement par « non seulement frais, mais encore bon marché » — tournure correcte mais \textbf{soutenue}. En français publicitaire, on préfère souvent « à la fois frais et bon marché » ou « frais \emph{et} bon marché », plus fluide à l'affiche. Les deux sont acceptées au Bac ; choisissez selon le ton du texte. Erreur à éviter : ne pas traduire 而且 par « donc » ou « alors » — c'est une \textbf{addition}, pas une conséquence.
\end{attention}

\begin{attention}[Le piège des réductions : 八折 ne veut pas dire « moins 80 \% »]
En chinois, la réduction s'exprime en \textbf{dixièmes du prix conservé} : \zh{八折} (\emph{bā zhé}) signifie « on paie 80 \% du prix », donc « \textbf{moins 20 \%} » en français ; \zh{五折} (\emph{wǔ zhé}) = « moitié prix » ; \zh{九折} (\emph{jiǔ zhé}) = « moins 10 \% ». C'est le contresens le plus fréquent — et le plus sanctionné — dans les versions publicitaires. Structure : \textbf{numéral + 折} = « (numéral × 10) \% du prix initial ».
\end{attention}

\courssub{Le registre publicitaire français : rendre l'effet, pas seulement le sens}

\begin{propriete}[Traduction de l'effet perlocutoire]
Une affiche ne sert pas à informer : elle sert à \textbf{faire agir} (acheter, venir, scanner). La version doit donc rendre l'\cle{effet perlocutoire} du texte chinois, pas seulement son sens littéral. Une traduction exacte mais plate (« Venez et choisissez-achetez ») est moins bonne qu'une traduction adaptée (« Foncez ! »), car elle trahit la fonction du texte.
\end{propriete}

Trois outils du registre publicitaire français à mobiliser :

\begin{itemize}
\item \textbf{l'impératif} : 快来选购吧！ → « Venez vite choisir ! » / « Foncez ! » ;
\item \textbf{le superlatif ou l'intensif} : 音质超好 → « une qualité sonore exceptionnelle » (et non « super bonne », trop familier à l'écrit) ;
\item \textbf{la phrase nominale} (sans verbe) : 新款无线耳机 → « Nouveaux écouteurs sans fil » — le français publicitaire adore les titres sans verbe, exactement comme le chinois.
\end{itemize}

\begin{exemplebox}[Formules publicitaires types : chinois → français]
\zh{限量发售，售完即止。} \emph{Xiànliàng fāshòu, shòu wán jí zhǐ.} → « Édition limitée, vente jusqu'à épuisement des stocks. »\\[3pt]
\zh{扫码关注，惊喜不断。} \emph{Sǎo mǎ guānzhù, jīngxǐ bùduàn.} → « Scannez pour nous suivre : les surprises s'enchaînent ! »\\[3pt]
\zh{买一送一。} \emph{Mǎi yī sòng yī.} → « Un acheté, un offert. »\\[3pt]
\zh{先到先得。} \emph{Xiān dào xiān dé.} → « Premier arrivé, premier servi. »
\end{exemplebox}

Remarquez que le chinois publicitaire aime les \textbf{formules en quatre caractères} (售完即止, 先到先得) : le français les rend par des proverbes ou des formules toutes faites, jamais par un calque (« vendre finir alors arrêter » serait incompréhensible).

\begin{exoresolu}[Version guidée de l'affiche 2]
Traduisez en français publicitaire : \zh{喀麦隆芒果，又大又甜，天然种植。不但新鲜，而且便宜，快来选购吧！}
\tcblower
\textbf{Étape 1 — Lecture globale} : produit local (mangues du Cameroun), ton chaleureux et direct, cible : le client d'un marché.\\
\textbf{Étape 2 — Segmentation} : slogan + arguments (又大又甜，天然种植) ; argument renforcé (不但新鲜，而且便宜) ; appel (快来选购吧).\\
\textbf{Étape 3 — Mot à mot} : « Mangues du Cameroun, aussi grosses aussi sucrées, naturellement cultivées. Non seulement fraîches, mais encore bon marché, vite venez choisir-acheter ! »\\
\textbf{Étape 4 — Reformulation} : « Mangues du Cameroun, à la fois grosses et sucrées, cultivées naturellement. Fraîches \emph{et} bon marché : venez vite faire votre choix ! »\\
\textbf{Étape 5 — Vérification du ton} : phrases courtes, un impératif, un rythme d'affiche. On peut intensifier : « Foncez ! » à la place de « venez vite » si le ton général est très dynamique.\\
\textbf{Version finale retenue} : « Mangues du Cameroun : grosses, sucrées, cultivées naturellement. Fraîches et à petit prix — venez vite faire votre choix ! »
\end{exoresolu}

\begin{attention}[Pièges fréquents des francophones en version]
\begin{itemize}
\item \textbf{Calquer l'ordre chinois} : \zh{天然种植} ne se traduit pas « naturel plantation » mais « cultivées naturellement » (le français exige un participe ou une tournure passive).
\item \textbf{Traduire 吧 par « n'est-ce pas »} : dans un appel publicitaire, 吧 adoucit l'impératif ; le français le rend simplement par l'impératif avec point d'exclamation.
\item \textbf{Oublier les sous-entendus} : \zh{立减五十元} suppose « (le prix) est immédiatement réduit de 50 yuans » ; le français doit expliciter : « 50 yuans de réduction immédiate ».
\item \textbf{Traduire 惊喜不断 trop littéralement} : « les surprises ne se coupent pas » n'existe pas ; dites « les surprises s'enchaînent » ou « que des surprises ».
\end{itemize}
\end{attention}

\courssub{Relecture croisée en binôme : la grille fidélité / naturel}

En classe, la relecture se fait en binôme : chaque élève traduit seul, puis échange sa copie avec un voisin qui la corrige à l'aide de la grille suivante.

\begin{center}
\begin{tabularx}{\linewidth}{|X|X|X|}
\hline
\rowcolor{popblueL} Critère & Question de contrôle & Points \\
\hline
Fidélité — sens & Toutes les informations du chinois sont-elles présentes (produit, qualités, offre, appel) ? & /4 \\
\hline
Fidélité — structures & 不但…而且…, 又…又…, 限时 : les structures sont-elles rendues ? & /3 \\
\hline
Naturel — grammaire & Le français est-il correct (accords, temps, articles) ? & /3 \\
\hline
Naturel — registre & Cela « sonne »-t-il comme une vraie pub française (impératifs, rythme, formules) ? & /4 \\
\hline
Ton & L'affiche donne-t-elle envie d'acheter ? & /2 \\
\hline
Présentation & Ponctuation française, typographie soignée. & /1 \\
\hline
\end{tabularx}
\end{center}

\begin{experience}[Relecture croisée en binôme]
\begin{enumerate}
\item Chacun traduit l'affiche 1 (écouteurs) en 10 minutes, au propre.
\item Échangez les copies. Le relecteur annote avec la grille : il souligne en vert les réussites, en rouge les contresens, en bleu les tournures « pas naturelles ».
\item Le relecteur propose \textbf{une amélioration de style} minimum (ex. remplacer « Allez-y » par « Foncez ! »).
\item Chaque binôme lit sa meilleure version à la classe ; la classe vote pour l'affiche la plus convaincante.
\end{enumerate}
\end{experience}

\begin{savaistu}[Ce que les correcteurs du Bac valorisent]
À l'épreuve de chinois, une version qui « sonne » comme une vraie publicité française est toujours mieux notée qu'un calque exact mais terne. Les correcteurs apprécient les formules idiomatiques : « Foncez ! » plutôt que « Allez-y ! », « Un acheté, un offert » plutôt que « Achetez un, donnez un », « Jusqu'à épuisement des stocks » plutôt que « vendre jusqu'à ce que ce soit fini ». En Chine comme au Cameroun, la publicité est un art du rythme : les affiches chinoises jouent sur les formules de quatre caractères, les affiches camerounaises sur les slogans percutants en français ou en anglais (« Buyam sellam » au marché !). Traduire une pub, c'est donc traduire une \emph{énergie}, pas seulement des mots.
\end{savaistu}

\begin{exercice}[À vous de jouer]
Traduisez en français publicitaire soigné, en appliquant les 5 étapes de la méthode :\\
1. \zh{新款运动鞋，又轻又舒服，限时八折，快来买吧！} \emph{Xīn kuǎn yùndòngxié, yòu qīng yòu shūfu, xiànshí bā zhé, kuài lái mǎi ba!}\\
2. \zh{本地蜂蜜，天然纯正，不但好吃，而且健康。扫码下单，送货上门！} \emph{Běndì fēngmì, tiānrán chúnzhèng, búdàn hǎochī, érqiě jiànkāng. Sǎo mǎ xiàdān, sòng huò shàng mén!}\\
3. \zh{限量发售，售完即止，先到先得！} \emph{Xiànliàng fāshòu, shòu wán jí zhǐ, xiān dào xiān dé!}
\end{exercice}

\begin{corrige}[Exercice — Exercices de version]
\textbf{1.} Segmentation : slogan (新款运动鞋) ; arguments (又轻又舒服) ; offre (限时八折) ; appel (快来买吧). Version : « Nouvelles chaussures de sport, à la fois légères et confortables. Moins 20 \% pour une durée limitée : foncez ! » (Attention : 八折 = 80 \% du prix, donc « moins 20 \% » ou « 20 \% de réduction » — piège classique !)\\
\textbf{2.} « Miel local, pur et naturel. Non seulement délicieux, mais encore bon pour la santé ! Scannez le code pour commander : livraison à domicile ! » (送货上门 = « livrer la marchandise jusqu'à la porte » → « livraison à domicile ».)\\
\textbf{3.} « Édition limitée, jusqu'à épuisement des stocks : premier arrivé, premier servi ! » (formule toute faite du français pour 先到先得.)
\end{corrige}

\begin{aretenir}[L'essentiel de la section]
\begin{itemize}
\item La version publicitaire se fait en 5 étapes : lire globalement, segmenter, traduire mot à mot, reformuler en français publicitaire, vérifier le ton.
\item On repère d'abord les quatre zones : slogan, arguments, condition, appel à l'action.
\item Les mots polysémiques se traduisent selon le contexte : 优惠 → « promotion / prix réduit / avantage » ; 扫码 → « scanner / flasher le code » ; les marques (微信 WeChat) restent inchangées.
\item 不但…而且… : « non seulement… mais encore… » (soutenu) ou « à la fois… et… » (fluide).
\item Les réductions se lisent en dixièmes du prix : 八折 = « moins 20 \% », 五折 = « moitié prix ».
\item La version doit rendre l'\cle{effet perlocutoire} : inciter à acheter, avec impératifs, phrases nominales et formules idiomatiques françaises (« Foncez ! », « Un acheté, un offert »).
\end{itemize}
\end{aretenir}

\coursec{Grille d'évaluation et production finale}

Dans la section précédente, nous avons traduit du chinois vers le français des slogans et des arguments publicitaires. Il est maintenant temps de rassembler tout ce que vous avez appris dans cette leçon : comment votre affiche sera-t-elle notée ? Comment vérifier votre travail avant de le remettre ? Comment le présenter à l'oral ? Cette dernière section vous donne les outils de l'évaluation et de la production finale.

\courssub{La grille d'évaluation type Bac}

Au baccalauréat, une production écrite en chinois n'est pas notée « au feeling ». L'examinateur utilise une grille critériée : cinq critères, chacun évalué sur quatre niveaux. Connaître cette grille, c'est connaître les règles du jeu.

\begin{definition}[Les cinq critères de l'affiche publicitaire]
\begin{enumerate}
\item \cle{Vocabulaire spécifique} : richesse et exactitude du lexique du marketing et du produit (\zh{有机} \emph{yǒujī} « bio », \zh{营养} \emph{yíngyǎng} « nutrition », \zh{订购} \emph{dìnggòu} « commander », \zh{折扣} \emph{zhékòu} « remise »…).
\item \cle{Structures grammaticales} : correction et variété des structures (comparatif \zh{比}, \zh{不但……而且……}, \zh{越……越……}, impératif).
\item \cle{Caractères et présentation} : écriture correcte, traits nets, mise en page soignée.
\item \cle{Cohérence et persuasion} : logique interne et force d'incitation à l'achat.
\item \cle{Originalité} : slogan créatif, angle culturel, élément visuel (logo, QR code fictif).
\end{enumerate}
\end{definition}

\begin{definition}[Critère « Cohérence et persuasion »]
Ce critère mesure deux choses : la \cle{logique interne} de l'affiche, qui doit suivre le schéma \textbf{problème → solution → action} (le consommateur a un besoin, le produit y répond, on lui dit quoi faire), et le \cle{ton engageant} : phrases courtes, rythme, impératifs, promesses crédibles. Une affiche belle mais qui ne dit pas où ni comment acheter perd des points sur ce critère.
\end{definition}

\begin{popfigure}
\begin{tikzpicture}[x=2.45cm,y=0.92cm, every node/.style={font=\scriptsize, align=center}]
\fill[popblue] (0,5) rectangle (1.05,6); \node[white, font=\scriptsize\bfseries] at (0.525,5.5) {Critère};
\fill[poporange] (1.05,5) rectangle (2,6); \node[white, font=\scriptsize\bfseries] at (1.525,5.5) {Insatisfaisant};
\fill[popgold] (2,5) rectangle (3,6); \node[white, font=\scriptsize\bfseries] at (2.5,5.5) {Passable};
\fill[popgreen] (3,5) rectangle (4,6); \node[white, font=\scriptsize\bfseries] at (3.5,5.5) {Bien};
\fill[poppurple] (4,5) rectangle (5,6); \node[white, font=\scriptsize\bfseries] at (4.5,5.5) {Très bien};
% Ligne 1 : Vocabulaire
\fill[popblueL] (0,4) rectangle (1.05,5); \node[font=\scriptsize\bfseries, align=center] at (0.525,4.5){Vocabulaire\\spécifique};
\node[fill=poporangeL, align=center] at (1.525,4.5){Mots génériques\\(好, 便宜)\\seulement};
\node[fill=popgoldL, align=center] at (2.5,4.5){3--4 mots\\du thème,\\1--2 erreurs};
\node[fill=popgreenL, align=center] at (3.5,4.5){6--8 mots\\du thème\\corrects};
\node[fill=poppinkL, align=center] at (4.5,4.5){Lexique riche\\et précis\\(有机, 折扣…)};
% Ligne 2 : Structures
\fill[popblueL] (0,3) rectangle (1.05,4); \node[font=\scriptsize\bfseries, align=center] at (0.525,3.5){Structures\\grammaticales};
\node[fill=poporangeL, align=center] at (1.525,3.5){Phrases\\télégraphiques\\ou fautives};
\node[fill=popgoldL, align=center] at (2.5,3.5){Structures\\simples\\correctes};
\node[fill=popgreenL, align=center] at (3.5,3.5){1--2 structures\\complexes\\correctes};
\node[fill=poppinkL, align=center] at (4.5,3.5){不但…而且…,\\越…越…\\bien employés};
% Ligne 3 : Caractères
\fill[popblueL] (0,2) rectangle (1.05,3); \node[font=\scriptsize\bfseries, align=center] at (0.525,2.5){Caractères et\\présentation};
\node[fill=poporangeL, align=center] at (1.525,2.5){Caractères\\inventés ou\\illisibles};
\node[fill=popgoldL, align=center] at (2.5,2.5){Lisible,\\quelques\\traits faux};
\node[fill=popgreenL, align=center] at (3.5,2.5){Traits nets,\\proportions\\correctes};
\node[fill=poppinkL, align=center] at (4.5,2.5){Calligraphie\\soignée, mise\\en page claire};
% Ligne 4 : Cohérence
\fill[popblueL] (0,1) rectangle (1.05,2); \node[font=\scriptsize\bfseries, align=center] at (0.525,1.5){Cohérence et\\persuasion};
\node[fill=poporangeL, align=center] at (1.525,1.5){Pas de logique,\\pas d'appel\\à l'action};
\node[fill=popgoldL, align=center] at (2.5,1.5){Idées présentes\\mais peu\\liées};
\node[fill=popgreenL, align=center] at (3.5,1.5){Problème →\\solution →\\action visible};
\node[fill=poppinkL, align=center] at (4.5,1.5){Parcours clair,\\ton rythmé,\\incitation forte};
% Ligne 5 : Originalité
\fill[popblueL] (0,0) rectangle (1.05,1); \node[font=\scriptsize\bfseries] at (0.525,0.5) {Originalité};
\node[fill=poporangeL, align=center] at (1.525,0.5){Copie du\\modèle du\\manuel};
\node[fill=popgoldL, align=center] at (2.5,0.5){Slogan simple\\mais personnel};
\node[fill=popgreenL, align=center] at (3.5,0.5){Slogan créatif,\\logo ou QR\\fictif};
\node[fill=poppinkL, align=center] at (4.5,0.5){Angle culturel\\fort, slogan\\mémorable};
\draw[popdark, thick] (0,0) grid (5,6);
\end{tikzpicture}
\legende{Grille d'évaluation de l'affiche publicitaire : 5 critères, 4 niveaux.}
\end{popfigure}

\begin{aretenir}[Comment lire la grille]
Chaque critère vaut le même poids. Un « Très bien » partout donne la note maximale ; un seul critère « Insatisfaisant » fait chuter fortement la moyenne. La stratégie gagnante : ne négliger \textbf{aucun} critère, même l'originalité, qui coûte peu d'efforts (un slogan court et un logo fictif suffisent).
\end{aretenir}

\courssub{Analyse de deux productions d'élèves}

Observons d'abord deux affiches rédigées par des élèves de Terminale sur le même sujet : « Crée une affiche pour un thé vert bio camerounais destiné au marché chinois. »

\begin{exemplebox}[Production forte — affiche de Larissa]
\textbf{Slogan :} \zh{绿色有机，健康之选} \emph{(Lǜsè yǒujī, jiànkāng zhī xuǎn.)} « Vert et bio, le choix de la santé. »\\[2pt]
\textbf{Argument :} \zh{我们的茶不但口感纯正，而且营养满分。} \emph{(Wǒmen de chá búdàn kǒugǎn chúnzhèng, érqiě yíngyǎng mǎnfēn.)} « Notre thé a non seulement un goût pur, mais il est aussi très nutritif. »\\[2pt]
\textbf{Appel à l'action :} \zh{扫码订购，首单九折！} \emph{(Sǎo mǎ dìnggòu, shǒu dān jiǔ zhé!)} « Scannez le code pour commander, 10 \% de réduction sur la première commande ! »\\[4pt]
\textbf{Analyse critériée :}
\begin{itemize}
\item \textbf{Vocabulaire} : Très bien — \zh{有机}, \zh{口感}, \zh{营养}, \zh{订购}, \zh{折} : lexique marketing précis.
\item \textbf{Structures} : Très bien — \zh{不但……而且……} correctement employé, slogan parallèle en deux segments de quatre caractères.
\item \textbf{Caractères} : Bien — traits nets, proportions équilibrées, aucun caractère douteux.
\item \textbf{Cohérence} : Très bien — promesse (santé) → preuve (goût et nutrition) → action (scanner, commander, réduction).
\item \textbf{Originalité} : Très bien — slogan en 8 caractères symétriques (4 + 4), très chinois dans l'esprit, QR code fictif dessiné.
\end{itemize}
\end{exemplebox}

\begin{exemplebox}[Production faible — affiche de Marc]
\zh{这个茶很好喝，便宜，买吧。} \emph{(Zhège chá hěn hǎohē, piányi, mǎi ba.)} « Ce thé est très bon à boire, pas cher, achetez-le. »\\[4pt]
\textbf{Analyse critériée :}
\begin{itemize}
\item \textbf{Vocabulaire} : Insatisfaisant — uniquement des mots génériques (\zh{好喝} \emph{hǎohē} « bon à boire », \zh{便宜} \emph{piányi} « pas cher ») ; aucun mot du thème marketing.
\item \textbf{Structures} : Passable — phrases correctes mais juxtaposées, aucun connecteur.
\item \textbf{Caractères} : Passable — lisible, mais \zh{喝} écrit avec la clé de la bouche 口 à peine visible.
\item \textbf{Cohérence} : Insatisfaisant — pas de problème posé, pas de preuve, l'appel à l'action \zh{买吧} \emph{(mǎi ba)} « achetez donc » est sec et peu engageant.
\item \textbf{Originalité} : Insatisfaisant — aucun slogan, aucun élément visuel.
\end{itemize}
\textbf{Conseil pour Marc :} garder ses idées (goût, prix) mais les habiller : \zh{好喝又健康} \emph{(hǎohē yòu jiànkāng)} « bon et sain » au lieu de \zh{很好喝}, ajouter \zh{快来试试吧！} \emph{(Kuài lái shìshi ba!)} « Venez vite essayer ! » et un slogan court.
\end{exemplebox}

\begin{attention}[Le piège du caractère inventé]
Ne « devine » jamais l'écriture d'un caractère. Un caractère mal écrit peut devenir un \cle{caractère fantôme} qui n'existe pas : c'est une faute grave au Bac. Si tu hésites, \textbf{remplace par un synonyme maîtrisé} : \zh{优} \emph{yōu} (excellent) → \zh{好} \emph{hǎo} ; \zh{购} \emph{gòu} (acheter) → \zh{买} \emph{mǎi} ; \zh{惠} \emph{huì} (avantage) → \zh{便宜} \emph{piányi}. Un mot simple et correct vaut toujours mieux qu'un mot savant et faux.
\end{attention}

\courssub{La check-list d'auto-évaluation avant remise}

Avant de rendre ta copie, passe chaque item au crible. Cette check-list reprend exactement ce que l'examinateur cherche.

\begin{methode}[S'auto-évaluer en 5 minutes]
\begin{enumerate}
\item \textbf{Les 12 caractères cibles} du thème sont-ils présents et corrects ? (\zh{广告}, \zh{产品}, \zh{质量}, \zh{价格}, \zh{有机}, \zh{健康}, \zh{订购}, \zh{折扣}…)
\item \textbf{Au moins 3 connecteurs ou structures de persuasion} ? (\zh{不但……而且……}, \zh{越……越……}, \zh{因为……所以……})
\item \textbf{Le slogan fait-il 8 caractères ou moins} ? (Un slogan long n'est pas un slogan ; l'idéal est deux blocs de 4 caractères.)
\item \textbf{L'appel à l'action} est-il explicite ? (\zh{快来买吧！} \emph{Kuài lái mǎi ba!} « Venez vite acheter ! », \zh{扫码订购！} \emph{Sǎo mǎ dìnggòu!} « Scannez le code et commandez ! »)
\item \textbf{Le QR code fictif et le logo} sont-ils dessinés ?
\item \textbf{Aucun caractère incertain} n'a été laissé tel quel ?
\item \textbf{La hiérarchie visuelle} est-elle claire : grand titre, argument moyen, informations petites ?
\item \textbf{Le schéma problème → solution → action} est-il respecté ?
\item \textbf{Le ton} est-il engageant (impératif, rythme, exclamation) ?
\item \textbf{Les valeurs chinoises} (santé, famille, prestige) sont-elles visées ?
\end{enumerate}
Compte tes cases cochées : \textbf{si moins de 8 sur 10 (80 \%) sont validées, ne remets pas encore ta copie} — révise d'abord la zone faible identifiée.
\end{methode}

\begin{exoresolu}[Appliquer la check-list à une affiche]
\textbf{Énoncé.} Voici l'affiche de Divine : \zh{喀麦隆咖啡，又香又浓。因为质量好，所以大家都喜欢。快来买吧！} \emph{(Kāmàilóng kāfēi, yòu xiāng yòu nóng. Yīnwèi zhìliàng hǎo, suǒyǐ dàjiā dōu xǐhuan. Kuài lái mǎi ba!)} « Café du Cameroun, à la fois parfumé et corsé. Parce que la qualité est bonne, tout le monde l'aime. Venez vite acheter ! » Vérifiez sa copie avec la check-list et proposez deux améliorations.
\tcblower
\textbf{Solution détaillée.}
\begin{enumerate}
\item Caractères cibles : \zh{咖啡}, \zh{质量}, \zh{香}, \zh{浓} présents — environ 6 sur 12 : \textbf{insuffisant}.
\item Connecteurs : \zh{又……又……} et \zh{因为……所以……} = 2 structures : \textbf{presque} (il en faut 3).
\item Slogan : \zh{喀麦隆咖啡，又香又浓} = 9 caractères hors ponctuation : \textbf{trop long}.
\item Appel à l'action : \zh{快来买吧！} : \textbf{validé}.
\item Schéma problème → solution → action : solution et action présentes, mais \textbf{pas de problème/besoin posé}.
\end{enumerate}
\textbf{Améliorations proposées :}
\begin{itemize}
\item Raccourcir et équilibrer le slogan : \zh{香浓咖啡，健康之选} \emph{(Xiāng nóng kāfēi, jiànkāng zhī xuǎn)} « Café parfumé, le choix de la santé » — 8 caractères en deux blocs symétriques (4 + 4) ; version encore plus courte : \zh{香浓喀麦隆} \emph{(Xiāng nóng Kāmàilóng)} « Le Cameroun tout en arômes » (5 caractères).
\item Ajouter une structure et du vocabulaire marketing : \zh{我们的咖啡不但香浓，而且价格优惠。} \emph{(Wǒmen de kāfēi búdàn xiāng nóng, érqiě jiàgé yōuhuì.)} « Notre café est non seulement parfumé, mais aussi à prix avantageux. » On gagne \zh{不但……而且……}, \zh{价格} et \zh{优惠}.
\end{itemize}
\end{exoresolu}

\courssub{Conseils de mise en page}

Une affiche se lit en trois secondes. La mise en page est donc un critère de notation à part entière.

\begin{itemize}
\item \textbf{Police} : en chinois imprimé, on utilise le style \cle{Heiti} (黑体 \emph{hēitǐ}, SimHei), aux traits droits et épais, pour les titres ; le style Songti (宋体 \emph{sòngtǐ}) convient aux petits textes. À la main, imite le Heiti : traits réguliers, sans fioritures.
\item \textbf{Taille hiérarchisée} : le slogan en très grand (3 fois la taille du texte), l'argument en taille moyenne, le prix et les informations pratiques en petit.
\item \textbf{Espace négatif} : ne remplis pas toute la feuille. Le vide autour du slogan le fait respirer et attire l'œil. Une affiche surchargée est une affiche illisible.
\item \textbf{Logo fictif} : un cercle, une feuille de thé stylisée, un grain de café — simple et placé en haut à gauche ou au centre.
\item \textbf{QR code fictif} : un carré avec trois petits carrés aux coins suffit ; écris à côté \zh{扫码订购} \emph{(sǎo mǎ dìnggòu)} « scannez pour commander ».
\end{itemize}

\begin{center}
\begin{tabularx}{\linewidth}{|X|X|X|X|}
\hline
\rowcolor{popblueL} Caractères & Pinyin & Nature & Traduction \\
\hline
广告 & guǎnggào & nom & publicité, annonce \\
\hline
产品 & chǎnpǐn & nom & produit \\
\hline
质量 & zhìliàng & nom & qualité \\
\hline
价格 & jiàgé & nom & prix \\
\hline
优惠 & yōuhuì & adjectif/nom & avantageux ; réduction \\
\hline
折扣 & zhékòu & nom & remise, rabais \\
\hline
订购 & dìnggòu & verbe & commander \\
\hline
健康 & jiànkāng & nom/adjectif & santé ; sain \\
\hline
有机 & yǒujī & adjectif & biologique, bio \\
\hline
口感 & kǒugǎn & nom & goût (en bouche), texture \\
\hline
扫码 & sǎo mǎ & verbe + obj. & scanner le code (QR) \\
\hline
选择 & xuǎnzé & nom/verbe & choix ; choisir \\
\hline
\end{tabularx}
\end{center}

\courssub{Le volet interculturel : adapter le message aux valeurs chinoises}

Une bonne affiche n'est pas seulement correcte : elle parle à la culture du destinataire. Pour le marché chinois, trois valeurs dominent la publicité alimentaire.

\begin{itemize}
\item \textbf{La santé} (\zh{健康} \emph{jiànkāng}) : l'argument « bio » (\zh{有机} \emph{yǒujī}) et « naturel » (\zh{天然} \emph{tiānrán}) est roi. \zh{绿色有机，健康之选} touche directement cette corde.
\item \textbf{La famille} (\zh{家庭} \emph{jiātíng}) : montrer le produit partagé en famille. Exemple : \zh{全家人的健康选择} \emph{(Quánjiā rén de jiànkāng xuǎnzé)} « Le choix santé de toute la famille ».
\item \textbf{Le prestige} (\zh{面子} \emph{miànzi}, le « prestige/face ») : offrir un produit rare ou importé valorise celui qui l'offre. Exemple : \zh{来自喀麦隆高原的珍品} \emph{(Láizì Kāmàilóng gāoyuán de zhēnpǐn)} « Un trésor venu des hauts plateaux du Cameroun ».
\end{itemize}

\begin{savaistu}[Pourquoi le chiffre 9 dans « 九折 » ?]
En chinois, \zh{九折} \emph{(jiǔ zhé)} signifie littéralement « neuf dixièmes » : le client paie 90 \% du prix, soit 10 \% de réduction. Attention, c'est l'inverse du raisonnement français : \zh{八折} \emph{(bā zhé)} = 20 \% de réduction, \zh{五折} \emph{(wǔ zhé)} = moitié prix ! De plus, le chiffre 9 (\zh{九}, jiǔ) est homophone de \zh{久} (jiǔ, « longtemps ») : il symbolise la durée et est très apprécié dans le commerce chinois.
\end{savaistu}

\begin{savaistu}[L'examinateur lit tes caractères comme un calligraphe]
Au Bac, la lisibilité des caractères est réellement notée : traits nets, proportions équilibrées (chaque caractère occupe un carré imaginaire), et surtout \textbf{aucun « caractère fantôme »} — un caractère inventé qui n'existe dans aucun dictionnaire. Un trait de trop au-dessus de \zh{买} \emph{(mǎi, acheter)} le transforme en \zh{卖} \emph{(mài, vendre)} : le sens complet de ta phrase change ! Entraîne-toi à écrire lentement les 12 caractères cibles avant l'épreuve.
\end{savaistu}

\courssub{Préparation à l'oral : présenter son affiche en une minute}

L'affiche écrite débouche souvent sur une présentation orale. Une minute, c'est environ 120 caractères prononcés : il faut un plan strict.

\begin{methode}[Présenter son affiche en 1 minute]
\begin{enumerate}
\item \textbf{Accroche (10 s)} : nommer le produit et lire le slogan. \zh{大家好！我给大家介绍我们的新产品：喀麦隆绿茶。} \emph{(Dàjiā hǎo! Wǒ gěi dàjiā jièshào wǒmen de xīn chǎnpǐn : Kāmàilóng lǜchá.)} « Bonjour à tous ! Je vous présente notre nouveau produit : le thé vert du Cameroun. »
\item \textbf{Arguments (30 s)} : deux arguments avec connecteurs. \zh{我们的茶不但绿色有机，而且口感纯正。越喝越健康！} \emph{(Wǒmen de chá búdàn lǜsè yǒujī, érqiě kǒugǎn chúnzhèng. Yuè hē yuè jiànkāng!)} « Notre thé est non seulement vert et bio, mais il a aussi un goût pur. Plus on le boit, plus on est en bonne santé ! »
\item \textbf{Offre et action (15 s)} : prix, réduction, appel à l'action. \zh{现在扫码订购，首单九折！} \emph{(Xiànzài sǎo mǎ dìnggòu, shǒu dān jiǔ zhé!)} « Scannez maintenant le code pour commander : 10 \% de réduction sur la première commande ! »
\item \textbf{Conclusion (5 s)} : remercier. \zh{谢谢大家！} \emph{(Xièxie dàjiā!)} « Merci à tous ! »
\end{enumerate}
Parle lentement, soigne les tons, et montre les éléments de l'affiche du doigt au moment où tu les cites.
\end{methode}

\begin{experience}[Concours d'affiches en classe]
Par groupes de trois : (1) chaque élève finalise son affiche et la vérifie avec la check-list ; (2) chacun présente son affiche en une minute devant son groupe ; (3) le groupe note chaque présentation avec la grille à cinq critères et élit la meilleure affiche ; (4) les gagnants présentent devant la classe entière, qui vote pour « l'affiche la plus persuasive ». Le professeur clôt par un commentaire critérié collectif.
\end{experience}

\begin{exercice}[Production finale]
Sujet : tu es stagiaire dans une entreprise de Douala qui exporte du miel bio vers la Chine. Rédige une affiche publicitaire complète en chinois : slogan (8 caractères maximum, idéalement 4 + 4), deux arguments avec connecteurs, une offre promotionnelle avec \zh{折}, un appel à l'action avec \zh{扫码}, un logo et un QR code fictifs. Vérifie ta production avec la check-list des 10 points, puis prépare ta présentation orale d'une minute.
\end{exercice}

\begin{corrige}[Production finale — éléments de réponse]
Exemple de production attendue (niveau « Très bien ») :\\
Slogan : \zh{天然蜂蜜，健康之选} \emph{(Tiānrán fēngmì, jiànkāng zhī xuǎn)} « Miel naturel, le choix de la santé » (8 caractères, 4 + 4 ; version courte : \zh{天然好蜂蜜} \emph{Tiānrán hǎo fēngmì}, 5 caractères).\\
Arguments : \zh{我们的蜂蜜不但味道纯正，而且营养丰富。} \emph{(Wǒmen de fēngmì búdàn wèidào chúnzhèng, érqiě yíngyǎng fēngfù.)} « Notre miel a non seulement un goût pur, mais il est aussi très nutritif. »\\
Offre : \zh{本周订购，一律八折！} \emph{(Běn zhōu dìnggòu, yílǜ bā zhé!)} « Cette semaine, 20 \% de réduction sur toute commande ! »\\
Action : \zh{快扫码订购吧！} \emph{(Kuài sǎo mǎ dìnggòu ba!)} « Scannez vite le code pour commander ! »\\
Barème : vocabulaire du thème (\zh{有机}, \zh{营养}, \zh{订购}, \zh{折}) / 4 ; structures (\zh{不但……而且……}, impératif) / 4 ; caractères et mise en page / 4 ; cohérence problème → solution → action / 4 ; originalité (slogan, logo, QR) / 4.
\end{corrige}

\newpage

\coursec{Activité d'intégration}

\courssub{Objectifs de l'activité}

À l'issue de cette activité d'intégration, tu seras capable de :

\begin{itemize}
\item mobiliser le \cle{vocabulaire du marketing} (produit, qualité, prix, promotion, client) dans une tâche authentique ;
\item rédiger une \cle{affiche publicitaire} complète en chinois : slogan, arguments, appel à l'action ;
\item écrire correctement les \cle{13 caractères cibles} du thème en respectant l'ordre des traits ;
\item traduire ta propre production (version chinois → français) ;
\item présenter un \cle{argumentaire oral} d'une minute devant la classe ;
\item évaluer la production d'un autre groupe à l'aide d'une grille simplifiée.
\end{itemize}

\courssub{Situation}

Le ministère du Commerce du Cameroun organise à Douala le salon « \cle{Made in Cameroon} ». Cette année, un pavillon spécial est réservé aux exportations vers la Chine. Ton lycée a été invité à participer au concours de la meilleure affiche publicitaire bilingue destinée au marché chinois. Le jury, composé d'enseignants et de commerçants chinois installés à Douala, choisira l'affiche la plus claire, la plus correcte et la plus convaincante.

Voici le message officiel du concours, affiché au tableau :

\begin{textebox}[Annonce du concours — 比赛通知]
比赛通知\\
Bǐsài tōngzhī\\
« Avis de concours »\\[0.4em]
亲爱的同学们：\\
Qīn'ài de tóngxuémen :\\
« Chers camarades, »\\[0.4em]
请你们为喀麦隆产品做一张中文广告海报。\\
Qǐng nǐmen wèi Kāmàilóng chǎnpǐn zuò yì zhāng Zhōngwén guǎnggào hǎibào.\\
« Veuillez créer une affiche publicitaire en chinois pour un produit camerounais. »\\[0.4em]
海报要有口号、三个优点和一个行动号召。\\
Hǎibào yào yǒu kǒuhào, sān ge yōudiǎn hé yí ge xíngdòng hàozhào.\\
« L'affiche doit contenir un slogan, trois avantages et un appel à l'action. »\\[0.4em]
每组还要用法语介绍自己的海报，时间一分钟。\\
Měi zǔ hái yào yòng Fǎyǔ jièshào zìjǐ de hǎibào, shíjiān yì fēnzhōng.\\
« Chaque groupe doit aussi présenter son affiche en français, en une minute. »\\[0.4em]
祝你们成功！\\
Zhù nǐmen chénggōng !\\
« Bonne chance à tous ! »
\end{textebox}

\begin{definition}[Repérage dans l'annonce]
Note le classificateur \zh{张} (zhāng), utilisé pour les objets plats et en feuille : \zh{一张海报} « une affiche », comme \zh{一张纸} « une feuille de papier ». Observe aussi la structure \zh{为 + personne + verbe} : \zh{为喀麦隆产品做……} « faire... pour les produits camerounais ». Enfin, \zh{祝 + quelqu'un + nom/adjectif} est la formule de souhait : \zh{祝你们成功！} « Je vous souhaite le succès ! »
\end{definition}

\courssub{Consignes}

Travail en groupes de 3 ou 4 élèves. Durée totale : 2 heures.

\begin{enumerate}
\item \textbf{Choix du produit (10 min).} Choisissez un produit camerounais dans la liste : le café Arabica de l'Ouest, le beurre de karité du Nord, le tissu Ndop de Bamenda, l'ananas Victoria de la péninsule du Cap Cameroun. Justifiez votre choix en une phrase en chinois : \zh{我们选择……，因为……} (Wǒmen xuǎnzé..., yīnwèi...) « Nous choisissons..., parce que... ».

\item \textbf{Compréhension et repérage (15 min).} Relisez le modèle d'affiche étudié dans la leçon. Soulignez le slogan, les trois arguments et l'appel à l'action. Notez les connecteurs utilisés (\zh{而且}, \zh{所以}, \zh{不仅……还……}).

\item \textbf{Rédaction de l'affiche en chinois (35 min).} Sur une feuille de brouillon, rédigez : un \cle{slogan} court (8 à 12 caractères), \cle{trois arguments} (qualité, prix, origine, santé...) avec au moins deux connecteurs différents, et un \cle{appel à l'action} (\zh{快来买吧！} « Venez vite acheter ! »). Ajoutez un QR code fictif dessiné à la main avec la mention \zh{扫码了解更多} (Sǎo mǎ liǎojiě gèng duō) « Scannez le code pour en savoir plus ».

\item \textbf{Écriture des caractères cibles (20 min).} Sur le grand format A3, écrivez au propre les 13 caractères du thème en respectant l'ordre des traits : 广告产品优质便宜买卖促销牌 (guǎng gào chǎn pǐn yōu zhì pián yi mǎi mài cù xiāo pái). Chaque caractère est écrit quatre fois, grand et lisible.

\item \textbf{Version française (15 min).} Traduisez votre affiche en français, en bas de la feuille A3. Vérifiez que chaque phrase chinoise a son équivalent exact.

\item \textbf{Présentation orale (15 min).} Chaque groupe présente son affiche en une minute : un élève lit le slogan, un autre les arguments, un troisième l'appel à l'action. Parlez lentement et faites attention aux tons.

\item \textbf{Évaluation par les pairs (10 min).} Chaque groupe évalue une autre affiche avec la grille simplifiée (voir Bilan) et donne un conseil d'amélioration.
\end{enumerate}

\courssub{Ressources à mobiliser}

\begin{aretenir}[Boîte à outils de l'affiche publicitaire]
\begin{itemize}
\item \textbf{Slogan} : phrase courte, souvent avec \zh{最} (zuì, « le plus ») : \zh{最好的咖啡！} (Zuì hǎo de kāfēi !) « Le meilleur café ! »
\item \textbf{Arguments} : \zh{不仅……还……} (bùjǐn... hái...) « non seulement... mais aussi... » ; \zh{又……又……} (yòu... yòu...) « à la fois... et... » ; \zh{因为……所以……} « parce que... donc... ».
\item \textbf{Comparatif} : Sujet + \zh{比} + B + adjectif : \zh{我们的咖啡比别的咖啡香。} (Wǒmen de kāfēi bǐ bié de kāfēi xiāng.) « Notre café est plus parfumé que les autres. »
\item \textbf{Appel à l'action} : \zh{快来……吧！} (Kuài lái... ba !) « Venez vite... ! » ; \zh{不要错过！} (Bú yào cuòguò !) « Ne manquez pas ! »
\end{itemize}
\end{aretenir}

\begin{propriete}[Deux façons de cumuler les qualités : 又……又…… et 不仅……还/而且……]
\textbf{1.} Sujet + \zh{又} + adjectif 1 + \zh{又} + adjectif 2 : cumul simple de deux qualités simultanées, registre courant.\\
\zh{这个菠萝又甜又便宜。} (Zhège bōluó yòu tián yòu piányi.) « Cet ananas est à la fois sucré et bon marché. »\\[0.3em]
\textbf{2.} Sujet + \zh{不仅} + qualité 1 + \zh{还/而且} + qualité 2 : insistance, la deuxième qualité ajoute à la première ; registre plus soigné, idéal pour un argumentaire.\\
\zh{我们的咖啡不仅香，而且质量好。} « Notre café est non seulement parfumé, mais aussi de bonne qualité. »\\[0.3em]
\textbf{Attention} : \zh{又……又……} relie deux adjectifs ou deux verbes psychologiques, jamais deux propositions complètes ; pour relier deux phrases, utilise \zh{不仅……而且……}.
\end{propriete}

\begin{center}
\begin{tabularx}{\linewidth}{|X|X|X|X|}
\hline
\rowcolor{popblueL} Caractères & Pinyin & Nature & Traduction \\
\hline
海报 & hǎibào & n. & affiche, poster \\
\hline
口号 & kǒuhào & n. & slogan \\
\hline
产品 & chǎnpǐn & n. & produit \\
\hline
优质 & yōuzhì & adj. & de qualité supérieure \\
\hline
便宜 & piányi & adj. & bon marché \\
\hline
促销 & cùxiāo & v./n. & promouvoir ; promotion \\
\hline
品牌 & pǐnpái & n. & marque \\
\hline
质量 & zhìliàng & n. & qualité \\
\hline
香 & xiāng & adj. & parfumé, odorant \\
\hline
新鲜 & xīnxiān & adj. & frais \\
\hline
不要错过 & bú yào cuòguò & expr. & ne manquez pas \\
\hline
\end{tabularx}
\end{center}

\begin{attention}[Pièges fréquents]
Attention à ne pas confondre \zh{买} (mǎi, acheter) et \zh{卖} (mài, vendre) : 卖 a le trait supplémentaire 十 en haut. Dans le slogan, l'adjectif précède le nom avec \zh{的} : \zh{优质的咖啡} « un café de qualité supérieure », et non l'inverse. N'oubliez pas le 3\textsuperscript{e} ton de \zh{买} (mǎi) et le 4\textsuperscript{e} ton de \zh{卖} (mài) à l'oral : confondre les deux inverse tout le sens de l'affiche ! Enfin, \zh{便宜} se prononce piányi (yi au ton neutre), et non piányì.
\end{attention}

\begin{savaistu}[Le QR code, roi du marketing chinois]
En Chine, presque toutes les affiches publicitaires comportent un QR code (\zh{二维码}, èrwéimǎ) : les consommateurs le scannent avec leur téléphone pour payer, obtenir une réduction ou visiter la boutique en ligne. La mention \zh{扫码了解更多} « scannez le code pour en savoir plus » est donc un marqueur d'authenticité qui donnera à ton affiche un vrai air chinois. Au Cameroun, le paiement mobile se développe aussi rapidement : les deux pays partagent cette culture du « tout-smartphone » dans le commerce.
\end{savaistu}

\courssub{Proposition de production et corrigé commenté}

\begin{exoresolu}[Affiche modèle : le café Arabica de l'Ouest]
Énoncé : rédiger l'affiche complète du groupe « Café Arabica » (slogan, trois arguments, appel à l'action, mention du QR code) et sa version française.
\tcblower
\textbf{Affiche en chinois :}\\[0.3em]
喀麦隆咖啡，世界的味道！\\
Kāmàilóng kāfēi, shìjiè de wèidao !\\
« Café du Cameroun, le goût du monde ! »\\[0.3em]
我们的咖啡不仅香，而且质量好。\\
Wǒmen de kāfēi bùjǐn xiāng, érqiě zhìliàng hǎo.\\
« Notre café est non seulement parfumé, mais aussi de bonne qualité. »\\[0.3em]
它又便宜又新鲜，因为农民直接卖给您。\\
Tā yòu piányi yòu xīnxiān, yīnwèi nóngmín zhíjiē mài gěi nín.\\
« Il est à la fois bon marché et frais, parce que les paysans vous le vendent directement. »\\[0.3em]
买喀麦隆咖啡，就是帮助我们的农民！\\
Mǎi Kāmàilóng kāfēi, jiùshì bāngzhù wǒmen de nóngmín !\\
« Acheter le café du Cameroun, c'est aider nos paysans ! »\\[0.3em]
快来买吧！不要错过！扫码了解更多。\\
Kuài lái mǎi ba ! Bú yào cuòguò ! Sǎo mǎ liǎojiě gèng duō.\\
« Venez vite acheter ! Ne manquez pas ! Scannez le code pour en savoir plus. »\\[0.5em]
\textbf{Commentaire des choix de langue :}
\begin{itemize}
\item Le slogan utilise la structure parallèle nom + \zh{的} + nom, typique des publicités chinoises, courte et rythmée (deux groupes de cinq syllabes).
\item Premier argument : \zh{不仅……而且……} pour cumuler deux qualités ; \zh{香} (xiāng, parfumé) est l'adjectif le plus naturel pour le café.
\item Deuxième argument : \zh{又……又……} avec la cause introduite par \zh{因为} ; le verbe \zh{卖} est bien distingué de \zh{买}, et \zh{卖给您} respecte la structure verbe + \zh{给} + personne.
\item Troisième argument : \zh{就是} (« c'est précisément ») donne un ton émotionnel, efficace en marketing.
\item L'appel final cumule deux impératifs avec \zh{吧} et \zh{不要}, puis la mention moderne du QR code, très courante en Chine.
\end{itemize}
\end{exoresolu}

\courssub{Bilan de l'activité}

\begin{methode}[Grille simplifiée d'évaluation par les pairs]
Attribuez de 0 à 2 points par critère (0 = absent, 1 = partiel, 2 = réussi) :
\begin{enumerate}
\item Le slogan est court, correct et accrocheur.
\item Les trois arguments utilisent au moins deux connecteurs différents.
\item L'appel à l'action est présent et bien formulé.
\item Les 13 caractères cibles sont écrits correctement (traits, proportions).
\item La version française correspond exactement au texte chinois.
\item La présentation orale est claire, avec des tons corrects.
\end{enumerate}
Total sur 12. Donnez ensuite un conseil précis : « Pour améliorer votre affiche, vous pouvez... ».
\end{methode}

Cette activité t'a permis de réunir toutes les compétences de la leçon : comprendre une consigne authentique, choisir le bon lexique du marketing, construire des phrases de persuasion avec les connecteurs étudiés, écrire les caractères avec soin, traduire fidèlement et défendre ton produit à l'oral. Tu es maintenant prêt à concevoir une affiche publicitaire complète en chinois, compétence directement utile dans le monde économique camerounais et chinois d'aujourd'hui.

\newpage

\coursec{Méthodes et exercices résolus}

\coursec{Leçon 32 — Expression écrite : affiche publicitaire sur un produit (marketing)}

\begin{textebox}{Modèle d'affiche commenté : Miel naturel de Buea}
\zh{布埃亚的天然蜂蜜！}\\
\zh{我们的蜂蜜又甜又健康，不但好吃，而且不贵。}\\
\zh{每瓶只要八百法郎！}\\
\zh{在布埃亚的市场有卖。}\\
\zh{快来买吧，不要错过！}\\
\emph{Bù'āiyà de tiānrán fēngmì! Wǒmen de fēngmì yòu tián yòu jiànkāng, búdàn hǎochī, érqiě bú guì. Měi píng zhǐyào bābǎi fǎláng! Zài Bù'āiyà de shìchǎng yǒu mài. Kuài lái mǎi ba, búyào cuòguò!}\\
« Le miel naturel de Buea ! Notre miel est à la fois doux et bon pour la santé ; non seulement il est délicieux, mais en plus il n'est pas cher. Seulement 800 francs le pot ! En vente au marché de Buea. Venez vite acheter, ne manquez pas ! »\\[0.5em]
\textbf{Analyse structurelle :}
\begin{enumerate}
\item \textbf{Titre} : \zh{布埃亚的天然蜂蜜！} (Nom du produit + origine + adjectif qualitatif).
\item \textbf{Arguments} : \zh{又甜又健康} (structure \zh{又……又……}) + \zh{不但好吃，而且不贵} (structure \zh{不但……而且……}).
\item \textbf{Prix} : \zh{每瓶只要八百法郎} (classificateur \zh{瓶} + \zh{只要}).
\item \textbf{Lieu} : \zh{在布埃亚的市场有卖} (complément de lieu \zh{在……} AVANT le verbe \zh{有卖}).
\item \textbf{Slogan} : \zh{快来买吧，不要错过！} (impératifs \zh{快来} + \zh{不要} + particule \zh{吧}).
\end{enumerate}
\end{textebox}

\courssub{Fiche méthode 1 : Produire une affiche publicitaire simple en chinois}

\begin{methode}[Rédiger une affiche publicitaire en 6 étapes]
Une affiche publicitaire (海报 \emph{hǎibào}) n'est pas une rédaction : c'est un texte court, percutant, organisé en blocs. Voici la démarche à suivre au Bac.
\begin{enumerate}
\item \textbf{Identifier le produit et la cible.} Avant d'écrire, réponds en français : qu'est-ce que je vends ? À qui ? (ex. : du jus de mangue, aux lycéens de Yaoundé).
\item \textbf{Écrire un titre accrocheur.} Court, avec un adjectif de qualité : \zh{新鲜芒果汁！} \emph{Xīnxiān mángguǒ zhī!} « Jus de mangue frais ! »
\item \textbf{Donner 2 ou 3 arguments de vente} avec les structures de persuasion : \zh{又……又……} (à la fois... et...), \zh{不但……而且……} (non seulement... mais aussi...), \zh{越……越……} (plus... plus...).
\item \textbf{Ajouter les informations pratiques} : prix (\zh{每瓶五百法郎}), lieu de vente (\zh{在……有卖}), horaires.
\item \textbf{Terminer par un slogan impératif} : \zh{快来买吧！} \emph{Kuài lái mǎi ba!} « Venez vite acheter ! » ; \zh{不要错过！} \emph{Búyào cuòguò!} « Ne manquez pas ! »
\item \textbf{Relire} : vérifier les tons, les classificateurs (\zh{一瓶果汁}, \zh{一件衬衫}) et l'ordre des mots (complément de lieu AVANT le verbe : \zh{在雅温得的市场有卖}).
\end{enumerate}
\textbf{Plan-type à retenir :} Titre accrocheur → qualités du produit → prix → lieu → slogan final.
\end{methode}

\courssub{Fiche méthode 2 : Maîtriser les structures et connecteurs de la persuasion}

\begin{methode}[Les 5 structures qui font gagner des points]
Pour convaincre, le correcteur attend des structures variées. Retiens ces cinq schémas :
\begin{enumerate}
\item \cle{又 + adj. + 又 + adj.} : deux qualités simultanées. \zh{这个手机又便宜又好用。} \emph{Zhège shǒujī yòu piányi yòu hǎoyòng.} « Ce téléphone est à la fois bon marché et pratique. »
\item \cle{不但……而且……} : gradation. \zh{这种茶不但香，而且对身体好。} \emph{Zhè zhǒng chá búdàn xiāng, érqiě duì shēntǐ hǎo.} « Ce thé est non seulement parfumé, mais aussi bon pour la santé. »
\item \cle{越 + verbe/adj. + 越 + verbe/adj.} : progression. \zh{越喝越想喝。} \emph{Yuè hē yuè xiǎng hē.} « Plus on en boit, plus on veut en boire. »
\item \cle{只要 + prix} : « seulement pour ». \zh{每件只要两千法郎！} \emph{Měi jiàn zhǐyào liǎngqiān fǎláng!} « Seulement 2000 francs pièce ! »
\item \cle{对 + nom + 好/有好处} : bénéfice. \zh{吃水果对身体有好处。} \emph{Chī shuǐguǒ duì shēntǐ yǒu hǎochù.} « Manger des fruits est bon pour la santé. »
\end{enumerate}
Réflexe : dans une affiche, utilise au moins DEUX de ces structures, jamais la même deux fois.
\end{methode}

\begin{exoresolu}[Transformer des phrases simples en phrases publicitaires]
Énoncé : Réécris chaque phrase en chinois avec la structure entre parenthèses pour la rendre plus publicitaire.\\
1. Ce savon est bon et pas cher. (又……又……)\\
2. Ce jus est délicieux. Il est aussi bon pour la santé. (不但……而且……)\\
3. Ces chaussures coûtent 5000 FCFA. (只要)
\tcblower
\textbf{Solution détaillée :}\\
1. On combine les deux adjectifs \zh{好} et \zh{便宜} avec \zh{又……又……} :\\
\zh{这个肥皂又好又便宜。} \emph{Zhège féizào yòu hǎo yòu piányi.} « Ce savon est à la fois bon et bon marché. »\\
Justification : \zh{又……又……} encadre deux adjectifs de qualité ; on ne met PAS \zh{是} devant les adjectifs en chinois.\\
2. On fusionne les deux phrases avec \zh{不但……而且……} :\\
\zh{这种果汁不但好喝，而且对身体有好处。} \emph{Zhè zhǒng guǒzhī búdàn hǎohē, érqiě duì shēntǐ yǒu hǎochù.} « Ce jus est non seulement délicieux, mais aussi bon pour la santé. »\\
Justification : le sujet \zh{这种果汁} est placé avant \zh{不但} car il est identique dans les deux propositions ; \zh{而且} introduit la seconde.\\
3. On ajoute \zh{只要} pour l'effet promotionnel, avec le classificateur des chaussures \zh{双} (paire) :\\
\zh{这双鞋只要五千法郎！} \emph{Zhè shuāng xié zhǐyào wǔqiān fǎláng!} « Cette paire de chaussures, seulement 5000 francs ! »\\
Justification : \zh{只} (seulement) se place avant le verbe \zh{要} ; pour des chaussures, le classificateur est \zh{双} (paire), pas \zh{个}.\\
\textbf{Conclusion :} chaque transformation remplace une énumération plate par une structure de persuasion idiomatique — c'est exactement ce que le correcteur récompense.
\end{exoresolu}

\courssub{Fiche méthode 3 : Écrire correctement les caractères du thème (marketing)}

\begin{methode}[Traits, radicaux et caractères proches du lexique de la vente]
\begin{enumerate}
\item \textbf{Identifier le radical (clé)} avant d'écrire : il donne le sens et guide l'ordre des traits. Ex. : \zh{买} (acheter) et \zh{卖} (vendre) partagent la base \zh{卖} = \zh{买} + \zh{十} au-dessus : on VEND ce qu'on met « en plus » sur l'étal.
\item \textbf{Respecter l'ordre des traits} : haut → bas, gauche → droite, horizontal avant vertical, extérieur avant intérieur, on ferme la boîte en dernier (ex. : \zh{国} : on écrit le contenu \zh{玉} avant le trait de fermeture du bas).
\item \textbf{Mémoriser les paires confondues} du thème : \zh{买 mǎi} (acheter) / \zh{卖 mài} (vendre) ; \zh{便 pián} de \zh{便宜} / \zh{使 shǐ} (utiliser) ; \zh{钱 qián} (argent, radical \zh{钅} du métal) / \zh{线 xiàn} (fil) ; \zh{品 pǐn} (trois bouches \zh{口}) / \zh{晶 jīng} (trois soleils \zh{日}) ; \zh{客 kè} (toit \zh{宀} + \zh{夂} + \zh{口}) / \zh{容 róng} (toit \zh{宀} + \zh{谷}) ; \zh{宜 yí} (toit \zh{宀} + \zh{且}) / \zh{谊 yì} (parole \zh{讠} + \zh{义}).
\item \textbf{Écrire chaque caractère 5 fois} en prononçant le pinyin à voix haute : mémoire musculaire + mémoire auditive des tons.
\end{enumerate}
\end{methode}

\begin{center}\begin{tabularx}{\linewidth}{|X|X|X|X|}\hline
\rowcolor{popblueL} Caractères & Pinyin & Nature & Traduction \\ \hline
广告 & guǎnggào & nom & publicité \\ \hline
海报 & hǎibào & nom & affiche \\ \hline
产品 & chǎnpǐn & nom & produit \\ \hline
便宜 & piányi & adj. & bon marché \\ \hline
质量 & zhìliàng & nom & qualité \\ \hline
价格 & jiàgé & nom & prix \\ \hline
顾客 & gùkè & nom & client \\ \hline
打折 & dǎzhé & v. & faire une réduction \\ \hline
\end{tabularx}\end{center}

\begin{exoresolu}[Corriger des caractères mal écrits]
Énoncé : Voici trois phrases d'une affiche d'élève contenant chacune une erreur de caractère. Retrouve l'erreur et réécris la phrase correctement.\\
1. \zh{我们的产晶质量很好。}\\
2. \zh{欢迎顾容来买东西！}\\
3. \zh{这个商店在打折，东西很便谊。}
\tcblower
\textbf{Solution détaillée :}\\
1. \zh{产晶} est faux : \zh{晶} (jīng, cristal, trois soleils \zh{日}) est un homophone proche de \zh{品} (pǐn, trois bouches \zh{口}). Le mot correct est \zh{产品} (produit). Phrase corrigée : \zh{我们的产品质量很好。} \emph{Wǒmen de chǎnpǐn zhìliàng hěn hǎo.} « La qualité de nos produits est très bonne. »\\
2. \zh{顾容} est faux : \zh{容} (róng, contenir) remplace \zh{客} (kè, client). Les deux ont le toit \zh{宀}, mais le bas diffère : \zh{客} contient \zh{夂} + \zh{口}. Phrase corrigée : \zh{欢迎顾客来买东西！} \emph{Huānyíng gùkè lái mǎi dōngxi!} « Bienvenue aux clients, venez faire vos achats ! »\\
3. \zh{便谊} est faux : \zh{谊} (yì, amitié, radical de la parole \zh{讠}) est confondu avec \zh{宜} (yí, convenable, toit \zh{宀}). Le mot correct est \zh{便宜}. Phrase corrigée : \zh{这个商店在打折，东西很便宜。} \emph{Zhège shāngdiàn zài dǎzhé, dōngxi hěn piányi.} « Ce magasin fait des réductions, les articles sont très bon marché. »\\
\textbf{Conclusion :} au Bac, un caractère proche mal écrit coûte un point entier ; la stratégie gagnante est de mémoriser le MOT entier (\zh{产品}, \zh{顾客}, \zh{便宜}) et non le caractère isolé.
\end{exoresolu}

\begin{savaistu}[Marketing en Chine et au Cameroun : convergence et divergence]
\textbf{Chine :} L'affiche papier (海报) reste présente, mais le marketing est dominé par le numérique : \zh{二维码} (codes QR) sur chaque affiche pour paiement \zh{微信支付} / \zh{支付宝}, vente en direct \zh{直播带货} (livestreaming e-commerce), influenceurs \zh{网红}. Les slogans utilisent souvent des jeux de mots sur les chiffres (ex. \zh{666} « fluide, top », \zh{888} « prospérité »).
\textbf{Cameroun :} Affiches peintes à la main sur murs, panneaux en bois au marché, annonces radio/TV, bouche-à-oreille, WhatsApp Business. Le paiement cash (FCFA) ou Mobile Money (MTN MoMo, Orange Money) domine.
\textbf{Point commun :} L'importance du \textbf{relationnel} (\zh{关系} \emph{guānxi} / « connaissance ») et de la \textbf{confiance} (\zh{信任} \emph{xìnrèn}) pour fidéliser le client (\zh{回头客} \emph{huítóukè} / « client fidèle »).
\end{savaistu}

\begin{experience}[Jeu de rôle : « Vendez votre produit ! »]
\textbf{Consigne :} Par binôme. L'élève A est vendeur(se) au marché de Douala (vêtements, jus, téléphone d'occasion...). L'élève B est client(e).\\
\textbf{Étapes :}
\begin{enumerate}
\item A présente son produit avec une mini-affiche orale (titre, 2 qualités, prix, lieu) en utilisant \zh{又……又……} et \zh{只要}.
\item B pose 2 questions : \zh{质量怎么样？} (Qualité ?) \zh{能不能便宜点？} (Moins cher ?).
\item A répond avec \zh{不但……而且……} et propose un geste commercial \zh{给你打个折}.
\item Échangez les rôles.
\end{enumerate}
\textbf{Évaluation par les pairs :} Prononciation des tons (mǎi/mài), classificateurs corrects, fluidité des structures de persuasion.
\end{experience}

\courssub{Fiche méthode 4 : Thème — traduire du français vers le chinois}

\begin{methode}[Traduire une phrase publicitaire (français → chinois)]
\begin{enumerate}
\item \textbf{Repérer le squelette} : sujet + verbe/adjectif. En chinois, l'adjectif seul fait office de prédicat, SANS \zh{是} : « le produit est bon » → \zh{产品很好} (jamais \zh{产品是很好}).
\item \textbf{Placer les compléments au bon endroit} : lieu et moyen AVANT le verbe (\zh{在这个商店买}), durée et compléments de degré APRÈS.
\item \textbf{Choisir le classificateur} : \zh{瓶} (bouteille/pot), \zh{件} (vêtement, article), \zh{双} (paire), \zh{包} (sachet), \zh{公斤} (kilo), \zh{个} (défaut).
\item \textbf{Traduire les marqueurs publicitaires} : « seulement » = \zh{只要} ; « bienvenue » = \zh{欢迎来} ; « ne manquez pas » = \zh{不要错过}.
\item \textbf{Vérifier les tons} en relisant à voix haute : \zh{mǎi} (3\textsuperscript{e} ton, acheter) ≠ \zh{mài} (4\textsuperscript{e} ton, vendre).
\end{enumerate}
\end{methode}

\begin{exoresolu}[Thème : phrases publicitaires]
Énoncé : Traduis en chinois.\\
1. Nos produits sont à la fois bons et bon marché.\\
2. Ce magasin fait une réduction : seulement 1000 francs le kilo de mangues !\\
3. Bienvenue au marché de Douala !
\tcblower
\textbf{Solution détaillée :}\\
1. Squelette : sujet \zh{我们的产品} + double adjectif avec \zh{又……又……}, sans \zh{是} :\\
\zh{我们的产品又好又便宜。} \emph{Wǒmen de chǎnpǐn yòu hǎo yòu piányi.}\\
2. « Faire une réduction » = \zh{打折} ; le complément de prix suit \zh{只要} ; classificateur du kilo : \zh{公斤}.\\
\zh{这个商店在打折：芒果每公斤只要一千法郎！} \emph{Zhège shāngdiàn zài dǎzhé: mángguǒ měi gōngjīn zhǐyào yìqiān fǎláng!}\\
Justification : \zh{在} + verbe marque l'action en cours ; \zh{每} (chaque) + classificateur exprime « le/par ».\\
3. « Bienvenue à + lieu » = \zh{欢迎来 + lieu} :\\
\zh{欢迎来杜阿拉的市场！} \emph{Huānyíng lái Dù'ālā de shìchǎng!}\\
Justification : en chinois, on dit « bienvenue VENIR au marché » ; le verbe \zh{来} est obligatoire, c'est un piège de calque du français.\\
\textbf{Conclusion :} chaque phrase applique une règle précise (pas de \zh{是} devant l'adjectif, place de \zh{在}, \zh{欢迎来}) : la traduction est naturelle, pas calquée.
\end{exoresolu}

\courssub{Fiche méthode 5 : Version — traduire du chinois vers le français}

\begin{methode}[Traduire une affiche chinoise (chinois → français)]
\begin{enumerate}
\item \textbf{Lire l'affiche en entier} une première fois pour identifier le produit, le prix et le lieu.
\item \textbf{Découper chaque phrase} en blocs : sujet / compléments / verbe / complément d'objet.
\item \textbf{Traduire les structures de persuasion par leurs équivalents français} : \zh{又……又……} → « à la fois... et... » ; \zh{不但……而且……} → « non seulement... mais aussi... » ; \zh{只要} → « seulement » ; \zh{越……越……} → « plus... plus... ».
\item \textbf{Rendre le ton publicitaire} en français : phrases courtes, impératifs, points d'exclamation.
\item \textbf{Relire} : le français doit être idiomatique, pas mot à mot (\zh{有卖} = « en vente », pas « il y a vendre »).
\end{enumerate}
\end{methode}

\begin{exoresolu}[Version : une affiche de Yaoundé]
Énoncé : Traduis en français l'affiche suivante.\\
\zh{新鲜咖啡！我们的咖啡又香又便宜，每包只要一千五百法郎。喝咖啡对身体有好处。在雅温得的大市场有卖。快来买吧！}\\
\emph{Xīnxiān kāfēi! Wǒmen de kāfēi yòu xiāng yòu piányi, měi bāo zhǐyào yìqiān wǔbǎi fǎláng. Hē kāfēi duì shēntǐ yǒu hǎochù. Zài Yǎwēndé de dà shìchǎng yǒu mài. Kuài lái mǎi ba!}
\tcblower
\textbf{Solution détaillée, bloc par bloc :}\\
1. \zh{新鲜咖啡！} : titre accrocheur → « Café frais ! »\\
2. \zh{我们的咖啡又香又便宜} : structure \zh{又……又……} → « Notre café est à la fois parfumé et bon marché » ; \zh{香} = parfumé (pour les boissons et la nourriture).\\
3. \zh{每包只要一千五百法郎} : \zh{每} + classificateur \zh{包} (paquet, sachet) + \zh{只要} → « seulement 1500 francs le paquet ».\\
4. \zh{喝咖啡对身体有好处} : \zh{对……有好处} → « boire du café est bon pour la santé ».\\
5. \zh{在雅温得的大市场有卖} : complément de lieu + \zh{有卖} → « en vente au grand marché de Yaoundé ».\\
6. \zh{快来买吧！} : impératif + particule d'exhortation \zh{吧} → « Venez vite l'acheter ! »\\
\textbf{Traduction complète rédigée :} « Café frais ! Notre café est à la fois parfumé et bon marché : seulement 1500 francs le paquet. Boire du café est bon pour la santé. En vente au grand marché de Yaoundé. Venez vite l'acheter ! »\\
\textbf{Conclusion :} la version respecte le sens de chaque structure et restitue le style publicitaire (phrases courtes, impératif final) : double exigence du correcteur.
\end{exoresolu}

\begin{aretenir}[Grille d'évaluation de l'affiche publicitaire (Bac — 10 pts)]
\begin{center}\begin{tabularx}{\linewidth}{|X|X|X|}\hline
\rowcolor{popblueL} Critère & Indicateurs de réussite & Points \\ \hline
\textbf{Respect du format} & Titre distinct, paragraphes courts, slogan final visible & 1 \\ \hline
\textbf{Vocabulaire thème} & 5+ mots du lexique marketing (\zh{产品, 价格, 打折, 顾客, 质量}...) & 2 \\ \hline
\textbf{Structures de persuasion} & Au moins 2 structures différentes correctes (\zh{又……又……}, \zh{不但……而且……}, \zh{只要}, \zh{对……有好处}) & 3 \\ \hline
\textbf{Grammaire / Ordre des mots} & Lieu avant verbe (\zh{在……}), pas de \zh{是} + adj., classificateurs justes & 2 \\ \hline
\textbf{Caractères / Pinyin} & Écriture soignée, pas de caractères proches, tons justes si pinyin fourni & 2 \\ \hline
\end{tabularx}\end{center}
\textbf{Seuil de réussite :} 6/10. \textbf{Excellence :} 9-10/10 (zéro erreur de caractère, structures variées, style naturel).
\end{aretenir}

\begin{exercice}[Production finale : votre affiche]
\textbf{Consigne :} Tu gères un stand de \textbf{ndolé surgelé prêt à cuisiner} (emballé sous vide) au marché Mfoundi à Yaoundé. Rédige une affiche publicitaire en chinois (6-8 lignes) pour attirer les clients pressés (fonctionnaires, étudiants).\\
\textbf{Contraintes obligatoires :}
\begin{enumerate}
\item Titre avec adjectif (ex. \zh{正宗} « authentique », \zh{方便} « pratique »).
\item 2 qualités : goût + praticité (structures \zh{又……又……} OU \zh{不但……而且……}).
\item Prix : 2500 FCFA le sachet (classificateur \zh{包}).
\item Lieu : Marché Mfoundi, Yaoundé (\zh{在雅温得的姆富迪市场}).
\item Slogan final impératif.
\item Utilise \zh{只要} pour le prix.
\end{enumerate}
\textbf{Mots utiles :} \zh{正宗 zhèngzōng} (authentique), \zh{方便 fāngbiàn} (pratique), \zh{冷冻 lěngdòng} (surgelé), \zh{真空包装 zhēnkōng bāozhuāng} (sous vide), \zh{热饭 rèfàn} (riz chaud).
\end{exercice}

\begin{corrige}[Exercice Production finale]
\textbf{Proposition de correction (modèle) :}\\
\zh{雅温得正宗冷冻 ndolé！}\\
\zh{我们的 ndolé 又好吃又方便，不但正宗，而且真空包装。}\\
\zh{每包只要二千五百法郎！}\\
\zh{在雅温得的姆富迪市场有卖。}\\
\zh{买回家热一热，马上开吃！快来买吧！}\\
\emph{Yǎwēndé zhèngzōng lěngdòng ndolé! Wǒmen de ndolé yòu hǎochī yòu fāngbiàn, búdàn zhèngzōng, érqiě zhēnkōng bāozhuāng. Měi bāo zhǐyào èrqiān wǔbǎi fǎláng! Zài Yǎwēndé de Mǔfùdí shìchǎng yǒu mài. Mǎi huíjiā rè yī rè, mǎshàng kāi chī! Kuài lái mǎi ba!}\\
« Ndolé authentique surgelé de Yaoundé ! Notre ndolé est à la fois délicieux et pratique, non seulement authentique, mais aussi emballé sous vide. Seulement 2500 francs le sachet ! En vente au marché Mfoundi de Yaoundé. Achetez, rentrez, réchauffez, mangez tout de suite ! Venez vite acheter ! »\\
\textbf{Points forts du modèle :} Vocabulaire précis (\zh{正宗, 真空包装, 马上}), deux structures de persuasion (\zh{又……又……}, \zh{不但……而且……}), classificateur \zh{包}, lieu avant verbe, slogan dynamique avec séquence verbale \zh{买回家热一热，马上开吃}.
\end{corrige}

\begin{attention}[Les 5 erreurs qui coûtent des points au Bac]
\begin{enumerate}
\item \textbf{Confondre \zh{买} et \zh{卖}.} Acheter = \zh{买} \emph{mǎi} (3\textsuperscript{e} ton) ; vendre = \zh{卖} \emph{mài} (4\textsuperscript{e} ton, avec \zh{十} au-dessus). Erreur de ton ET de caractère : double pénalité.
\item \textbf{Mettre \zh{是} devant un adjectif.} « Le produit est bon » se dit \zh{产品很好}, JAMAIS \zh{产品是很好}. En chinois, l'adjectif est lui-même le prédicat.
\item \textbf{Placer le complément de lieu après le verbe.} Le français dit « acheter au marché » ; le chinois exige \zh{在市场买东西} (lieu AVANT le verbe). Calque du français sanctionné.
\item \textbf{Oublier ou mal choisir le classificateur.} On écrit \zh{一瓶果汁}, \zh{一双鞋}, \zh{一包咖啡}, \zh{一公斤芒果} : jamais \zh{一个果汁}. Le classificateur \zh{个} n'est qu'un recours, pas une solution systématique.
\item \textbf{Écrire des caractères proches à la place du mot juste.} \zh{品}/\zh{晶}, \zh{客}/\zh{容}, \zh{宜}/\zh{谊}, \zh{广}/\zh{床} : mémorise les mots entiers (\zh{产品}, \zh{顾客}, \zh{便宜}, \zh{广告}) et vérifie chaque radical à la relecture.
\end{enumerate}
\end{attention}

\newpage

\coursec{Exercices}

\coursec{Leçon 32 — Expression écrite : affiche publicitaire sur un produit (marketing)}

\courssub{Analyse d'un modèle d'affiche commenté}

\begin{textebox}[Modèle d'affiche : Le poivre de Penja]
\zh{彭雅胡椒——喀麦隆的骄傲！}\\
\zh{彭雅胡椒是世界上最有名的胡椒之一。它又香又辣，味道特别好。}\\
\zh{我们的胡椒是农民亲手种的，百分之百天然，质量有保证。}\\
\zh{现在购买，打八折！}\\
\zh{欢迎大家来我们的展台品尝、购买！}\\
\zh{喀麦隆彭雅胡椒合作社}\\
(Péngyǎ hújiāo —— Kāmàilóng de jiāo'ào!\\
Péngyǎ hújiāo shì shìjiè shàng zuì yǒumíng de hújiāo zhī yī. Tā yòu xiāng yòu là, wèidao tèbié hǎo.\\
Wǒmen de hújiāo shì nóngmín qīnshǒu zhòng de, bǎi fēn zhī bǎi tiānrán, zhìliàng yǒu bǎozhèng.\\
Xiànzài gòumǎi, dǎ bā zhé!\\
Huānyíng dàjiā lái wǒmen de zhǎntái pǐncháng, gòumǎi!\\
Kāmàilóng Péngyǎ Hújiāo Hézuòshè)
\end{textebox}

\begin{aretenir}[Structure type d'une affiche publicitaire (HSK 3-4)]
Une affiche efficace suit un schéma logique en \cle{5 étapes} :
\begin{enumerate}
\item \textbf{Accroche / Slogan} : Court, percutant, souvent nominal. Ex: \zh{彭雅胡椒——喀麦隆的骄傲！}
\item \textbf{Présentation du produit} : Nom, origine, statut. Ex: \zh{是世界上最有名的胡椒之一。}
\item \textbf{Arguments de vente (Qualités)} : Structures \zh{又……又……}, \zh{最}, adjectifs qualitatifs. Ex: \zh{又香又辣}, \zh{百分之百天然}, \zh{质量有保证}。
\item \textbf{Offre commerciale / Incitation} : Prix, réduction (\zh{打折}, \zh{买一送一}), urgence (\zh{现在}). Ex: \zh{现在购买，打八折！}
\item \textbf{Appel à l'action (Call to action) + Signature} : Verbes d'action (\zh{欢迎}, \zh{快来}, \zh{购买}, \zh{品尝}), lieu, émetteur. Ex: \zh{欢迎大家来我们的展台……喀麦隆彭雅胡椒合作社}
\end{enumerate}
\end{aretenir}

\begin{attention}[Pièges à éviter pour l'examen]
\begin{itemize}
\item \textbf{Oubli du \zh{的} (de)} après un adjectif qualificatif : \zh{天然的产品} (pas \zh{天然产品} à l'écrit formel), \zh{最好的蜂蜜}.
\item \textbf{Mauvaise place de \zh{最} (zuì)} : Il se place \textbf{devant} l'adjectif/verbe, jamais après. \zh{最好} (zuì hǎo), \zh{最便宜} (zuì piányi).
\item \textbf{Confusion \zh{买} (mǎi) / \zh{卖} (mài)} : Tons différents (3ème vs 4ème). \zh{买} = acheter (argent sort), \zh{卖} = vendre (argent rentre).
\item \textbf{Classificateurs obligatoires} : \zh{一个展台}, \zh{一种产品}, \zh{一盒茶叶}. Ne pas écrire \zh{买产品} sans contexte générique, préférer \zh{买这个产品} ou \zh{买我们的产品}.
\end{itemize}
\end{attention}

\courssub{Structures et connecteurs pour la persuasion}

\begin{propriete}[Structure comparée : \zh{比} (bǐ) vs \zh{最} (zuì)]
\begin{center}
\begin{tabularx}{\linewidth}{|X|X|X|}
\hline
\rowcolor{popblueL} Structure & Exemple (Caractères + Pinyin) & Traduction \\
\hline
Sujet + \zh{比} + Objet de comparaison + Adjectif & \zh{我们的茶比别的茶便宜。} (Wǒmen de chá bǐ biéde chá piányi.) & Notre thé est moins cher que les autres thés. \\
\hline
Sujet + \zh{比} + Objet + \zh{更} / \zh{还} + Adjectif & \zh{这个品牌比那个品牌更有名。} (Zhège pǐnpái bǐ nàge pǐnpái gèng yǒumíng.) & Cette marque est encore plus célèbre que celle-là. \\
\hline
Sujet + \zh{最} + Adjectif (+ \zh{的} + Nom) & \zh{这是市场上最好的咖啡。} (Zhè shì shìchǎng shàng zuì hǎo de kāfēi.) & C'est le meilleur café du marché. \\
\hline
Sujet + \zh{是} + Groupe nominal + \zh{最} + Adjectif + \zh{的} & \zh{彭雅胡椒是世界上最有名的胡椒之一。} & Le poivre de Penja est l'un des poivres les plus célèbres au monde. \\
\hline
\end{tabularx}
\end{center}
\end{propriete}

\begin{propriete}[Structure \zh{又……又……} (yòu... yòu...) : Double qualité]
\emph{Emploi :} Énumère deux qualités \textbf{simultanées}, généralement positives (ou négatives), pour insister sur la valeur du produit. L'adjectif est souvent monosyllabique ou bisyllabique.
\begin{center}
\begin{tabularx}{\linewidth}{|X|X|X|}
\hline
\rowcolor{popgreenL} Exemple & Pinyin & Traduction \\
\hline
\zh{这个产品又便宜又好用。} & Zhège chǎnpǐn yòu piányi yòu hǎoyòng. & Ce produit est à la fois bon marché et pratique. \\
\zh{喀麦隆可可又香又天然。} & Kāmàilóng kěkě yòu xiāng yòu tiānrán. & Le cacao camerounais est à la fois parfumé et naturel. \\
\zh{我们的蜂蜜又甜又健康。} & Wǒmen de fēngmì yòu tián yòu jiànkāng. & Notre miel est à la fois doux et sain. \\
\hline
\end{tabularx}
\end{center}
\end{propriete}

\begin{attention}[Erreur fréquente] Ne pas confondre avec \zh{既……又……} (jì... yòu...) plus formel/écrit, ni avec \zh{一边……一边……} (yībiān... yībiān...) pour actions simultanées. \zh{又……又……} est le standard pour les affiches publicitaires (style percutant).
\end{attention}


\begin{propriete}[Formules d'appel à l'action (Call to Action)]
\begin{center}
\begin{tabularx}{\linewidth}{|X|X|X|}
\hline
\rowcolor{poporangeL} Formule & Pinyin & Nuance / Contexte \\
\hline
\zh{欢迎……} & Huānyíng... & Accueil poli, standard (ex: \zh{欢迎大家来品尝}). \\
\zh{快来……吧！} & Kuài lái... ba! & Urgence, enthousiasme, oral / affiche dynamique (ex: \zh{快来买吧！}). \\
\zh{赶快……吧！} & Gǎnkuài... ba! & Très pressant, "Dépêchez-vous !" (promo limitée). \\
\zh{请……} & Qǐng... & Formel, courtois (lettre, email pro). \\
\zh{别错过！} & Bié cuòguò! & "Ne manquez pas !" (souvent en bas d'affiche). \\
\hline
\end{tabularx}
\end{center}
\end{propriete}

\begin{definition}[Connecteurs logiques utiles pour l'affiche]
\begin{center}
\begin{tabularx}{\linewidth}{|X|X|X|}
\hline
\rowcolor{poppurpleL} Connecteur & Pinyin & Fonction / Exemple \\
\hline
\zh{而且} & érqiě & Addition (et en plus) : \zh{质量好，而且不贵。} \\
\zh{尤其是} & yóuqíshì & Focus / Mise en relief : \zh{大家都爱喝，尤其是年轻人。} \\
\zh{因为……所以……} & yīnwèi... suǒyǐ... & Cause / Conséquence : \zh{因为天然，所以健康。} \\
\zh{现在} & xiànzài & Temporel / Urgence : \zh{现在购买有优惠。} \\
\zh{总之} & zǒngzhī & Conclusion : \zh{总之，值得买。} \\
\hline
\end{tabularx}
\end{center}
\end{definition}

\begin{experience}[Atelier oral : « Pitch ton produit »]
En binôme. L'élève A tire une carte « Produit camerounais » (Ndolé en conserve, Café Arabica, Miel d'Oku, Tissu Ndop, Jus de gingembre). L'élève B est un acheteur chinois au salon de Shanghai. A doit présenter le produit en 1 minute en utilisant : 1 slogan + \zh{又……又……} + \zh{最} + 1 offre (\zh{打折}/\zh{买一送一}) + \zh{快来……吧}. B note les structures entendues. Inversion des rôles.
\end{experience}

\courssub{Écriture correcte des caractères du thème}

\begin{methode}[Rappel méthodologique : Ordre des traits et radicaux]
\begin{enumerate}
\item \textbf{Ordre général} : Haut $\rightarrow$ Bas ; Gauche $\rightarrow$ Droite ; Horizontal $\rightarrow$ Vertical ; Extérieur $\rightarrow$ Intérieur $\rightarrow$ Fermeture.
\item \textbf{Identification du radical (clé)} : Permet de classer le caractère au dictionnaire et d'en deviner le sens (sémantique) ou le son (phonétique).
\item \textbf{Caractères proches} : Attention aux paires visuellement similaires (ex: \zh{买}/\zh{卖}, \zh{产}/\zh{生}).
\end{enumerate}
\end{methode}

\begin{exemplebox}[Fiches caractères cibles (10 caractères obligatoires)]
\begin{center}
\begin{tabularx}{\linewidth}{|c|X|X|X|X|}
\hline
\rowcolor{popblueL} Car. & Pinyin & Radical (Clé) & Nb traits & Points de vigilance / Caractères proches \\
\hline
\zh{产} & chǎn & \zh{生} (shēng, vie) & 6 & Haut : \zh{立} (debout) + \zh{生}. Sens : « naître, produire ». \textbf{Proche} : \zh{生} (shēng, naître/vivre) — \zh{产} a un trait horizontal de plus en bas. \\
\hline
\zh{品} & pǐn & \zh{口} (kǒu, bouche) & 9 & Trois \zh{口} empilés. Sens originel : « goûter, qualité ». \textbf{Proche} : \zh{呂} (lǚ, nom) — disposition différente. \\
\hline
\zh{质} & zhì & \zh{贝} (bèi, coquillage/argent) & 15 & Simplifié de \zh{質}. Haut : \zh{斤} (hache/poids) + \zh{贝}. Sens : « substance, qualité ». \textbf{Proche} : \zh{贿} (huì, corruption) — même phonétique \zh{质} mais clé \zh{贝} à gauche. \\
\hline
\zh{量} & liáng / liàng & \zh{里} (lǐ, village/lieu) & 12 & Haut : \zh{日} (soleil) + \zh{里}. \zh{liáng} (mesurer), \zh{liàng} (quantité). \textbf{Proche} : \zh{里} (seul). \\
\hline
\zh{广} & guǎng & \zh{广} (guǎng, toit/abri) & 3 & Caractère-clé. « Large, vaste ». \textbf{Proche} : \zh{庚} (gēng), \zh{厂} (chǎng, usine - clé \zh{厂}). \\
\hline
\zh{告} & gào & \zh{口} (kǒu, bouche) & 7 & \zh{牛} (bœuf) au-dessus de \zh{口}. « Annoncer, informer ». \textbf{Proche} : \zh{造} (zào, fabriquer - clé \zh{辶}). \\
\hline
\zh{买} & mǎi & \zh{买} (mǎi - clé propre) & 6 & \zh{卖} sans le trait du haut. « Acheter » (3ème ton). Mnémo : on \textbf{reçoit} (traits vers le bas) l'objet. \\
\hline
\zh{卖} & mài & \zh{十} (shí, dix) / \zh{买} & 8 & \zh{买} + trait horizontal en haut (\zh{十}). « Vendre » (4ème ton). Mnémo : on \textbf{donne} (trait du haut qui part) l'objet. \\
\hline
\zh{价} & jià & \zh{人} (rén, homme) & 8 & Simplifié de \zh{價}. \zh{人} + \zh{襾} (couvrir) + \zh{贝} (argent). « Prix ». \textbf{Proche} : \zh{嘉} (jiā, louange - même phonétique). \\
\hline
\zh{客} & kè & \zh{宀} (mián, toit) & 9 & \zh{宀} + \zh{各}. « Client, invité ». \textbf{Proche} : \zh{容} (róng, contenir - même structure mais \zh{谷} au lieu de \zh{各}). \\
\hline
\end{tabularx}
\end{center}
\end{exemplebox}

\begin{exercice}[Entraînement à l'écriture : Dictée de caractères]
L'enseignant dicte les 10 caractères ci-dessus (dans le désordre). L'élève écrit sur papier quadrillé (style \emph{Tiánzìgé} 田字格). Auto-correction : vérifier l'ordre des traits (numérotation), la proportion dans le carré, la présence du radical.
\end{exercice}

\courssub{Exercices de thème (Français $\rightarrow$ Chinois)}

\begin{exercice}[Thème : Phrases isolées — Vocabulaire et structures de la leçon]
Traduis en chinois (Caractères + Pinyin avec tons).
\begin{enumerate}
\item Notre poivre est naturel et de très bonne qualité.
\item Bienvenue au salon alimentaire de Shanghai !
\item Ce produit est à la fois bon et bon marché.
\item Venez vite acheter le café du Cameroun !
\item Le cacao camerounais est le meilleur du monde.
\item La qualité de nos produits est garantie.
\item Il y a une réduction de 20\% maintenant.
\item Ce thé est plus parfumé que celui-là.
\end{enumerate}
\end{exercice}

\begin{exercice}[Thème : Mini-paragraphe cohérent]
Traduis le paragraphe suivant en un seul paragraphe chinois fluide (connecteurs attendus).
\begin{quote}
« Notre coopérative vend du miel naturel. Ce miel est à la fois doux et sain. C'est le meilleur miel du marché. Achetez maintenant, il y a une réduction. Bienvenue pour goûter ! »
\end{quote}
\end{exercice}

\courssub{Exercices de version (Chinois $\rightarrow$ Français)}

\begin{exercice}[Version : Phrases publicitaires]
Traduis en français soigné (style marketing).
\begin{enumerate}
\item \zh{我们的产品质量最好，价格最便宜。}
\item \zh{喀麦隆咖啡又香又浓，深受中国顾客喜爱。}
\item \zh{新产品上市，买二送一，快来抢购吧！}
\item \zh{欢迎各位客商来我们的展位洽谈合作。}
\item \zh{因为天然，所以放心；因为专业，所以信赖。}
\end{enumerate}
\end{exercice}

\begin{exercice}[Version : Texte complet (Type Bac)]
Traduis l'affiche suivante en français.
\begin{textebox}[Affiche : Thé du Cameroun]
\zh{好消息！好消息！}\\
\zh{喀麦隆茶叶又香又便宜，是最受欢迎的饮料之一。}\\
\zh{今天购买，买一送一！}\\
\zh{快来我们的商店看看吧！}\\
(Hǎo xiāoxi! Hǎo xiāoxi!\\
Kāmàilóng cháyè yòu xiāng yòu piányi, shì zuì shòu huānyíng de yǐnliào zhī yī.\\
Jīntiān gòumǎi, mǎi yī sòng yī!\\
Kuài lái wǒmen de shāngdiàn kànkan ba!)
\end{textebox}
\end{exercice}

\courssub{Grille d'évaluation et production finale (Type Bac)}

\begin{aretenir}[Grille d'évaluation détaillée — Expression écrite : Affiche publicitaire (20 pts)]
\begin{center}
\begin{tabularx}{\linewidth}{|X|c|X|}
\hline
\rowcolor{popblueL} \textbf{Critère} & \textbf{Points} & \textbf{Indicateurs de réussite} \\
\hline
Respect des consignes \& Longueur (80-100 car.) & 4 & Longueur respectée $\pm$ 10 car. ; Tous les éléments obligatoires présents (Slogan, Produit, Qualités, Offre, Appel à l'action, Nom structure). \\
\hline
Correction grammaticale \& Structures & 6 & \textbf{3 structures imposées} utilisées correctement : \zh{最}, \zh{又……又……}, \zh{快来……吧}/\zh{欢迎}, \zh{比}. Pas d'erreur de syntaxe majeure (ordre mots, \zh{的}/\zh{得}/\zh{地}, classificateurs). \\
\hline
Vocabulaire du thème (Caractères cibles) & 4 & \textbf{5/10 caractères cibles} employés correctement : \zh{产}, \zh{品}, \zh{质}, \zh{量}, \zh{广}, \zh{告}, \zh{买}, \zh{卖}, \zh{价}, \zh{客}. Orthographe exacte. \\
\hline
Écriture des caractères (Calligraphie/Lisibilité) & 4 & Traits nets, ordre respecté, proportions équilibrées dans le carré imaginaire, radicaux identifiables. \\
\hline
Effet persuasif \& Présentation visuelle & 2 & Ton enthousiaste, vocabulaire varié (pas que \zh{好}), mise en page aérée (slogan isolé, paragraphes), ponctuation chinoise correcte. \\
\hline
\textbf{TOTAL} & \textbf{20} & \\
\hline
\end{tabularx}
\end{center}
\end{aretenir}

\begin{exercice}[Sujet Type Bac — Expression écrite : Affiche publicitaire]
\textbf{Sujet.} Tu représentes la \textbf{Coopérative du Poivre de Penja} (\zh{喀麦隆彭雅胡椒合作社}). Tu participes au \textbf{Salon Alimentaire International de Shanghai} (\zh{上海国际食品展}). Rédige en chinois une affiche publicitaire pour attirer les clients chinois vers ton stand.

\textbf{Consignes obligatoires :}
\begin{itemize}
\item \textbf{Longueur :} 80 à 100 caractères chinois (ponctuation non comptée).
\item \textbf{Structures imposées (min. 3 distinctes) :}
      \begin{enumerate}
      \item Superlatif avec \zh{最} (ex: \zh{最好的}, \zh{最有名}).
      \item Double qualité avec \zh{又……又……} (ex: \zh{又香又辣}).
      \item Appel à l'action avec \zh{快来……吧} OU \zh{欢迎……}.
      \item Comparaison avec \zh{比} (ex: \zh{比别的胡椒好}).
      \end{enumerate}
\item \textbf{Vocabulaire cible (min. 5 caractères) :} \zh{产}, \zh{品}, \zh{质}, \zh{量}, \zh{广}, \zh{告}, \zh{买}, \zh{卖}, \zh{价}, \zh{客}.
\item \textbf{Contenu obligatoire :} Slogan + Présentation produit/qualités + Offre commerciale (réduction/cadeau) + Appel à l'action + Nom de la coopérative.
\item \textbf{Présentation :} Écriture soignée (traits, radicaux), ponctuation chinoise (，。！), mise en page d'affiche (titre centré, paragraphes).
\end{itemize}
\end{exercice}

\begin{savaistu}[Le saviez-vous ? Le marketing chinois et les chiffres]
En Chine, les chiffres ont une forte connotation culturelle dans le marketing et les prix :
\begin{itemize}
\item \zh{8} (bā) : Son proche de \zh{发} (fā, « s'enrichir, prospérer »). \textbf{Très porte-bonheur}. Prix finissant par 8 (ex: 88 元, 188 元), dates d'ouverture (08/08), téléphone avec beaucoup de 8.
\item \zh{6} (liù) : Son proche de \zh{溜} (liū, « couler, glisser ») $\rightarrow$ « tout se passe bien ». \zh{666} (liùliùliù) = « Excellent, fluide » (argot internet/jeunes).
\item \zh{9} (jiǔ) : Son proche de \zh{久} (jiǔ, « long, durable »). Utilisé pour les mariages, longévité. \zh{99} (jiǔjiǔ) = « Pour toujours ».
\item \zh{4} (sì) : Son proche de \zh{死} (sǐ, « mort »). \textbf{Tabou}. Pas d'étage 4 dans les immeubles, pas de prix en 44, 14, 24...
\end{itemize}
\textbf{Application pour votre affiche :} Préférez \zh{打八折} (20\% remise, son « prospérité ») ou \zh{价格八十八元} plutôt que \zh{打四折} ou \zh{四十四元}.
\end{savaistu}

\begin{popfigure}
\begin{tikzpicture}[
  node distance=1.2cm and 1.5cm,
  box/.style={draw=popblue, thick, rounded corners, fill=popblueL, text width=3.2cm, align=center, font=\small},
  arrow/.style={-Stealth, thick, popdark}
]
\node[box, align=center] (slogan){\textbf{1. SLOGAN}\\\zh{彭雅胡椒——喀麦隆的骄傲！}};
\node[box, below=of slogan, align=center] (prod){\textbf{2. PRODUIT}\\\zh{世界上最有名的胡椒之一}};
\node[box, below=of prod, align=center] (qual){\textbf{3. QUALITÉS}\\\zh{又香又辣 }\\\zh{100\% 天然 }\\\zh{质量有保证}};
\node[box, below=of qual, align=center] (offre){\textbf{4. OFFRE}\\\zh{现在购买 打八折}};
\node[box, below=of offre, align=center] (action){\textbf{5. ACTION}\\\zh{欢迎来展台品尝、购买}};
\node[box, below=of action, align=center] (sign){\textbf{6. SIGNATURE}\\\zh{喀麦隆彭雅胡椒合作社}};

\draw[arrow] (slogan) -- (prod);
\draw[arrow] (prod) -- (qual);
\draw[arrow] (qual) -- (offre);
\draw[arrow] (offre) -- (action);
\draw[arrow] (action) -- (sign);
\end{tikzpicture}
\legende{Schéma de la structure logique d'une affiche publicitaire (Type Bac)}
\end{popfigure}

\newpage

\coursec{Corrigés des exercices}

\begin{corrige}[Exercice 1 — Dictée de caractères]
\textbf{Corrigé de la dictée} (ordre de dictée proposé : \zh{买}, \zh{卖}, \zh{产}, \zh{品}, \zh{客}, \zh{价}, \zh{广}, \zh{告}, \zh{量}, \zh{质}).

\begin{center}
\begin{tabularx}{\linewidth}{|c|X|X|X|}
\hline
\rowcolor{popblueL} Caractère & Pinyin & Radical & Vérification attendue \\
\hline
\zh{买} & mǎi & \zh{买} & 6 traits ; \textbf{sans} le trait horizontal du haut ; 3\textsuperscript{e} ton. \\
\hline
\zh{卖} & mài & \zh{十} & 8 traits ; \zh{买} + \zh{十} au-dessus ; 4\textsuperscript{e} ton. \\
\hline
\zh{产} & chǎn & \zh{生} & 6 traits ; \zh{立} au-dessus de \zh{生}. \\
\hline
\zh{品} & pǐn & \zh{口} & 9 traits ; trois \zh{口} : un en haut, deux en bas ; proportions équilibrées. \\
\hline
\zh{客} & kè & \zh{宀} & 9 traits ; \zh{宀} (toit) d'abord, puis \zh{各} en dessous. \\
\hline
\zh{价} & jià & \zh{人} & 8 traits ; clé \zh{人} à gauche écrite en premier (gauche $\rightarrow$ droite). \\
\hline
\zh{广} & guǎng & \zh{广} & 3 traits ; point en haut, horizontal, puis le trait tombant à gauche. \\
\hline
\zh{告} & gào & \zh{口} & 7 traits ; \zh{牛} en haut (haut $\rightarrow$ bas), \zh{口} en bas. \\
\hline
\zh{量} & liàng & \zh{里} & 12 traits ; \zh{日} en haut, \zh{里} en bas ; ne pas oublier le trait horizontal du milieu. \\
\hline
\zh{质} & zhì & \zh{贝} & 15 traits ; partie haute \zh{斤} d'abord, \zh{贝} en bas ; \zh{贝} se ferme en dernier. \\
\hline
\end{tabularx}
\end{center}

\textbf{Barème indicatif} : 5 points par caractère (ordre des traits 2 pts ; radical bien placé 1 pt ; proportions 1 pt ; lisibilité 1 pt), soit 50 points, ramenés à 10.

\textbf{Points de vigilance} : la confusion \zh{买}/\zh{卖} est l'erreur la plus fréquente — vérifier le sens dicté (acheter = \zh{买}, vendre = \zh{卖}) ; dans \zh{量}, le composant du bas est \zh{里} et non \zh{重}.
\end{corrige}

\begin{corrige}[Exercice 2 — Thème : phrases isolées]
\begin{enumerate}
\item \zh{我们的胡椒天然，质量很好。} (Wǒmen de hújiāo tiānrán, zhìliàng hěn hǎo.)\\
\emph{Justification :} \zh{的} obligatoire après le possessif \zh{我们} ; \zh{很} devant l'adjectif \zh{好} pour une phrase affirmative équilibrée.
\item \zh{欢迎来到上海食品展览会！} (Huānyíng lái dào Shànghǎi Shípǐn Zhǎnlǎnhuì!)\\
\emph{Justification :} \zh{欢迎来到} + lieu = « bienvenue à/au » ; majuscules aux noms propres en pinyin.
\item \zh{这个产品又好又便宜。} (Zhège chǎnpǐn yòu hǎo yòu piányi.)\\
\emph{Justification :} structure \zh{又……又……} pour deux qualités simultanées ; classificateur \zh{个} avec \zh{这}.
\item \zh{快来买喀麦隆咖啡吧！} (Kuài lái mǎi Kāmàilóng kāfēi ba!)\\
\emph{Justification :} \zh{快来……吧} = appel à l'action ; \zh{买} (mǎi, acheter) et non \zh{卖} (mài, vendre).
\item \zh{喀麦隆可可是世界上最好的。} (Kāmàilóng kěkě shì shìjiè shàng zuì hǎo de.)\\
\emph{Justification :} superlatif \zh{最} + adjectif + \zh{的} ; \zh{世界上} = « au monde ». Variante acceptée : \zh{……是世界上最好的可可。}
\item \zh{我们产品的质量有保证。} (Wǒmen chǎnpǐn de zhìliàng yǒu bǎozhèng.)\\
\emph{Justification :} \zh{有保证} = « est garantie » ; \zh{的} entre \zh{产品} et \zh{质量}.
\item \zh{现在打八折。} (Xiànzài dǎ bā zhé.)\\
\emph{Justification :} en chinois, \zh{打八折} signifie « payer 80 \% du prix », donc 20 \% de réduction. Piège classique : ne pas traduire par \zh{打二折} (qui voudrait dire 80 \% de réduction !).
\item \zh{这个茶比那个茶更香。} (Zhège chá bǐ nàge chá gèng xiāng.)\\
\emph{Justification :} structure comparative Sujet + \zh{比} + objet de comparaison + (\zh{更}) + adjectif ; \zh{更} renforce l'écart.
\end{enumerate}
\textbf{Barème indicatif} : 2 pts par phrase (1 pt structure, 1 pt vocabulaire/tons), soit 16 pts.
\end{corrige}

\begin{corrige}[Exercice 3 — Thème : mini-paragraphe cohérent]
\textbf{Traduction modèle :}

\zh{我们合作社卖天然蜂蜜。这种蜂蜜又甜又健康，是市场上最好的蜂蜜。现在购买打九折，欢迎大家来品尝！}

(Wǒmen hézuòshè mài tiānrán fēngmì. Zhè zhǒng fēngmì yòu tián yòu jiànkāng, shì shìchǎng shàng zuì hǎo de fēngmì. Xiànzài gòumǎi dǎ jiǔ zhé, huānyíng dàjiā lái pǐncháng!)

« Notre coopérative vend du miel naturel. Ce miel est à la fois doux et sain, c'est le meilleur miel du marché. Achetez maintenant, 10 \% de réduction ; bienvenue à tous pour goûter ! »

\textbf{Justifications grammaticales :}
\begin{itemize}
\item \zh{卖} (mài) = vendre : c'est la coopérative qui vend, donc \zh{卖} et non \zh{买}.
\item \zh{这种} : \zh{种} (zhǒng) est le classificateur des types/espèces de produits.
\item \zh{又甜又健康} : double qualité avec \zh{又……又……}.
\item \zh{最好的蜂蜜} : superlatif \zh{最} + adjectif + \zh{的} + nom.
\item \zh{打九折} : « une réduction » non précisée en français ; on choisit 10 \% (\zh{九折} = payer 90 \%). Toute valeur cohérente est acceptée (\zh{打八折}, \zh{有优惠}).
\item \zh{欢迎大家来品尝} : appel à l'action standard en fin d'affiche.
\end{itemize}
\textbf{Barème indicatif} : 10 pts (structures imposées 4 pts ; vocabulaire 3 pts ; cohérence/ponctuation chinoise 2 pts ; caractères 1 pt).
\end{corrige}

\begin{corrige}[Exercice 4 — Version : phrases publicitaires]
\begin{enumerate}
\item \zh{我们的产品质量最好，价格最便宜。}\\
« Nos produits sont de la meilleure qualité et aux prix les plus bas. »\\
\emph{Note :} double superlatif \zh{最} ; \zh{价格} = prix, \zh{便宜} = bon marché.
\item \zh{喀麦隆咖啡又香又浓，深受中国顾客喜爱。}\\
« Le café du Cameroun, à la fois parfumé et corsé, est très apprécié des clients chinois. »\\
\emph{Note :} \zh{深受……喜爱} = « être profondément aimé par… » (tournure passive élégante, niveau HSK 4).
\item \zh{新产品上市，买二送一，快来抢购吧！}\\
« Nouveau produit en rayon : deux achetés, un offert ! Dépêchez-vous de l'acheter ! »\\
\emph{Note :} \zh{上市} = « être lancé sur le marché » ; \zh{买二送一} = « pour deux achetés, un offert » ; \zh{抢购} = « se ruer pour acheter ».
\item \zh{欢迎各位客商来我们的展位洽谈合作。}\\
« Nous souhaitons la bienvenue à tous les hommes d'affaires venus sur notre stand pour discuter d'une coopération. »\\
\emph{Note :} \zh{客商} = clients commerçants ; \zh{洽谈} = négocier, discuter (registre formel).
\item \zh{因为天然，所以放心；因为专业，所以信赖。}\\
« Parce que c'est naturel, vous êtes rassurés ; parce que c'est professionnel, vous pouvez nous faire confiance. »\\
\emph{Note :} structure parallélique \zh{因为……所以……} répétée, typique des slogans chinois ; rendre le parallélisme en français.
\end{enumerate}
\textbf{Barème indicatif} : 2 pts par phrase (fidélité 1 pt ; qualité du français marketing 1 pt), soit 10 pts.
\end{corrige}

\begin{corrige}[Exercice 5 — Version : texte complet « Thé du Cameroun »]
\textbf{Traduction modèle :}

« Bonne nouvelle ! Bonne nouvelle !\\
Le thé du Cameroun est à la fois parfumé et bon marché ; c'est l'une des boissons les plus appréciées.\\
Aujourd'hui, pour tout achat : un acheté, un offert !\\
Venez vite jeter un coup d'œil dans notre magasin ! »

\textbf{Justifications :}
\begin{itemize}
\item \zh{好消息} répété deux fois : procédé d'accroche oral ; on garde la répétition en français pour l'effet publicitaire.
\item \zh{又香又便宜} : double qualité (\zh{又……又……}) → « à la fois parfumé et bon marché ».
\item \zh{最受欢迎的饮料之一} : \zh{最} + \zh{受……欢迎} (être populaire auprès de) + \zh{之一} (l'un parmi) → « l'une des boissons les plus populaires ».
\item \zh{买一送一} : formule commerciale figée → « un acheté, un offert ».
\item \zh{看看吧} : \zh{看看} (redoublement adoucissant) + \zh{吧} (invitation) → « venez jeter un coup d'œil », ton léger et chaleureux.
\end{itemize}
\textbf{Barème indicatif} : 10 pts (fidélité du sens 5 pts ; rendu du style publicitaire 3 pts ; correction du français 2 pts).
\end{corrige}

\begin{corrige}[Exercice 6 — Sujet Type Bac : affiche publicitaire (Poivre de Penja)]
\textbf{Proposition de rédaction modèle (92 caractères hors ponctuation) :}

\begin{textebox}[Corrigé modèle — Affiche Poivre de Penja]
\zh{彭雅胡椒——喀麦隆的骄傲！}\\
\zh{我们的产品质量最好，又香又辣，比别的胡椒更天然。}\\
\zh{展会特价，买二送一！}\\
\zh{欢迎各位客商来展台品尝、购买！}\\
\zh{喀麦隆彭雅胡椒合作社}\\
(Péngyǎ hújiāo —— Kāmàilóng de jiāo'ào!\\
Wǒmen de chǎnpǐn zhìliàng zuì hǎo, yòu xiāng yòu là, bǐ biéde hújiāo gèng tiānrán.\\
Zhǎnhuì tèjià, mǎi èr sòng yī!\\
Huānyíng gèwèi kèshāng lái zhǎntái pǐncháng, gòumǎi!\\
Kāmàilóng Péngyǎ Hújiāo Hézuòshè)
\end{textebox}

\textbf{Analyse du modèle :}
\begin{itemize}
\item \textbf{Structures imposées (4/3) :} \zh{质量最好} (superlatif \zh{最}) ; \zh{又香又辣} (double qualité) ; \zh{比别的胡椒更天然} (comparatif \zh{比} + \zh{更}) ; \zh{欢迎各位客商……} (appel à l'action).
\item \textbf{Caractères cibles (5/5) :} \zh{产} (\zh{产品}), \zh{质} et \zh{量} (\zh{质量}), \zh{买} (\zh{买二送一}, \zh{购买}), \zh{客} (\zh{客商}).
\item \textbf{Contenu :} slogan (ligne 1) ; présentation et qualités (ligne 2) ; offre commerciale \zh{买二送一} (ligne 3) ; appel à l'action + lieu (ligne 4) ; signature (ligne 5).
\end{itemize}

\textbf{Barème indicatif (20 pts) :}
\begin{center}
\begin{tabularx}{\linewidth}{|X|c|}
\hline
\rowcolor{popblueL} Critère & Points \\
\hline
Respect des consignes et longueur (80-100 caractères, 5 éléments présents) & 4 \\
\hline
Correction grammaticale et 3 structures imposées (\zh{最}, \zh{又……又……}, \zh{比}, \zh{欢迎}/\zh{快来……吧}) & 6 \\
\hline
Vocabulaire du thème : au moins 5 caractères cibles correctement écrits & 4 \\
\hline
Écriture des caractères (traits, radicaux, proportions, lisibilité) & 4 \\
\hline
Effet persuasif et présentation visuelle (mise en page, ponctuation chinoise) & 2 \\
\hline
\textbf{Total} & \textbf{20} \\
\hline
\end{tabularx}
\end{center}

\textbf{Points de vigilance :}
\begin{itemize}
\item Ne pas écrire \zh{打四折} pour « 40 \% de réduction » : \zh{打四折} signifie payer 40 \% du prix, soit 60 \% de réduction. Pour 40 \% de réduction, écrire \zh{打六折}.
\item \zh{最} se place toujours \textbf{devant} l'adjectif : \zh{最好}, jamais \zh{好最}.
\item \zh{买} (mǎi, acheter) / \zh{卖} (mài, vendre) : vérifier le sens et le ton.
\item Ne pas oublier le \zh{的} dans \zh{最好的产品}, \zh{我们的胡椒}.
\item Ponctuation chinoise pleine chasse : ， 。 ！ — et signature de l'affiche obligatoire.
\end{itemize}
\end{corrige}

\newpage

\coursec{Fiche bilan}

\begin{popfigure}
\begin{tikzpicture}[mindmap, grow cyclic, every node/.style={concept, text=white, font=\small},
  root concept/.append style={concept color=popink, font=\bfseries},
  level 1/.append style={level distance=3.6cm, sibling angle=51.4},
  level 1 concept/.append style={font=\footnotesize}]
\node[root concept, align=center]{Affiche publicitaire\\(marketing)}
  child[concept color=popblue]{node[align=center]{Lexique du produit\\产品 质量 价格}}
  child[concept color=poporange]{node[align=center]{Verbes de persuasion\\买 试 欢迎 别错过}}
  child[concept color=popgreen]{node[align=center]{Structures\\太……了 / 又……又……}}
  child[concept color=poppurple]{node[align=center]{Connecteurs\\而且 所以 如果……就}}
  child[concept color=poppink]{node[align=center]{Mise en page\\titre + slogan + infos}}
  child[concept color=popgold]{node[align=center]{Caractères\\radicaux 讠 贝 钅}}
  child[concept color=popmuted]{node[align=center]{Traduction\\thème et version}};
\end{tikzpicture}
\legende{Carte mentale de la leçon : rédiger une affiche publicitaire en chinois}
\end{popfigure}

\begin{definition}[Lexique clé de la publicité]
\begin{itemize}
  \item \zh{产品} chǎnpǐn : produit ; \zh{质量} zhìliàng : qualité ; \zh{价格} jiàgé : prix ; \zh{便宜} piányi : bon marché
  \item \zh{广告} guǎnggào : publicité ; \zh{海报} hǎibào : affiche ; \zh{顾客} gùkè : client ; \zh{优惠} yōuhuì : réduction, avantage
  \item \zh{欢迎} huānyíng : bienvenue, accueillir ; \zh{试一试} shì yi shì : essayer (forme redoublée adoucie) ; \zh{别错过} bié cuòguò : ne manquez pas
  \item \zh{新款} xīnkuǎn : nouveau modèle ; \zh{名牌} míngpái : grande marque ; \zh{打折} dǎzhé : faire une remise
\end{itemize}
\end{definition}

\begin{propriete}[Grammaire à connaître par cœur]
\begin{center}
\begin{tabularx}{\linewidth}{|X|X|X|}
\hline
\rowcolor{popblueL} Structure & Exemple & Traduction \\
\hline
太 + adj. + 了 & \zh{太好了！} Tài hǎo le ! & C'est vraiment bien ! \\
\hline
又 + adj. + 又 + adj. & \zh{又便宜又好} yòu piányi yòu hǎo & à la fois bon marché et bon \\
\hline
Impératif / invitation & \zh{快来买吧！} Kuài lái mǎi ba ! & Venez vite acheter ! \\
\hline
别 + verbe (interdiction) & \zh{别错过！} Bié cuòguò ! & Ne manquez pas ! \\
\hline
如果……就…… & \zh{如果你现在买，就可以打折。} Rúguǒ nǐ xiànzài mǎi, jiù kěyǐ dǎzhé. & Si tu achètes maintenant, tu peux avoir une remise. \\
\hline
只要……就…… & \zh{只要一百块，就能买到。} Zhǐyào yìbǎi kuài, jiù néng mǎi dào. & Pour seulement 100 yuans, on peut l'acheter. \\
\hline
\end{tabularx}
\end{center}
\end{propriete}

\begin{attention}[Pièges fréquents des francophones]
\begin{itemize}
  \item \zh{太……了} exige le \zh{了} final : on écrit \zh{太好了}, jamais \zh{太好} seul dans une exclamation.
  \item \zh{又……又……} ne s'emploie qu'avec des adjectifs ou verbes psychologiques, pas avec des verbes d'action : on dit \zh{又便宜又好}, mais pas \zh{又买又吃} au sens de « en même temps ».
  \item Ne pas confondre \zh{买} mǎi (acheter, 3\ieme{} ton) et \zh{卖} mài (vendre, 4\ieme{} ton) : inversion catastrophique dans une publicité !
  \item \zh{别} exprime l'interdiction (« ne... pas »), \zh{不} la simple négation : \zh{别错过} (ne manquez pas) / \zh{不错} (pas mal).
\end{itemize}
\end{attention}

\begin{methode}[Plan type d'une affiche publicitaire]
\begin{enumerate}
  \item Titre accrocheur : \zh{新产品来了！} Xīn chǎnpǐn lái le ! (Le nouveau produit est arrivé !)
  \item Slogan avec \zh{太……了} ou \zh{又……又……} : \zh{又好看又耐用！} Yòu hǎokàn yòu nàiyòng ! (À la fois beau et solide !)
  \item Arguments : qualité, prix, avantages : \zh{质量好，价格低。} Zhìliàng hǎo, jiàgé dī. (Bonne qualité, prix bas.)
  \item Appel à l'action : \zh{欢迎光临！快来试一试！} Huānyíng guānglín ! Kuài lái shì yi shì ! (Bienvenue ! Venez vite essayer !)
  \item Informations pratiques : lieu, date, prix : \zh{地址、时间、价格} dìzhǐ, shíjiān, jiàgé (adresse, horaires, prix).
\end{enumerate}
\end{methode}

\begin{propriete}[Écriture des caractères du thème]
\begin{itemize}
  \item \zh{告} (gào, dans \zh{广告}) : clé \zh{口} (bouche) ; écrire \zh{牛} au-dessus, \zh{口} en dessous, de haut en bas.
  \item \zh{质} (zhì, dans \zh{质量}) : clé \zh{贝} (coquillage, argent) ; attention à ne pas confondre avec \zh{盾} dùn (bouclier).
  \item \zh{价} (jià, dans \zh{价格}) : clé \zh{亻} (homme) à gauche, puis \zh{介} à droite ; gauche avant droite.
  \item \zh{便} (pián, dans \zh{便宜}) : clé \zh{亻} ; ne pas confondre avec \zh{使} shǐ (utiliser).
  \item \zh{错} (cuò, dans \zh{别错过}) : clé \zh{钅} (métal) ; pour la partie \zh{昔}, horizontal avant vertical.
\end{itemize}
\end{propriete}

\begin{methode}[Méthode express de traduction]
\begin{itemize}
  \item Thème (français → chinois) : repérer sujet + verbe, placer le complément de temps ou de lieu avant le verbe, vérifier les tons et le classificateur.
  \item Version (chinois → français) : traduire d'abord le sens global, puis rendre le slogan en français naturel (jamais mot à mot).
\end{itemize}
\end{methode}

\begin{aretenir}[Checklist avant le Bac]
\begin{itemize}
  \item $\square$ Je connais le lexique du produit : \zh{产品、质量、价格、广告、优惠、顾客}.
  \item $\square$ Je sais écrire sans erreur \zh{广告、质量、便宜、欢迎、错过} (traits et radicaux).
  \item $\square$ Je maîtrise la structure d'exclamation \zh{太……了}.
  \item $\square$ Je sais utiliser \zh{又……又……} pour cumuler deux qualités.
  \item $\square$ Je sais formuler un appel à l'action : \zh{快来……吧！} et \zh{别错过！}.
  \item $\square$ Je place correctement les compléments de temps et de lieu avant le verbe.
  \item $\square$ Je sais construire une phrase de condition avec \zh{如果……就……}.
  \item $\square$ Mon affiche a un titre, un slogan, des arguments et des informations pratiques.
  \item $\square$ Je vérifie les tons du pinyin sur chaque mot (\zh{mǎi} acheter / \zh{mài} vendre).
  \item $\square$ Je n'oublie pas le classificateur : \zh{一个产品、一件衣服、一部手机}.
  \item $\square$ Je relis ma production : ordre des mots, caractères proches, ponctuation chinoise 。！
  \item $\square$ Je sais traduire un slogan en français naturel, sans mot à mot.
\end{itemize}
\end{aretenir}

\findecours

\end{document}
